% Generated mechanically from frozen input inputs/p10/ja/formal-defos.tex.
% Do not edit this generated file; edit the bound locale source instead.


% OK, start here.
%

\setcounter{chapter}{89}
\chapter{形式変形理論}



\phantomsection
\label{formal-defos-section-phantom}




\section{序論}
\label{formal-defos-section-introduction}

\noindent
本章では，Stacks プロジェクトで後に適用できる形で形式変形理論を展開し，
Rim \cite[Exposee VI]{SGA7-I} および Schlessinger \cite{Sch} に密接に従う。
この話題を初めて学ぶ読者には，まず Schlessinger の論文を読むことを強く勧める。
同論文はほとんどの応用に十分なほど一般的であり，実際，この種の形式変形理論を
用いる論文の大部分で Schlessinger の結果が使われているからである。

\medskip\noindent
$\Lambda$ を剰余体 $k$ をもつ完備 Noether 局所環とし，
$\mathcal{C}_\Lambda$ を，剰余体 $k$ をもつ Artin 局所
$\Lambda$-代数の圏とする。$F(k)$ が一点集合となるような関手
$F : \mathcal{C}_\Lambda \to \textit{Sets}$ が与えられたとき，
Schlessinger の論文では，次を成り立たせる条件 (H1)--(H4) が導入された：
\begin{enumerate}
\item (H1)--(H3) が成り立つことと，$F$ が「hull」をもつこととは同値である。
\item (H1)--(H4) が成り立つことと，$F$ が pro-表現可能であることとは同値である。
\end{enumerate}
本章の目的は，Rim の論文とまったく同じく，これらの結果を次の二方向に一般化する
ことである：
\begin{enumerate}
\item[(A)] 関手 $F$ を，$\mathcal{C}_\Lambda$ 上で群亜群をファイバーとする
余ファイバー圏 $\mathcal{F}$ に置き換える。第 \ref{formal-defos-section-CLambda} 節を参照せよ。
\item[(B)] $\Lambda$ を Noether 環とし，$\Lambda \to k$ を体への有限環写像とする。
圏 $\mathcal{C}_\Lambda$ は，同一視 $A/\mathfrak m_A = k$ を指定した
Artin 局所 $\Lambda$-代数 $A$ の圏である。
\end{enumerate}
$F(k)$ が一点集合であるという条件に対応するのは，$\mathcal{F}(k)$ が自明な
群亜群であるという条件である。$\mathcal{F}$ がこの条件を満たすとき，これを
{\it 前変形圏}と呼ぶが，一般にはこの仮定を置かない。Rim の論文
\cite[Exposee VI]{SGA7-I} が本稿の結果の原典である。また，群亜群を用いた
変形理論を論じる有用な論文 \cite{Talpo-Vistoli} にも触れておく。ただし，
そこで扱われる一般性は本章よりも低い。

\medskip\noindent
圏 $\mathcal{C}_\Lambda$ の「完備化」
$\widehat{\mathcal{C}}_\Lambda$ は重要な役割を果たす。
$\widehat{\mathcal{C}}_\Lambda$ の対象は，剰余体を $k$ と同一視した
Noether 完備局所 $\Lambda$-代数 $R$ である。第
\ref{formal-defos-section-category-completion-CLambda} 節を参照せよ。
一方で $\mathcal{C}_\Lambda \subset \widehat{\mathcal{C}}_\Lambda$ は
厳密充満部分圏であり，他方で $\widehat{\mathcal{C}}_\Lambda$ は
$\mathcal{C}_\Lambda$ の pro-対象の圏の充満部分圏である。関手
$\mathcal{C}_\Lambda \to \textit{Sets}$ が {\it pro-表現可能}であるとは，
$R \in \Ob(\widehat{\mathcal{C}}_\Lambda)$ に対する表現可能関手
$\underline{R} = \Mor_{\widehat{\mathcal{C}}_\Lambda}(R, -)$ の
$\mathcal{C}_\Lambda$ への制限と同型であることをいう。

\medskip\noindent
{\it 群亜群をファイバーとする余ファイバー圏}は，群亜群をファイバーとする
ファイバー圏の双対である。これを第 \ref{formal-defos-section-preliminary} 節で導入する。
$\mathcal{C}_\Lambda$ 上で群亜群をファイバーとする余ファイバー圏の射が
{\it 滑らか}であるとは，対象についての無限小持ち上げ判定法を満たすことをいう。
第 \ref{formal-defos-section-smooth-morphisms} 節を参照せよ。
これは形式的に滑らかな環写像の定義に対応する。『代数』定義
\ref{algebra-definition-formally-smooth} を参照せよ。また，これは
『表現可能性の判定法』第 \ref{criteria-section-formally-smooth} 節の概念と
ちょうど双対である。これは重要な概念である。というのも，最終的には，
群亜群をファイバーとするある種の余ファイバー圏が，滑らかな pro-表現可能
表示をもつことを証明したいからであり，これは『代数スタック』第
\ref{algebraic-section-stack-to-presentation} 節および第
\ref{algebraic-section-smooth-groupoid-gives-algebraic-stack} 節における
代数スタックの特徴付けに似ている。$\mathcal{C}_\Lambda$ 上で群亜群を
ファイバーとする余ファイバー圏 $\mathcal{F}$ の {\it versal 形式対象}とは，
対応する射
$\underline{\xi} : \underline{R}|_{\mathcal{C}_\Lambda} \to \mathcal{F}$
が滑らかとなるような，完備化の対象
$\xi \in \widehat{\mathcal{F}}(R)$ のことである。

\medskip\noindent
第 \ref{formal-defos-section-schlessinger-conditions} 節では，Schlessinger の (H1) と (H2) を
一般化する条件 (S1) と (S2) を $\mathcal{F}$ に対して定義する。
Schlessinger の (H3)，すなわち $\mathcal{F}$ の接空間が有限次元であるという
条件には名称を与えない。この理論の展開における重要な段階は，(S1)，(S2)，
(H3) を満たす前変形圏に versal 形式対象が存在することである。補題
\ref{formal-defos-lemma-versal-object-existence} を参照せよ。
関手 $F : \mathcal{C}_\Lambda \to \textit{Sets}$ に対する
Schlessinger の {\it hull} とは，本章の用語では，接空間の誘導写像
$d\underline{\xi} : T\underline{R}|_{\mathcal{C}_\Lambda} \to TF$
が同型となるような versal 形式対象
$\xi \in \widehat{F}(R)$ のことである。
文献では hull を「miniversal」対象と呼ぶことが多いが，本章ではそう呼ばない。
その理由は次のとおりである。関手が hull をもたずに versal 形式対象をもつことが
ありうる。さらに，第 \ref{formal-defos-section-minimal-versal} 節で示すように，前変形圏が
versal 形式対象をもつならば，定義 \ref{formal-defos-definition-minimal-versal} の意味で
{\it 極小}なものを必ずもち，それは同型を除いて一意である。補題
\ref{formal-defos-lemma-minimal-versal} を参照せよ。しかし，極小 versal 形式対象が
接空間上の同型を誘導しないこともありうるのである（例
\ref{formal-defos-example-do-not-get-S2} および \ref{formal-defos-example-smooth-continued} を参照）。

\medskip\noindent
上で指摘した相違を念頭に置けば，定理
\ref{formal-defos-theorem-miniversal-object-existence} は上の (1) の直接の一般化である。
$\mathcal{F}$ が関手である場合には Schlessinger の結果を回収し，条件 (S1) と
(S2) のもとで，接空間上の写像
$d\underline{\xi} : T\underline{R}|_{\mathcal{C}_\Lambda} \to TF$
によって極小 versal 形式対象を特徴付ける。

\medskip\noindent
第 \ref{formal-defos-section-RS-condition} 節では，Schlessinger の (H4) を一般化する
Rim の条件 (RS) を $\mathcal{F}$ に対して定義する。{\it 変形圏}とは，
(RS) を満たす前変形圏として定義される。pro-表現可能関手に対応するのは，
$\mathcal{C}_\Lambda$ 上で群亜群をファイバーとする余ファイバー圏のうち，
$\mathcal{C}_\Lambda$ 上の {\it 関手における滑らかな pro-表現可能群亜群による表示}
をもつものである。定義 \ref{formal-defos-definition-groupoid-in-functors}，
\ref{formal-defos-definition-prorepresentable-groupoid-in-functors}，および
\ref{formal-defos-definition-smooth-groupoid-in-functors} を参照せよ。
この表示の概念は，$\mathcal{F}$ のファイバーの群亜群構造を考慮に入れている。
定理 \ref{formal-defos-theorem-presentation-deformation-groupoid} では，$\mathcal{F}$ が
関手における滑らかな pro-表現可能群亜群による表示をもつことと，
$\mathcal{F}$ の接空間および無限小自己同型空間が有限次元であることとが
同値であることを証明する。これは上の (2) の一般化であり，$\mathcal{F}$ が
関手である場合には Schlessinger の結果に帰着する。最後の第
\ref{formal-defos-section-minimality} 節では，極小 versal 形式対象を用いて，
関手における滑らかな pro-表現可能群亜群による極小表示を
（同型を除いて一意に）構成する方法を論じる。

\medskip\noindent
また，Schlessinger の条件について次の概念的説明も得られる。
前変形圏 $\mathcal{F}$ が (RS) を満たすなら，同型類の付随関手
$\overline{\mathcal{F}}: \mathcal{C}_\Lambda \to \textit{Sets}$ は
(H1) と (H2) を満たす（補題 \ref{formal-defos-lemma-RS-implies-S1-S2} および
\ref{formal-defos-lemma-S1-S2-associated-functor}）。逆に，関手
$F : \mathcal{C}_\Lambda \to \textit{Sets}$ が，群亜群をファイバーとする
余ファイバー圏 $\mathcal{F}$ の同型類の関手として自然に生じるならば，
実際の例では，$F$ が (H1) と (H2) を満たすことを示す議論により，
$\mathcal{F}$ が (RS) を満たすことも示されるようである。例については，
『変形問題』第 \ref{examples-defos-section-introduction} 節で論じる。
さらに，$\mathcal{F}$ が (RS) を満たすなら，$\overline{\mathcal{F}}$ に
対する条件 (H4) は，$\mathcal{F}$ の対象の自己同型を延長することによって
簡潔に解釈できる（補題 \ref{formal-defos-lemma-RS-associated-functor}）。
これらの観察は，(RS) を基本的な変形理論的貼り合わせ条件とみなすべきことを
示唆している。




\section{記法と規約}
\label{formal-defos-section-notations-conventions}

\noindent
環は単位元 $1$ をもつ可換環とする。局所環 $A$ の極大イデアルを
$\mathfrak{m}_A$ で表す。正の整数の集合を
$\mathbf{N} = \{1, 2, 3, \ldots\}$ で表す。$U$ が圏
$\mathcal{C}$ の対象なら，関手
$\Mor_\mathcal{C}(U, -): \mathcal{C} \to \textit{Sets}$ を
$\underline{U}$ で表す。評注 \ref{formal-defos-remarks-cofibered-groupoids}
（\ref{formal-defos-item-definition-yoneda}）を参照せよ。注意：これは，他章でときどき
$\underline{U}$ を $h_U(-) = \Mor_\mathcal{C}(-, U)$ の意味で用いる記法と
衝突する可能性がある。

\medskip\noindent
本章を通じて $\Lambda$ は Noether 環，$\Lambda \to k$ は
$\Lambda$ から体への有限環写像である。この写像の核を
$\mathfrak m_\Lambda$，像 $k' \subset k$ と表す。すると
$\mathfrak m_\Lambda$ は極大イデアル，
$k' = \Lambda/\mathfrak m_\Lambda$ は体であり，拡大 $k/k'$ は有限である。
（\ref{formal-defos-equation-k-prime}）の周囲の議論を参照せよ。



\section{基底圏}
\label{formal-defos-section-CLambda}

\noindent
動機。形式変形理論の重要な応用の一つは，代数空間による表現可能性の
判定法にある。局所 Noether 基底 $S$ と関手
$F : (\Sch/S)_{fppf}^{opp} \to \textit{Sets}$
が与えられているとする。$k$ を $S$ 上有限型の体，すなわち有限型射
$\Spec(k) \to S$ が与えられた体とする。Artin の判定法の一つは，任意の元
$x \in F(\Spec(k))$ に対し，三つ組 $(S, k, x)$ に付随する前変形関手が
pro-表現可能であるべきだというものである。『射』補題
\ref{morphisms-lemma-point-finite-type} により，$k$ が $S$ 上有限型であるという
条件は，$k$ が有限 $\Lambda$-代数となるようなアフィン開集合
$\Spec(\Lambda) \subset S$ が存在することを意味する。このことから，本章を
通じて次の基底圏を用いる理由が分かる。

\begin{definition}
\label{formal-defos-definition-CLambda}
$\Lambda$ を Noether 環，$\Lambda \to k$ を有限環写像とし，$k$ は体であると
する。{\it $\mathcal{C}_\Lambda$} を次の圏として定義する。
\begin{enumerate}
\item 対象は対 $(A, \varphi)$ であって，$A$ は Artin 局所
$\Lambda$-代数，$\varphi : A/\mathfrak m_A \to k$ は
$\Lambda$-代数同型である。
\item 射 $f : (B, \psi) \to (A, \varphi)$ は，
$\varphi \circ (f \bmod \mathfrak m) = \psi$ を満たす局所
$\Lambda$-代数準同型である。
\end{enumerate}
$\Lambda$ が Noether 完備局所環であり，$k$ がその剰余体である場合を
{\it 古典的場合}と呼ぶ。
\end{definition}

\noindent
$\Lambda \to k$ が全射で，$A$ が Artin 局所 $\Lambda$-代数なら，
同一視 $\varphi$ は，存在する限り一意であることに注意せよ。さらにこの場合，
任意の $\Lambda$-代数写像 $A \to B$ はこれらの同一視と両立する。したがって
この場合，$\mathcal{C}_\Lambda$ は，剰余体が「$k$ である」局所 Artin
$\Lambda$-代数の圏にすぎない。記法を濫用して，一般の場合にも
$\mathcal{C}_\Lambda$ の対象を単に $A$ と表す。さらに，しばしば
$A/\mathfrak m = k$ と書く。すなわち，$\mathcal{C}_\Lambda$ のすべての環の
剰余体が $k$ であるかのように扱う（$\mathcal{C}_\Lambda$ の環写像はすべて
指定された同一視と両立するので，これにより問題が生じることはない）。
本章の残りを通じて，基底環 $\Lambda$ と体 $k$ を固定する。
圏 $\mathcal{C}_\Lambda$ は，以下で考える余ファイバー圏の基底圏となる。

\begin{definition}
\label{formal-defos-definition-small-extension}
$f: B \to A$ を $\mathcal{C}_\Lambda$ の環写像とする。$f$ が全射で，
$\Ker(f)$ が $\mathfrak{m}_B$ により零化される非零単項イデアルであるとき，
これを {\it 小拡大}と呼ぶ。
\end{definition}

\noindent
次の補題により，$\mathcal{C}_\Lambda$ の全射環写像を含む議論は，しばしば
小拡大の場合に帰着できる。

\begin{lemma}
\label{formal-defos-lemma-factor-small-extension}
$f: B \to A$ を $\mathcal{C}_\Lambda$ の全射環写像とする。このとき $f$ は
小拡大の合成として分解できる。
\end{lemma}

\begin{proof}
$I$ を $f$ の核とする。$B$ は Artin なので，極大イデアル
$\mathfrak{m}_B$ は冪零である。$\mathfrak{m}_B^n = 0$ と書く。このとき
$f$ は分解
$$
B = B/I\mathfrak{m}_B^{n-1} \to B/I\mathfrak{m}_B^{n-2} \to
\ldots \to B/I \cong A
$$
をもち，その各写像は全射で，核は極大イデアルにより零化される。したがって，
$f$ 自身がそのような写像である場合，すなわち\ $I$ が $\mathfrak{m}_B$ により
零化される場合に補題を証明すれば十分である。この場合 $I$ は有限次元
$k$-ベクトル空間である。『代数』補題
\ref{algebra-lemma-artinian-finite-length} を参照せよ。
$I$ の $k$-ベクトル空間としての基底 $x_1, \ldots, x_n$ を取れば，$f$ は
$$
B \to B/(x_1) \to \ldots \to  B/(x_1, \ldots, x_n) \cong  A
$$
という小拡大の合成に分解される。
\end{proof}


\noindent
次の補題は，剰余体 $k$ をもつ局所 $\Lambda$-代数上の加群の長さを，
$\Lambda$ 上の長さによって計算できることを述べる。主張の記法を説明するため，
固定した有限環写像 $\Lambda \to k$ の像を $k' \subset k$ とする。
$k' \subset k$ は有限環拡大であることに注意せよ。したがって $k'$ は体であり，
$k/k'$ は有限体拡大である。『代数』補題
\ref{algebra-lemma-integral-under-field} を参照せよ。さらに
$\Lambda \to k'$ は全射なので，その核は極大イデアル
$\mathfrak m_\Lambda$ である。したがって
\begin{equation}
\label{formal-defos-equation-k-prime}
[k : k'] = [k : \Lambda/\mathfrak m_\Lambda] < \infty
\end{equation}
であり，古典的場合には $k = k'$ である。本章を通じて
$k' = \Lambda/\mathfrak m_\Lambda$ という記法を固定する。

\begin{lemma}
\label{formal-defos-lemma-length}
$A$ を剰余体 $k$ をもつ局所 $\Lambda$-代数とし，$M$ を $A$-加群とする。
このとき
$[k : k'] \text{length}_A(M) = \text{length}_\Lambda(M)$
である。古典的場合には
$\text{length}_A(M) = \text{length}_\Lambda(M)$
である。
\end{lemma}

\begin{proof}
$M$ が単純 $A$-加群なら，$A$-加群として $M \cong k$ である。『代数』補題
\ref{algebra-lemma-characterize-length-1} を参照せよ。この場合
$\text{length}_A(M) = 1$ であり，
% P10-E0420: preserve the frozen reversed field-degree token; correction is sidecar-only.
$\text{length}_\Lambda(M) = [k' : k]$ である。『代数』補題
\ref{algebra-lemma-dimension-is-length} を参照せよ。
$\text{length}_A(M)$ が有限なら，単純商をもつ $A$-部分加群による $M$ の
フィルトレーションを選び，加法性を用いれば結果を得る。『代数』補題
\ref{algebra-lemma-length-additive} を参照せよ。
$\text{length}_A(M)$ が無限なら，明らかな不等式
$\text{length}_A(M) \leq \text{length}_\Lambda(M)$ から結果を得る。
\end{proof}

\begin{lemma}
\label{formal-defos-lemma-surjective}
% P10-E0421: preserve the frozen unnamed map and its subsequent undefined f.
$A \to B$ を $\mathcal{C}_\Lambda$ の環写像とする。次は同値である。
\begin{enumerate}
\item $f$ は全射である。
\item $\mathfrak m_A/\mathfrak m_A^2 \to \mathfrak m_B/\mathfrak m_B^2$
は全射である。
\item $\mathfrak m_A/(\mathfrak m_\Lambda A + \mathfrak m_A^2)
\to \mathfrak m_B/(\mathfrak m_\Lambda B + \mathfrak m_B^2)$ は全射である。
\end{enumerate}
\end{lemma}

\begin{proof}
$\mathcal{C}_\Lambda$ の任意の環写像 $f : A \to B$ に対し，
$f(\mathfrak m_A) \subset \mathfrak m_B$ である。例えば，
$\mathfrak m_A$ と $\mathfrak m_B$ はそれぞれ $A$ と $B$ の冪零元全体だからである。
$f$ が全射であるとする。$y \in \mathfrak m_B$ とし，$f(x) = y$ となる
$x \in A$ を選ぶ。$f$ は同型
$A/\mathfrak m_A \to B/\mathfrak m_B$ を誘導するので，
$x \in \mathfrak m_A$ である。したがって誘導写像
$\mathfrak m_A/\mathfrak m_A^2 \to \mathfrak m_B/\mathfrak m_B^2$
は全射である。これにより (1) が (2) を含意することが分かる。

\medskip\noindent
(2) が (3) を含意することは明らかである。写像 $A \to B$ から，完全行をもつ
標準的な可換図式
$$
\xymatrix{
\mathfrak m_\Lambda/\mathfrak m_\Lambda^2 \otimes_{k'} k \ar[r] \ar[d] &
\mathfrak m_A/\mathfrak m_A^2 \ar[r] \ar[d] &
\mathfrak m_A/(\mathfrak m_\Lambda A + \mathfrak m_A^2) \ar[r] \ar[d] & 0 \\
\mathfrak m_\Lambda/\mathfrak m_\Lambda^2 \otimes_{k'} k \ar[r] &
\mathfrak m_B/\mathfrak m_B^2 \ar[r] &
\mathfrak m_B/(\mathfrak m_\Lambda B + \mathfrak m_B^2) \ar[r] & 0
}
$$
が得られる。したがって (3) が成り立てば (2) も成り立つ。

\medskip\noindent
(2) を仮定する。$A \to B$ が全射であることを示すには，Nakayama の補題
（『代数』補題 \ref{algebra-lemma-NAK}）により，
$A/\mathfrak m_A \to B/\mathfrak m_AB$ が全射であることを示せば十分である
（$\mathfrak m_A$ は冪零イデアルであることに注意せよ）。
$k = A/\mathfrak m_A = B/\mathfrak m_B$ なので，
$\mathfrak m_AB \to \mathfrak m_B$ が全射であることを示せば十分である。
もう一度 Nakayama の補題を適用すると，これは
$\mathfrak m_AB/\mathfrak m_A\mathfrak m_B \to \mathfrak m_B/\mathfrak m_B^2$
が全射であることを示せば十分であり，まさにこれを仮定している。
\end{proof}

\noindent
$A \to B$ が $\mathcal{C}_\Lambda$ の環写像なら，写像
$\mathfrak m_A/(\mathfrak m_\Lambda A + \mathfrak m_A^2)
\to \mathfrak m_B/(\mathfrak m_\Lambda B + \mathfrak m_B^2)$
は相対余接空間上の写像である。形式的な定義を与える。

\begin{definition}
\label{formal-defos-definition-tangent-space-ring}
$R \to S$ を局所環の局所準同型とする。
% P10-E0422: preserve the frozen reversed R-over-S wording; correction is sidecar-only.
$R$ の $S$ 上の {\it 相対余接空間}\footnote{注意：後に見るように，本章の
一般的な設定では，$\Lambda$ 上の対象 $A \in \mathcal{C}_\Lambda$ の接空間を，
相対余接空間の $k$-線形双対として単純に定義すべきではない。実際，相対余接空間の
正しい定義は $\Omega_{S/R} \otimes_S S/\mathfrak m_S$ である。}とは，
$S/\mathfrak m_S$-ベクトル空間
$\mathfrak m_S/(\mathfrak m_R S + \mathfrak m_S^2)$ のことである。
\end{definition}

\noindent
$f_1: A_1 \to A$ と $f_2: A_2 \to A$ を二つの環写像とすると，
ファイバー積 $A_1 \times_A A_2$ は，$A$ への二つの射影が等しい元からなる
$A_1 \times A_2$ の部分環である。本章を通じて，$f_1$ と $f_2$ が
$\mathcal{C}_\Lambda$ に属するとき，このようなファイバー積を含む条件を考える。
ファイバー積が常に $\mathcal{C}_\Lambda$ の対象となるわけではない。

\begin{example}
\label{formal-defos-example-fibre-product}
$p$ を素数，$n \in \mathbf{N}$ とする。
$\Lambda = \mathbf{F}_p(t_1, t_2, \ldots, t_n)$，
$k = \mathbf{F}_p(x_1, \ldots, x_n)$ とし，写像 $\Lambda \to k$ を
$t_i \mapsto x_i^p$ で与える。$A = k[\epsilon] = k[x]/(x^2)$ とする。
このとき $A$ は $\mathcal{C}_\Lambda$ の対象である。
$D : k \to k$ を $\Lambda$ 上の $k$ の導分，例えば
$D = \partial/\partial x_i$ とする。このとき写像
$$
f_D : k \longrightarrow k[\epsilon], \quad
a \mapsto a + D(a)\epsilon
$$
は $\mathcal{C}_\Lambda$ の射である。$A_1 = A_2 = k$ とし，
$f_1 = f_{\partial/\partial x_1}$ および $f_2(a) = a$ と置く。このとき
$A_1 \times_A A_2 = \{a \in k \mid \partial/\partial x_1(a) = 0\}$
であり，これは $k$ に全射ではない。したがってファイバー積は
$\mathcal{C}_\Lambda$ の対象ではない。
\end{example}


\noindent
この問題が生じうるのは，剰余体拡大 $k/k'$（\ref{formal-defos-equation-k-prime}）が
非分離であり，しかも $f_1$ も $f_2$ も全射でない場合に限ることが分かる。

\begin{lemma}
\label{formal-defos-lemma-fiber-product-CLambda}
$f_1 : A_1 \to A$ と $f_2 : A_2 \to A$ を
$\mathcal{C}_\Lambda$ の環写像とする。このとき次が成り立つ。
\begin{enumerate}
\item $f_1$ または $f_2$ が全射なら，
$A_1 \times_A A_2$ は $\mathcal{C}_\Lambda$ に属する。
\item $f_2$ が小拡大なら，$A_1 \times_A A_2 \to A_1$ も小拡大である。
\item 体拡大 $k/k'$ が分離的なら，
$A_1 \times_A A_2$ は $\mathcal{C}_\Lambda$ に属する。
\end{enumerate}
\end{lemma}

\begin{proof}
環 $A_1 \times_A A_2$ は，写像 $\Lambda \to A_1$ と
$\Lambda \to A_2$ から誘導される写像
$\Lambda \to A_1 \times_A A_2$ によって $\Lambda$-代数となる。これは唯一の
極大イデアル
$$
\mathfrak m_{A_1} \times_{\mathfrak m_A} \mathfrak m_{A_2} =
\Ker(A_1 \times_A A_2 \longrightarrow k)
$$
をもつ局所環である。環が Artin であることと，自身の上の加群として有限長で
あることとは同値である。『代数』補題
\ref{algebra-lemma-artinian-finite-length} を参照せよ。
$A_1$ と $A_2$ は Artin なので，補題 \ref{formal-defos-lemma-length} により
$\text{length}_\Lambda(A_1)$ と $\text{length}_\Lambda(A_2)$，したがって
$\text{length}_\Lambda(A_1 \times A_2)$ はすべて有限である。
$A_1 \times_A A_2 \subset A_1 \times A_2$ は $\Lambda$-部分加群なので，
これは
$\text{length}_{A_1 \times_A A_2}(A_1 \times_A A_2) \leq
\text{length}_\Lambda(A_1 \times_A A_2)$
が有限であることを含意する。ゆえに $A_1 \times_A A_2$ は Artin である。
したがって $A_1 \times_A A_2$ が $\mathcal{C}_\Lambda$ の対象となるのを
妨げうるのは，上の写像
$A_1 \times_A A_2 \to A \to A/\mathfrak m_A = k$ によって，その剰余体が
$k$ の真部分体へ写る可能性だけである。

\medskip\noindent
(1) の証明。$f_2$ が全射なら，射影
$A_1 \times_A A_2 \to A_1$ は全射である。したがって合成
$A_1 \times_A A_2 \to A_1 \to A_1/\mathfrak m_{A_1} = k$ は全射であり，
$A_1 \times_A A_2$ は $\mathcal{C}_\Lambda$ の対象である。

\medskip\noindent
(2) の証明。$f_2$ が小拡大なら，$A_2 \to A$ と
$A_1 \times_A A_2  \to A_1$ はいずれも全射で，同じ核をもつ。したがって
$A_1 \times_A A_2  \to A_1$ の核は $1$-次元 $k$-ベクトル空間であり，
$A_1 \times_A A_2  \to A_1$ は小拡大である。

\medskip\noindent
(3) の証明。$k = k'(\overline{x})$ となる $\overline{x} \in k$ を選ぶ
（『体』補題 \ref{fields-lemma-primitive-element} を参照せよ）。
$P'(T) \in k'[T]$ を $k'$ 上の $\overline{x}$ の最小多項式とする。
$k/k'$ は分離的なので，
$\text{d}P/\text{d}T(\overline{x}) \not = 0$ である。全射
$\Lambda[T] \to k'[T]$ により $P'$ に写るモニック多項式
$P \in \Lambda[T]$ を選ぶ。$A, A_1, A_2$ は Hensel 環なので，
『代数』補題 \ref{algebra-lemma-local-dimension-zero-henselian} により，
$P(x) = 0, P(x_1) = 0, P(x_2) = 0$ を満たし，$k$ における
$x, x_1, x_2$ の像が $\overline{x}$ となるような
$x, x_1, x_2 \in A, A_1, A_2$ が取れる。一意性（『代数』補題
\ref{algebra-lemma-uniqueness}）により $x_1, x_2$ は $x \in A$ に
写るので，$(x_1, x_2) \in A_1 \times_A A_2$ である。したがって
$A_1 \times_A A_2$ の剰余体は $k$ の $k'$ 上の生成元を含み，結論を得る。
\end{proof}

\noindent
次に $\mathcal{C}_\Lambda$ における本質的全射を定義する。
$\mathcal{C}_\Lambda$ の全射が本質的であるための必要十分条件は補題
\ref{formal-defos-lemma-essential-surjection} で与える。

\begin{definition}
\label{formal-defos-definition-essential-surjection}
$f: B \to A$ を $\mathcal{C}_\Lambda$ の環写像とする。$f$ が次の性質を
もつとき，これを {\it 本質的全射}と呼ぶ。
\begin{enumerate}
\item $f$ は全射である。
\item $g: C \to B$ が $\mathcal{C}_\Lambda$ の環写像で，
$f \circ g$ が全射なら，$g$ も全射である。
\end{enumerate}
\end{definition}

\noindent
補題 \ref{formal-defos-lemma-surjective} を用いると，$\mathcal{C}_\Lambda$ における
本質的全射を次のように特徴付けられる。

\begin{lemma}
\label{formal-defos-lemma-essential-surjection-mod-squares}
$f: B \to A$ を $\mathcal{C}_\Lambda$ の環写像とする。次は同値である。
\begin{enumerate}
\item $f$ は本質的全射である。
\item 写像 $B/\mathfrak m_B^2 \to A/\mathfrak m_A^2$ は本質的全射である。
\item 写像
$B/(\mathfrak m_\Lambda B + \mathfrak m_B^2) \to
A/(\mathfrak m_\Lambda A + \mathfrak m_A^2)$ は本質的全射である。
\end{enumerate}
\end{lemma}

\begin{proof}
(3) を仮定する。$C \to B$ を $\mathcal{C}_\Lambda$ の環写像とし，
$C \to A$ は全射であるとする。このとき
$C \to A/(\mathfrak m_\Lambda A + \mathfrak m_A^2)$ も全射である。
仮定により $C \to B/(\mathfrak m_\Lambda B + \mathfrak m_B^2)$ も全射である。
したがって補題 \ref{formal-defos-lemma-surjective} を二度適用すると，$C \to B$ は全射である。

\medskip\noindent
(1) を仮定する。$C \to B/(\mathfrak m_\Lambda B + \mathfrak m_B^2)$ を
$\mathcal{C}_\Lambda$ の射とし，
$C \to A/(\mathfrak m_\Lambda A + \mathfrak m_A^2)$ は全射であるとする。
補題 \ref{formal-defos-lemma-fiber-product-CLambda} により $\mathcal{C}_\Lambda$ の対象となる
$C' = C \times_{B/(\mathfrak m_\Lambda B + \mathfrak m_B^2)} B$
を置く。$C' \to A/(\mathfrak m_\Lambda A + \mathfrak m_A^2)$ は依然として
全射であることに注意せよ。したがって補題 \ref{formal-defos-lemma-surjective} により
$C' \to A$ は全射である。ゆえに仮定から $C' \to B$ は全射である。
これは $C' \to B/(\mathfrak m_\Lambda B + \mathfrak m_B^2)$ が全射である
ことを含意し，$C'$ の構成から
$C \to B/(\mathfrak m_\Lambda B + \mathfrak m_B^2)$ が全射であることを含意する。

\medskip\noindent
第一段落で (3) $\Rightarrow$ (1) を，第二段落で (1) $\Rightarrow$ (3) を
証明した。(2) と (3) の同値性は (1) と (3) の同値性の特別な場合なので，
証明は完了した。
\end{proof}

\noindent
$\mathcal{C}_\Lambda$ における本質的全射をさらに解析するため，記法を導入する。
$A$ を $\mathcal{C}_\Lambda$ の対象，またはより一般に，$\Lambda$-代数全射
$A \to k$ を備えた任意の $\Lambda$-代数とする。標準的な完全列
\begin{equation}
\label{formal-defos-equation-sequence}
\mathfrak m_A/\mathfrak m_A^2 \xrightarrow{\text{d}_A}
\Omega_{A/\Lambda} \otimes_A k \to
\Omega_{k/\Lambda} \to 0
\end{equation}
がある。『代数』補題 \ref{algebra-lemma-differential-seq} を参照せよ。
（\ref{formal-defos-equation-k-prime}）の $k'$ に対し
$\Omega_{k/\Lambda} = \Omega_{k/k'}$ であることに注意せよ。
$H_1(L_{k/\Lambda})$ を，$\Lambda$ 上の $k$ の素朴な余接複体の第一ホモロジー
加群とする。『代数』定義 \ref{algebra-definition-naive-cotangent-complex} を
参照せよ。このとき（\ref{formal-defos-equation-sequence}）は完全列
\begin{equation}
\label{formal-defos-equation-sequence-extended}
H_1(L_{k/\Lambda}) \to
\mathfrak m_A/\mathfrak m_A^2 \xrightarrow{\text{d}_A}
\Omega_{A/\Lambda} \otimes_A k \to
\Omega_{k/\Lambda} \to 0,
\end{equation}
へ延長できる。『代数』補題 \ref{algebra-lemma-exact-sequence-NL} を参照せよ。
$B \to A$ を $\mathcal{C}_\Lambda$ の環写像，またはより一般に，$k$ への
$\Lambda$-代数全射を備えた $\Lambda$-代数の写像とすると，完全行をもつ
可換図式
\begin{equation}
\label{formal-defos-equation-sequence-functorial}
\vcenter{
\xymatrix{
H_1(L_{k/\Lambda}) \ar[r] \ar@{=}[d] &
\mathfrak m_B/\mathfrak m_B^2 \ar[r]_{\text{d}_B} \ar[d] &
\Omega_{B/\Lambda} \otimes_B k \ar[r] \ar[d] &
\Omega_{k/\Lambda} \ar[r] \ar@{=}[d] & 0 \\
H_1(L_{k/\Lambda}) \ar[r] &
\mathfrak m_A/\mathfrak m_A^2 \ar[r]^{\text{d}_A} &
\Omega_{A/\Lambda} \otimes_A k \ar[r] &
\Omega_{k/\Lambda} \ar[r] & 0
}
}
\end{equation}
が得られる。

\begin{lemma}
\label{formal-defos-lemma-H1-separable-case}
標準写像
$$
\mathfrak m_\Lambda/\mathfrak m_\Lambda^2 \longrightarrow H_1(L_{k/\Lambda}).
$$
が存在する。$k' \subset k$ が分離的なら（例えば $k$ の標数が零なら），
この写像は同型
$\mathfrak m_\Lambda/\mathfrak m_\Lambda^2 \otimes_{k'} k = H_1(L_{k/\Lambda})$
を誘導する。$k = k'$ なら（例えば古典的場合には），
$\mathfrak m_\Lambda/\mathfrak m_\Lambda^2 = H_1(L_{k/\Lambda})$ である。合成
$$
\mathfrak m_\Lambda/\mathfrak m_\Lambda^2 \longrightarrow
H_1(L_{k/\Lambda}) \longrightarrow \mathfrak m_A/\mathfrak m_A^2
$$
は標準写像 $\mathfrak m_\Lambda \to \mathfrak m_A$ から生じる。
\end{lemma}

\begin{proof}
$\Lambda \to k'$ は核 $\mathfrak m_\Lambda$ をもつ全射なので，
$H_1(L_{k'/\Lambda}) = \mathfrak m_\Lambda/\mathfrak m_\Lambda^2$ であることに
注意せよ。写像は素朴な余接複体の関手性から生じる。
$k' \subset k$ が分離的なら，$k' \to k$ は \'etale 環写像である。『代数』補題
\ref{algebra-lemma-etale-over-field} を参照せよ。したがって，その素朴な余接複体の
ホモロジー群は自明である。『代数』定義 \ref{algebra-definition-etale} を参照せよ。
環写像 $\Lambda \to k' \to k$ に『代数』補題
\ref{algebra-lemma-exact-sequence-NL} を適用すると，
$\mathfrak m_\Lambda/\mathfrak m_\Lambda^2 \otimes_{k'} k = H_1(L_{k/\Lambda})$
を得る。最後の主張の証明は省略する。
\end{proof}

\begin{lemma}
\label{formal-defos-lemma-essential-surjection}
$f: B \to A$ を $\mathcal{C}_\Lambda$ の環写像とし，記法は
（\ref{formal-defos-equation-sequence-functorial}）のとおりとする。
\begin{enumerate}
\item $f$ が全射であることを特徴付ける補題 \ref{formal-defos-lemma-surjective} の
同値な条件は，次の各条件とも同値である。
\begin{enumerate}
\item $\Im(\text{d}_B) \to \Im(\text{d}_A)$ は全射である。
\item 写像 $\Omega_{B/\Lambda} \otimes_B k \to
\Omega_{A/\Lambda} \otimes_A k$ は全射である。
\end{enumerate}
\item 次は同値である。
\begin{enumerate}
\item $f$ は本質的全射である
（補題 \ref{formal-defos-lemma-essential-surjection-mod-squares} を参照）。
\item 写像 $\Im(\text{d}_B) \to \Im(\text{d}_A)$ は同型である。
\item 写像 $\Omega_{B/\Lambda} \otimes_B k \to
\Omega_{A/\Lambda} \otimes_A k$ は同型である。
\end{enumerate}
\item $k/k'$ が分離的なら，$f$ が本質的全射であることと，写像
$\mathfrak m_B/(\mathfrak m_\Lambda B + \mathfrak m_B^2) \to
\mathfrak m_A/(\mathfrak m_\Lambda A + \mathfrak m_A^2)$
が同型であることとは同値である。
\item $f$ が小拡大なら，$f$ が本質的でないことと，
$f \circ s = \text{id}_A$ を満たす $\mathcal{C}_\Lambda$ の切断
$s: A \to B$ を $f$ がもつこととは同値である。
\end{enumerate}
\end{lemma}

\begin{proof}
(1) の証明。（\ref{formal-defos-equation-sequence-functorial}）から (1)(a) と (1)(b) は
同値である。また，
% P10-E0423: preserve the frozen reversed A-to-B surjectivity assertion.
$A \to B$ が全射なら，(1)(a) と (1)(b) は成り立つ。(1)(a) を仮定する。
$\text{d}_A$ の核は $H_1(L_{k/\Lambda})$ の像であり，これは
$\mathfrak m_B/\mathfrak m_B^2$ にも写るので，
$\mathfrak m_B/\mathfrak m_B^2 \to \mathfrak m_A/\mathfrak m_A^2$
は全射である。ゆえに補題 \ref{formal-defos-lemma-surjective} により $B \to A$ は全射である。
これで (1) の証明が完了した。

\medskip\noindent
(2) の証明。(2)(b) と (2)(c) の同値性は
（\ref{formal-defos-equation-sequence-functorial}）から直ちに従う。

\medskip\noindent
(2)(b) を仮定する。$g : C \to B$ を $\mathcal{C}_\Lambda$ の環写像とし，
$f \circ g$ は全射であるとする。補題 \ref{formal-defos-lemma-surjective} により
$\mathfrak m_C/\mathfrak m_C^2 \to \mathfrak m_A/\mathfrak m_A^2$ は全射である。
したがって $\Im(\text{d}_C) \to \Im(\text{d}_A)$ は全射であり，仮定から
$\Im(\text{d}_C) \to \Im(\text{d}_B)$ も全射である。(1) により
$C \to B$ は全射である。

\medskip\noindent
(2)(a) を仮定する。このとき $f$ は全射であり，
$\Omega_{B/\Lambda} \otimes_B k \to \Omega_{A/\Lambda} \otimes_A k$
は全射である。その核を $K$ とする。（\ref{formal-defos-equation-sequence-functorial}）により
$K = \text{d}_B(\Ker(\mathfrak m_B/\mathfrak m_B^2 \to
\mathfrak m_A/\mathfrak m_A^2))$ である。$k$-ベクトル空間の分裂
$$
\Omega_{B/\Lambda} \otimes_B k =
\Omega_{A/\Lambda} \otimes_A k \oplus K
$$
を選ぶ。写像 $\text{d} : B \to \Omega_{B/\Lambda}$ は，$K$ への射影を介して
写像 $D : B \to K$ を誘導する。
$C = \{b \in B \mid D(b) = 0\}$ と置く。Leibniz 則から，これは $B$ の
$\Lambda$-部分代数である。$\overline{x} \in k$ とし，$\overline{x}$ に写る
$x \in B$ を選ぶ。$D(x) \not = 0$ なら，$D(y) = D(x)$ を満たす元
$y \in \mathfrak m_B$ が取れる。したがって $x - y \in C$ は
$\overline{x}$ に写る元である。ゆえに $C \to k$ は全射で，$C$ は
$\mathcal{C}_\Lambda$ の対象である。同様に
$\omega \in \Im(\text{d}_A)$ を取る。(1) により，
$\text{d}_B(x)$ が $\omega$ に写るような $x \in \mathfrak m_B$ を取れる。
$D(x) \not = 0$ なら，$\mathfrak m_A/\mathfrak m_A^2$ で零に写り，
$D(y) = D(x)$ を満たす元 $y \in \mathfrak m_B$ が取れる。したがって
$z = x - y$ は $\mathfrak m_C$ の元で，その像
% P10-E0424: preserve the frozen C-over-k differential module.
$\text{d}_C(z) \in \Omega_{C/k} \otimes_C k$ は $\omega$ に写る。
ゆえに $\Im(\text{d}_C) \to \Im(\text{d}_A)$ は全射である。(1) により
$C \to A$ は全射である。したがって仮定から $C \to B$ は全射である。
ゆえに $D = 0$，すなわち $K = 0$，すなわち (2)(c) が成り立つ。
これで (2) の証明が完了した。

\medskip\noindent
(3) の証明。
% P10-E0425: preserve the frozen reversed separability extension k'/k.
$k'/k$ が分離的なら，補題 \ref{formal-defos-lemma-H1-separable-case} により
$H_1(L_{k/\Lambda}) =
\mathfrak m_\Lambda/\mathfrak m_\Lambda^2 \otimes_{k'} k$ である。
したがって $\Im(\text{d}_A) =
\mathfrak m_A/(\mathfrak m_\Lambda A + \mathfrak m_A^2)$ であり，$B$ についても
同様である。ゆえに (3) は (2) から従う。

\medskip\noindent
(4) の証明。$f$ の切断 $s$ は全射ではない（定義により小拡大の核は非自明で
ある）ので，$f$ は本質的全射ではない。逆に，$f$ は小拡大だが本質的全射で
ないとする。全射ではない $\mathcal{C}_\Lambda$ の環写像 $C \to B$ で，
$C \to A$ が全射となるものを選ぶ。$C' \subset B$ を $C \to B$ の像とする。
このとき $C' \not = B$ だが，$C'$ は $A$ に全射である。
$f : B \to A$ は小拡大なので，
$\text{length}_C(B) = \text{length}_C(A) + 1$ である。
$C'$ は $B$ の真部分環なので
$\text{length}_C(C') \leq \text{length}_C(A)$ である。しかし
$C' \to A$ は全射だから，実際には
$\text{length}_C(C') = \text{length}_C(A)$ であり，$C' \to A$ は同型である。
これが求める切断を与える。
\end{proof}

\begin{example}
\label{formal-defos-example-essential-surjection}
$\Lambda = k[[x]]$ を $k$ 上 $1$ 変数の冪級数環とする。$A = k$，
$B = \Lambda/(x^2)$ と置く。補題 \ref{formal-defos-lemma-essential-surjection} により
$B \to A$ は本質的全射である。実際，これは小拡大で，写像 $B \to A$ は
（圏 $\mathcal{C}_\Lambda$ において）右逆射をもたない。しかし写像
$$
k \cong \mathfrak m_B/\mathfrak m_B^2
\longrightarrow
\mathfrak m_A/\mathfrak m_A^2 = 0
$$
は同型ではない。したがって補題 \ref{formal-defos-lemma-essential-surjection} (3) では，
相対余接空間の写像
$\mathfrak m_B/(\mathfrak m_\Lambda B + \mathfrak m_B^2) \to
\mathfrak m_A/(\mathfrak m_\Lambda A + \mathfrak m_A^2)$
を考える必要がある。
\end{example}

\section{完備化された基底圏}
\label{formal-defos-section-category-completion-CLambda}

\noindent
以下で定める圏 $\mathcal{C}_\Lambda$ の「完備化」は，
$\mathcal{C}_\Lambda$ 上で群亜群をファイバーとする余ファイバー圏の完備化に
対する基底圏として働く（第 \ref{formal-defos-section-formal-objects} 節）。

\begin{definition}
\label{formal-defos-definition-completion-CLambda}
$\Lambda$ を Noether 環，$\Lambda \to k$ を有限環写像とし，$k$ は体であると
する。{\it $\widehat{\mathcal{C}}_\Lambda$} を次の圏として定義する。
\begin{enumerate}
\item 対象は対 $(R, \varphi)$ であって，$R$ は Noether 完備局所
$\Lambda$-代数，$\varphi : R/\mathfrak m_R \to k$ は
$\Lambda$-代数同型である。
\item 射 $f : (S, \psi) \to (R, \varphi)$ は，
$\varphi \circ (f \bmod \mathfrak m) = \psi$ を満たす局所
$\Lambda$-代数準同型である。
\end{enumerate}
\end{definition}

\noindent
定義 \ref{formal-defos-definition-CLambda} に続く議論と同様に，通常は
$\widehat{\mathcal{C}}_\Lambda$ の対象を単に $R$ と表し，同一視
$R/\mathfrak m_R = k$ は暗黙に了解する。本節では，前節における
$\mathcal{C}_\Lambda$ の議論と並行して，圏
$\widehat{\mathcal{C}}_\Lambda$ の対象と射の基本的性質を論じる。

\medskip\noindent
最初に，任意の対象 $A \in \mathcal{C}_\Lambda$ は
$\widehat{\mathcal{C}}_\Lambda$ の対象でもあることに注意する。実際，Artin
局所環は常に Noether であり，その極大イデアル（結局これは冪零イデアルで
ある）に関して完備である。さらに定義から明らかなように，
$\mathcal{C}_\Lambda \subset \widehat{\mathcal{C}}_\Lambda$
は Artin 環全体からなる厳密充満部分圏である。逆に，
$\widehat{\mathcal{C}}_\Lambda$ の任意の対象は
$\mathcal{C}_\Lambda$ の対象の極限であることが分かる。

\medskip\noindent
$R$ を $\widehat{\mathcal{C}}_\Lambda$ の対象とする。
$n \in \mathbf{N}$ に対して環 $R_n = R/\mathfrak m_R^n$ を考える。
これらは唯一の冪零素イデアルをもつ Noether 局所環であり，したがって
Artin である。『代数』命題
\ref{algebra-proposition-dimension-zero-ring} を参照せよ。環写像
$$
\ldots \to R_{n + 1} \to R_n \to \ldots \to R_2 \to R_1 = k
$$
はすべて全射である。定義により，$R$ の完備性は $R = \lim R_n$ を意味する。
$f : R \to S$ が $\widehat{\mathcal{C}}_\Lambda$ の環写像なら，与えられた
写像を極限にもつ環写像の系 $f_n : R_n \to S_n$ が得られる。

\begin{lemma}
\label{formal-defos-lemma-surjective-cotangent-space}
$f: R \to S$ を $\widehat{\mathcal{C}}_\Lambda$ の環写像とする。次は同値である。
\begin{enumerate}
\item $f$ は全射である。
\item 写像
$\mathfrak m_R/\mathfrak m_R^2 \to \mathfrak m_S/\mathfrak m_S^2$
は全射である。
\item 写像
$\mathfrak m_R/(\mathfrak m_\Lambda R + \mathfrak m_R^2) \to
\mathfrak m_S/(\mathfrak m_\Lambda S + \mathfrak m_S^2)$
は全射である。
\end{enumerate}
\end{lemma}

\begin{proof}
$n \geq 2$ に対し，相対余接空間の等式
$$
\mathfrak m_R/(\mathfrak m_\Lambda R + \mathfrak m_R^2)
=
\mathfrak m_{R_n}/(\mathfrak m_\Lambda R_n + \mathfrak m_{R_n}^2)
$$
があり，$S$ についても同様であることに注意せよ。したがって補題
\ref{formal-defos-lemma-surjective} により，すべての $n$ に対して $R_n \to S_n$ は
全射である。$K_n$ を $R_n \to S_n$ の核とする。このとき列
$$
0 \to K_n \to R_n \to S_n \to 0
$$
は有向逆系の完全列をなす。各 $K_n$ は Artin なので，系 $(K_n)$ は
Mittag--Leffler である。したがって『代数』補題
\ref{algebra-lemma-ML-exact-sequence} により，極限を取っても完全性は保たれる。
ゆえに $\lim R_n \to \lim S_n$ は全射，すなわち $f$ は全射である。
\end{proof}

\begin{lemma}
\label{formal-defos-lemma-CLambdahat-pushouts}
圏 $\widehat{\mathcal{C}}_\Lambda$ は押し出しをもつ。
\end{lemma}

\begin{proof}
$R \to S_1$ と $R \to S_2$ を $\widehat{\mathcal{C}}_\Lambda$ の射とする。
環 $C = S_1 \otimes_R S_2$ を考える。この環は有限生成極大イデアル
$\mathfrak m = \mathfrak m_{S_1} \otimes S_2 +
S_1 \otimes \mathfrak m_{S_2}$ をもち，その剰余体は $k$ である。
$C^\wedge$ を $\mathfrak m$ に関する $C$ の完備化とする。このとき
$C^\wedge$ は，極大イデアル
$\mathfrak m^\wedge = \mathfrak mC^\wedge$ に関して完備な Noether 環で，
その剰余体は $k$ と同一視される。『代数』補題
\ref{algebra-lemma-completion-Noetherian} を参照せよ。したがって
$C^\wedge$ は $\widehat{\mathcal{C}}_\Lambda$ の対象である。
そして $S_1 \to C^\wedge$ と $S_2 \to C^\wedge$ により，$C^\wedge$ は
$\widehat{\mathcal{C}}_\Lambda$ における $R$ 上の押し出しとなる
（詳細は省略する）。
\end{proof}

\noindent
次の補題は本章では用いない。

\begin{lemma}
\label{formal-defos-lemma-CLambdahat-coproducts}
圏 $\widehat{\mathcal{C}}_\Lambda$ は対象の対の余積をもつ。
\end{lemma}

\begin{proof}
$R$ と $S$ を $\widehat{\mathcal{C}}_\Lambda$ の対象とする。
環 $C = R \otimes_\Lambda S$ を考える。標準的な全射
$C \to R \otimes_\Lambda S \to k \otimes_\Lambda k \to k$
があり，最後の写像は乗法写像である。$C \to k$ の核は極大イデアル
$\mathfrak m$ である。$\mathfrak m$ は $\mathfrak m_R C$，
$\mathfrak m_S C$，および $k \otimes_\Lambda k \to k$ の核の生成元に
写る有限個の $C$ の元によって生成されることに注意せよ。したがって
$\mathfrak m$ は有限生成イデアルである。$C^\wedge$ を $\mathfrak m$ に
関する $C$ の完備化とする。このとき $C^\wedge$ は，剰余体 $k$ をもち，
極大イデアル $\mathfrak m^\wedge = \mathfrak mC^\wedge$ に関して完備な
Noether 環である。『代数』補題
\ref{algebra-lemma-completion-Noetherian} を参照せよ。したがって
$C^\wedge$ は $\widehat{\mathcal{C}}_\Lambda$ の対象である。
そして $R \to C^\wedge$ と $S \to C^\wedge$ により，$C^\wedge$ は
$\widehat{\mathcal{C}}_\Lambda$ における余積となる（詳細は省略する）。
\end{proof}

\noindent
圏における空余積はその圏の始対象である。古典的場合には
$\widehat{\mathcal{C}}_\Lambda$ は始対象，すなわち $\Lambda$ 自身をもつ。
より一般に $k' = k$ なら，$\mathfrak m_\Lambda$ に関する $\Lambda$ の
完備化 $\Lambda^\wedge$ が始対象である。さらに一般に，
$k' \subset k$ が分離的なら，$\widehat{\mathcal{C}}_\Lambda$ はやはり
始対象をもつ。実際，モニック多項式 $P \in \Lambda[T]$ で，
% P10-E0426: preserve the frozen lowercase p' image token.
$k \cong k'[T]/(P')$ かつ $p' \in k'[T]$ が $P$ の像となるものを選ぶ。
このとき $R = \Lambda^\wedge[T]/(P)$ は始対象である。補題
\ref{formal-defos-lemma-fiber-product-CLambda} の証明を参照せよ。

\medskip\noindent
上のような始対象 $R$ があれば，
$\mathcal{C}_\Lambda = \mathcal{C}_R$ および
$\widehat{\mathcal{C}}_\Lambda = \widehat{\mathcal{C}}_R$ であり，本章の議論
全体は事実上古典的場合に戻る。しかし $k' \subset k$ が非分離なら，
始対象は存在しない。

\begin{lemma}
\label{formal-defos-lemma-derivations-finite}
$S$ を $\widehat{\mathcal{C}}_\Lambda$ の対象とする。このとき
$\dim_k \text{Der}_\Lambda(S, k) < \infty$ である。
\end{lemma}

\begin{proof}
$x_1, \ldots, x_n \in \mathfrak m_S$ を，相対余接空間
$\mathfrak m_S/(\mathfrak m_\Lambda S + \mathfrak m_S^2)$ の
$k$-基底に写るように選ぶ。$y_1, \ldots, y_m \in S$ を，その $k$ における
像が $k'$ 上 $k$ を生成するように選ぶ。
$\dim_k \text{Der}_\Lambda(S, k) \leq n + m$ であると主張する。これを見るには，
$D(x_i) = 0$ かつ $D(y_j) = 0$ なら $D = 0$ であることを証明すれば十分である。
$a \in S$ とする。$k$ における像が $a$ の $k$ における像と同じになるような多項式
$P = \sum \lambda_J y^J$，$\lambda_J \in \Lambda$ を取れる。このとき，
すべての $j$ に対して $D(y_j) = 0$ という仮定により
$D(a - P) = D(a) - D(P) = D(a)$ である。したがって
$a \in \mathfrak m_S$ と仮定してよい。$a_i \in S$ を用いて
$a = \sum a_i x_i$ と書く。Leibniz 則により
$$
D(a) = \sum x_iD(a_i) + \sum a_iD(x_i) = \sum x_iD(a_i)
$$
である。ここで $D(x_i) = 0$ を仮定した。$k$ 上では $x_i$ による乗法が
零なので，$\sum x_iD(a_i) = 0$ である。
\end{proof}

\begin{lemma}
\label{formal-defos-lemma-derivations-surjective}
$f : R \to S$ を $\widehat{\mathcal{C}}_\Lambda$ の射とする。
$\text{Der}_\Lambda(S, k) \to \text{Der}_\Lambda(R, k)$ が単射なら，
$f$ は全射である。
\end{lemma}

\begin{proof}
$f$ が全射でないなら，補題 \ref{formal-defos-lemma-surjective-cotangent-space} により
$\mathfrak m_S/(\mathfrak m_R S + \mathfrak m_S^2)$ は非零である。
したがって $Q = S/(f(R) + \mathfrak m_R S + \mathfrak m_S^2)$ も非零である。
$Q$ は $f$ を介して $k = R/\mathfrak m_R$-ベクトル空間となることに注意せよ。
$S \to k$ を介して $Q$ を $S$-加群とする。商写像 $D : S \to Q$ は
$R$-導分である。実際，$a_1, a_2 \in S$ なら，ある
$b_1, b_2 \in R$ および $a_1', a_2' \in \mathfrak m_S$ を用いて
$a_1 = f(b_1) + a_1'$，$a_2 = f(b_2) + a_2'$ と書ける。
$i = 1, 2$ に対して $b_i$ と $a_i$ は $k$ において同じ像をもち，
\begin{align*}
a_1a_2 & = (f(b_1) + a_1')(f(b_2) + a_2') \\
& = f(b_1)a_2' + f(b_2)a_1' \\
& = f(b_1)(f(b_2) + a_2') + f(b_2)(f(b_1) + a_1') \\
& = f(b_1)a_2 + f(b_2)a_1
\end{align*}
が $Q$ において成り立つ。これは Leibniz 則を証明する。したがって
$D : S \to Q$ は，$R \to S$ との合成が零となる $\Lambda$-導分である。
$Q \not = 0$ なので，$R \to S$ との合成が零となる導分 $D : S \to k$ も
存在する。すなわち，
$\text{Der}_\Lambda(S, k) \to \text{Der}_\Lambda(R, k)$ は単射ではない。
\end{proof}

\begin{lemma}
\label{formal-defos-lemma-m-adic-topology}
$R$ を $\widehat{\mathcal{C}}_\Lambda$ の対象とする。$(J_n)$ を
$\mathfrak m_R^n \subset J_n$ を満たすイデアルの減少列とする。
$J = \bigcap J_n$ と置く。このとき列 $(J_n/J)$ は $R/J$ 上の
$\mathfrak m_{R/J}$-進位相を定める。
\end{lemma}

\begin{proof}
$\mathfrak m_{R/J}^n \subset J_n/J$ は明らかである。したがって，すべての
$n$ に対して $J_N/J \subset \mathfrak m_{R/J}^n$ となる $N$ が存在することを
示せば十分である。これは $J_N \subset \mathfrak m_R^n + J$ と同値である。
各 $n$ に対して環 $R/\mathfrak m_R^n$ は Artin なので，ある $N_n$ が存在して
$$
J_{N_n} + \mathfrak m_R^n = J_{N_n + 1} + \mathfrak m_R^n = \ldots
$$
となる。$E_n = (J_{N_n} + \mathfrak m_R^n)/\mathfrak m_R^n$ と置き，
$E = \lim E_n \subset \lim R/\mathfrak m_R^n = R$ と置く。
$E \subset J$ であることに注意せよ。実際，任意の $f \in E$ と任意の $m$ に
対して，すべての $n \gg 0$ について
$f \in J_m + \mathfrak m_R^n$ であるから，Krull の交叉定理により
$f \in J_m$ である。『代数』補題
\ref{algebra-lemma-intersect-powers-ideal-module-zero} を参照せよ。
遷移写像 $E_n \to E_{n - 1}$ はすべて全射なので，$J$ は $E_n$ に全射である。
% P10-E0427: the localized conclusion is grammatical; the frozen English defect remains sidecar-only.
したがって $N = N_n$ と取ればよい。
\end{proof}

\begin{lemma}
\label{formal-defos-lemma-limit-artinian}
$\ldots \to A_3 \to A_2 \to A_1$ を $\mathcal{C}_\Lambda$ における全射環写像の
列とする。$\dim_k (\mathfrak m_{A_n}/\mathfrak m_{A_n}^2)$ が有界なら，
$S = \lim A_n$ は $\widehat{\mathcal{C}}_\Lambda$ の対象であり，イデアル
$I_n = \Ker(S \to A_n)$ は $S$ 上の $\mathfrak m_S$-進位相を定める。
\end{lemma}

\begin{proof}
すべての $n$ に対して写像 $S \to A_n$ が全射であることを自由に用いる。
補題 \ref{formal-defos-lemma-surjective-cotangent-space} により，写像
$\mathfrak m_{A_{n + 1}}/\mathfrak m_{A_{n + 1}}^2 \to
\mathfrak m_{A_n}/\mathfrak m_{A_n}^2$ は全射であることに注意せよ。
したがって，十分大きな $n$ に対して次元
$\dim_k (\mathfrak m_{A_n}/\mathfrak m_{A_n}^2)$ はある整数，例えば $r$ に
安定する。ゆえに $x_1, \ldots, x_r \in \mathfrak m_S$ で，その $A_n$ に
おける像が $\mathfrak m_{A_n}$ を生成するものを取れる。さらに
$y_1, \ldots, y_t \in S$ を，その $k$ における像が $\Lambda$ 上 $k$ を
生成するように取る。このとき環写像
$P = \Lambda[z_1, \ldots, z_{r + t}] \to S$，$z_i \mapsto x_i$ および
$z_{r + j} \mapsto y_j$ が得られ，その合成 $P \to S \to A_n$ はすべての
$n$ に対して全射である。$\mathfrak m \subset P$ を $P \to k$ の核とする。
$R = P^\wedge$ を $P$ の $\mathfrak m$-進完備化とする。これは
$\widehat{\mathcal{C}}_\Lambda$ の対象である。両立する全射の系
$R \to A_n$ が依然としてあるので，写像 $R \to S$ を得る。
$J_n = \Ker(R \to A_n)$，$J = \bigcap J_n$ と置く。補題
\ref{formal-defos-lemma-m-adic-topology} により
$R/J = \lim R/J_n = \lim A_n = S$ であり，イデアル $J_n/J = I_n$ は
$\mathfrak m$-進位相を定める（各 $n$ に対して，ある $N_n$ について
$\mathfrak m_R^{N_n} \subset J_n$ であるが，必ずしも $N_n = n$ ではない。
したがって補題を適用する前にイデアル $J_n$ の番号を付け直す必要がありうる）。
\end{proof}

\begin{lemma}
\label{formal-defos-lemma-power-series}
$R', R \in \Ob(\widehat{\mathcal{C}}_\Lambda)$ とする。$R$ のあるイデアル $I$ に
対して $R = R' \oplus I$ であるとする。$x_1, \ldots, x_r \in I$ を
$I/\mathfrak m_R I$ の基底に写るように取る。
$S = R'[[X_1, \ldots, X_r]]$ と置き，$X_i$ を $x_i$ に写す
$R'$-代数写像 $S \to R$ を考える。すべての $n \gg 0$ に対して写像
$S/\mathfrak m_S^n \to R/\mathfrak m_R^n$ が
$\mathcal{C}_\Lambda$ における左逆射をもつと仮定する。このとき
$S \to R$ は同型である。
\end{lemma}

\begin{proof}
$R = R' \oplus I$ なので
$$
\mathfrak m_R/\mathfrak m_R^2 =
\mathfrak m_{R'}/\mathfrak m_{R'}^2 \oplus I/\mathfrak m_RI
$$
であり，同様に
$$
\mathfrak m_S/\mathfrak m_S^2 =
\mathfrak m_{R'}/\mathfrak m_{R'}^2 \oplus \bigoplus kX_i
$$
である。したがって $n > 1$ に対して写像
$S/\mathfrak m_S^n \to R/\mathfrak m_R^n$ は余接空間上の同型を誘導する。
ゆえに左逆射
$h_n : R/\mathfrak m_R^n \to S/\mathfrak m_S^n$ は補題
\ref{formal-defos-lemma-surjective-cotangent-space} により全射である。
$h_n$ は左逆射として単射でもあるので，同型である。したがって標準的全射
$S/\mathfrak m_S^n \to R/\mathfrak m_R^n$ はすべて同型であり，結論を得る。
\end{proof}



\section{群亜群をファイバーとする余ファイバー圏}
\label{formal-defos-section-preliminary}

\noindent
理論を展開するにあたり，群亜群をファイバーとする{\it 余ファイバー圏}を扱う。
群亜群をファイバーとする{\it ファイバー圏}の定義と基本的性質は既知と仮定する。
Categories, Section \ref{categories-section-fibred-groupoids} を参照せよ。

\begin{definition}
\label{formal-defos-definition-category-cofibred-groupoids}
$\mathcal{C}$ を圏とする。$\mathcal{C}$ 上で{\it 群亜群をファイバーとする
余ファイバー圏}とは，関手 $p: \mathcal{F} \to \mathcal{C}$ を備えた圏
$\mathcal{F}$ であって，$p^{opp}: \mathcal{F}^{opp} \to
\mathcal{C}^{opp}$ により $\mathcal{F}^{opp}$ が $\mathcal{C}^{opp}$ 上で
群亜群をファイバーとするファイバー圏となるものをいう。
\end{definition}

\noindent
明示的には，$p: \mathcal{F} \to \mathcal{C}$ が群亜群をファイバーとする
余ファイバー圏であるとは，次の二条件が成り立つことをいう：
\begin{enumerate}
\item $\mathcal{C}$ の任意の射 $f: U \to V$ と，$U$ 上にある任意の対象
$x$ に対し，$f$ 上にある $\mathcal{F}$ の射 $x \to y$ が存在する。
\item $\mathcal{F}$ の任意の二射 $a: x \to y$ と $b: x \to z$，および
$p(b) = f \circ p(a)$ を満たす任意の射 $f: p(y) \to p(z)$ に対し，
$f$ 上にあり $b = c \circ a$ を満たす $\mathcal F$ の射
$c: y \to z$ が一意に存在する。
\end{enumerate}

\begin{remarks}
\label{formal-defos-remarks-cofibered-groupoids}
群亜群をファイバーとするファイバー圏に関するすべての事柄は，余ファイバーの
場合へ直接移る。以下の注意は記法を固定するためのものである。
$\mathcal{C}$ を圏とする。
\begin{enumerate}
\item 関手 $p: \mathcal{F} \to \mathcal{C}$ を記法から省くことが多い。
\item $\mathcal{C}$ の対象 $U$ 上のファイバー圏を $\mathcal{F}(U)$ と書く。
その対象は $U$ 上にある $\mathcal{F}$ の対象であり，その射は
$\text{id}_U$ 上にある $\mathcal{F}$ の射である。
$x, y$ が $\mathcal{F}(U)$ の対象であるとき，
$\Mor_{\mathcal{F}(U)}(x, y)$ を $\Mor_U(x, y)$ と書くことがある。
\item ファイバー圏 $\mathcal{F}(U)$ は群亜群である。Categories, Lemma
\ref{categories-lemma-fibred-groupoids} を参照せよ。
したがって $\mathcal{F}(U)$ の射はすべて同型である。
$\Mor_{\mathcal{F}(U)}(x, x)$ を $\text{Aut}_U(x)$ と書くことがある。
\item
\label{formal-defos-item-pushforward}
$\mathcal{F}$ を $\mathcal{C}$ 上で群亜群をファイバーとする余ファイバー圏，
$f: U \to V$ を $\mathcal{C}$ の射，$x \in \Ob(\mathcal{F}(U))$ とする。
$f$ に沿う $x$ の{\it 順像}とは，$f$ 上にある $\mathcal{F}$ の射
$x \to y$ をいう。順像は一意な同型を除いて一意である（Categories,
Definition \ref{categories-definition-cartesian-over-C} に続く議論を参照せよ）。
$f$ に沿う $x$ の「その」順像を $x \to f_*x$ と書くことがある。
\item $\mathcal{F}$ に対する{\it 順像の選択}とは，上の各組 $(x, f)$ に対して
$f$ に沿う $x$ の順像を選ぶことをいう。選択公理により，$\mathcal{F}$ に
対するこのような順像の選択を行える。
\item $\mathcal{F}$ を $\mathcal{C}$ 上で群亜群をファイバーとする余ファイバー圏
とする。$\mathcal{F}$ に対する順像の選択を与えると，それに付随する擬関手
$\mathcal{C} \to \textit{Groupoids}$ がある。この構成を用いることはないので，
詳述しない。
\item
\label{formal-defos-item-cofibered-morphism}
$\mathcal{C}$ 上で群亜群をファイバーとする余ファイバー圏の射とは，
$\mathcal{C}$ への射影と可換する関手をいう。$\mathcal{F}$ と
$\mathcal{F}'$ が $\mathcal{C}$ 上で群亜群をファイバーとする余ファイバー圏
であるとき，$\mathcal{F}$ から $\mathcal{F}'$ への射を
$\Mor_\mathcal{C}(\mathcal{F}, \mathcal{F}')$ と書く。
\item
\label{formal-defos-item-definition-cofibered-groupoids-2-category}
群亜群をファイバーとする余ファイバー圏は $(2, 1)$-圏
$\text{Cof}(\mathcal{C})$ をなす。その 1-射は
(\ref{formal-defos-item-cofibered-morphism}) で述べた射である。
% P10-E0428: preserve the frozen bare C codomain; correction is sidecar-only.
$p : \mathcal{F} \to C$ と $p': \mathcal{F}' \to \mathcal{C}$ が群亜群を
ファイバーとする余ファイバー圏であり，$\varphi, \psi : \mathcal{F} \to
\mathcal{F}'$ が $1$-射であるとする。このとき 2-射
$t : \varphi \to \psi$ とは，すべての $x \in \Ob(\mathcal{F})$ に対して
$p'(t_x) = \text{id}_{p(x)}$ を満たす関手の射である。
\item
\label{formal-defos-item-construction-associated-cofibered-groupoid}
$F : \mathcal{C} \to \textit{Groupoids}$ を関手とする。$F$ に付随する，
群亜群をファイバーとする余ファイバー圏 $\mathcal{F} \to \mathcal{C}$ を
次のように得る。$\mathcal{F}$ の対象は，$U \in \Ob(\mathcal{C})$ かつ
$x \in \Ob(F(U))$ である組 $(U, x)$ である。射
$(U, x) \to (V, y)$ は，$f \in \Mor_\mathcal{C}(U, V)$ かつ
$a \in \Mor_{F(V)}(F(f)(x), y)$ である組 $(f, a)$ である。
関手 $\mathcal{F} \to \mathcal{C}$ は $(U, x)$ を $U$ に写す。
Categories, Section \ref{categories-section-presheaves-groupoids} を参照せよ。
\item
\label{formal-defos-item-associated-functor-isomorphism-classes}
$\mathcal{F}$ を $\mathcal{C}$ 上で群亜群をファイバーとする余ファイバー圏
とする。$U \in \Ob(\mathcal{C})$ に対し，
$\overline{\mathcal{F}}(U)$ を圏 $\mathcal{F}(U)$ の同型類の集合とする。
% P10-E0429: the Japanese wording masks the frozen English number error.
$f : U \to V$ が $\mathcal{C}$ の射であるとき，$x$ の同型類を
$x$ の順像 $f_*x$ の同型類へ写すことにより，集合の写像
$\overline{\mathcal{F}}(U) \to \overline{\mathcal{F}}(V)$ を得る。
% P10-E0430: Japanese supplies normal citation punctuation; the English defect remains sidecar-only.
(\ref{formal-defos-item-pushforward}) を参照せよ。このとき
$\overline{\mathcal{F}} : \mathcal{C} \to \textit{Sets}$ は関手である。
同様に，$\varphi : \mathcal{F} \to \mathcal{G}$ が余ファイバー圏の射で
あるとき，付随する関手の射を
$\overline{\varphi}: \overline{\mathcal{F}} \to  \overline{\mathcal{G}}$
と書く。
\item
\label{formal-defos-item-convention-cofibered-sets}
$F: \mathcal{C} \to \textit{Sets}$ を関手とする。集合を離散圏，すなわち
恒等射しかもたない群亜群とみなせる。このとき構成
(\ref{formal-defos-item-construction-associated-cofibered-groupoid}) は $F$ に，集合を
ファイバーとする余ファイバー圏を対応させる。これは関手
$\mathcal{C} \to \textit{Sets}$ の圏から，$\mathcal{C}$ 上で群亜群を
ファイバーとする余ファイバー圏の圏への充満忠実埋め込みを定める。
関手の圏をこの埋め込みによる像と同一視する。したがって
$F : \mathcal{C} \to \textit{Sets}$ が関手であるとき，付随する集合を
ファイバーとする余ファイバー圏も $F$ と書く。また
$\varphi : F \to G$ が関手の射であるとき，対応する，集合をファイバーとする
余ファイバー圏の射も引き続き $\varphi$ と書き，逆も同様とする。
Categories, Section \ref{categories-section-fibred-in-sets} を参照せよ。
\item
\label{formal-defos-item-definition-yoneda}
$U$ を $\mathcal{C}$ の対象とする。関手
$\Mor_\mathcal{C}(U, -): \mathcal{C} \to \textit{Sets}$ を
$\underline{U}$ と書く。これは $\mathcal C^{opp}$ から関手の圏
$\mathcal{C} \to \textit{Sets}$ への充満忠実埋め込みを定める。
したがって $f : U \to V$ が射であるとき，誘導される射
$\underline{V} \to \underline{U}$ も引き続き $f$ と書くことが正当化され，
逆も同様である。
\item
\label{formal-defos-item-fibre-product}
群亜群をファイバーとする余ファイバー圏のファイバー積：
$\mathcal{F} \to \mathcal{H}$ と $\mathcal{G} \to \mathcal{H}$ が
$\mathcal{C}_\Lambda$ 上で群亜群をファイバーとする余ファイバー圏の射で
あるとき，その 2-ファイバー積は，Categories, Lemma
\ref{categories-lemma-2-product-categories-over-C} で述べられた，
$\mathcal{C}_\Lambda$ 上の圏としての 2-ファイバー積の構成によって与えられる。
\item
\label{formal-defos-item-product}
群亜群をファイバーとする余ファイバー圏の積：$\mathcal{F}$ と
$\mathcal{G}$ が $\mathcal{C}_\Lambda$ 上で群亜群をファイバーとする
余ファイバー圏であるとき，その積を，Categories, Lemma
\ref{categories-lemma-2-product-categories-over-C} で述べられた
$2$-ファイバー積 $\mathcal{F} \times_{\mathcal{C}_\Lambda} \mathcal{G}$
として定義する。
\item
\label{formal-defos-item-definition-restricting-base-category}
基底圏の制限：$p : \mathcal{F} \to \mathcal{C}$ を群亜群をファイバーとする
余ファイバー圏とし，$\mathcal{C}'$ を $\mathcal{C}$ の充満部分圏とする。
制限 $\mathcal{F}|_{\mathcal{C}'}$ は，$\mathcal{C}'$ の対象上にある対象から
なる $\mathcal{F}$ の充満部分圏である。これは関手
$p|_{\mathcal{C}'}: \mathcal{F}|_{\mathcal{C}'} \to \mathcal{C}'$
により，群亜群をファイバーとする余ファイバー圏となる。
\end{enumerate}
\end{remarks}








\section{pro-表現可能関手と前変形圏}
\label{formal-defos-section-cofibered-groupoids}

\noindent
基本的な目標は，$\mathcal{C}_\Lambda$ および
$\widehat{\mathcal{C}}_\Lambda$ 上で群亜群をファイバーとする余ファイバー圏を
理解することである。$\mathcal{C}_\Lambda$ は
$\widehat{\mathcal{C}}_\Lambda$ の充満部分圏なので，
$\widehat{\mathcal{C}}_\Lambda$ 上で群亜群をファイバーとする余ファイバー圏を
$\mathcal{C}_\Lambda$ へ制限できる。Remarks
\ref{formal-defos-remarks-cofibered-groupoids}
(\ref{formal-defos-item-definition-restricting-base-category}) を参照せよ。
とくに関手，なかでも表現可能関手についてこれを行える。
この方法で得られる $\mathcal{C}_\Lambda$ 上の関手を
pro-表現可能関手と呼ぶ。

\begin{definition}
\label{formal-defos-definition-prorepresentable}
$F : \mathcal{C}_\Lambda \to \textit{Sets}$ を関手とする。ある
$R \in \Ob(\widehat{\mathcal{C}}_\Lambda)$ に対して関手の同型
$F \cong \underline{R}|_{\mathcal{C}_\Lambda}$ が存在するとき，
$F$ は{\it pro-表現可能}であるという。
\end{definition}

\noindent
$F : \mathcal{C}_\Lambda \to \textit{Sets}$ が
$R \in \Ob(\widehat{\mathcal{C}}_\Lambda)$ により pro-表現可能ならば
$$
F(k) = \Mor_{\widehat{\mathcal{C}}_\Lambda}(R, k) = \{*\}
$$
は一点集合であることに注意せよ。
% P10-E0431: Japanese renders the intended clause; the frozen English grammar defect remains sidecar-only.
変形理論に現れる，$\mathcal{C}_\Lambda$ 上で群亜群をファイバーとする
余ファイバー圏は，同様の条件を満たすことが多い。

\begin{definition}
\label{formal-defos-definition-predeformation-category}
{\it 前変形圏} $\mathcal{F}$ とは，$\mathcal{F}(k)$ が一つの対象と一つの射を
もつ圏と同値であるような，$\mathcal{C}_\Lambda$ 上で群亜群をファイバーとする
余ファイバー圏をいう。すなわち，$\mathcal{F}(k)$ は少なくとも一つの対象を
含み，任意の二対象の間に一意な射が存在する。
{\it 前変形圏の射}とは，$\mathcal{C}_\Lambda$ 上で群亜群をファイバーとする
余ファイバー圏の射をいう。
\end{definition}

\noindent
前変形圏は次の性質をもつ。$x_0 \in \Ob(\mathcal{F}(k))$ とする。
このとき $\mathcal{F}$ の各対象には $x_0$ への一意な射が備わる。
実際，$x$ が $A$ 上の $\mathcal{F}$ の対象ならば，商写像
$q : A \to k$ に沿う順像 $x \to q_*x$ を選べる。
一意な同型 $q_*x \to x_0$ が存在し，合成
$x \to q_*x \to x_0$ が求める射である。

\begin{remark}
\label{formal-defos-remark-predeformation-functor}
関手 $F: \mathcal{C}_\Lambda \to \textit{Sets}$ に付随する，集合をファイバーとする
余ファイバー圏が前変形圏であるとき，すなわち，\ $F(k)$ が一点集合であるとき，
この関手を{\it 前変形関手}という。したがって $\mathcal{F}$ が前変形圏ならば，
$\overline{\mathcal{F}}$ は前変形関手である。
\end{remark}

\begin{remark}
\label{formal-defos-remark-localize-cofibered-groupoid}
$p : \mathcal{F} \to \mathcal{C}_\Lambda$ を群亜群をファイバーとする
余ファイバー圏とし，$x \in \Ob(\mathcal{F}(k))$ とする。
$x$ 上の対象の圏を $\mathcal{F}_x$ と書く。
$\mathcal{F}_x$ の対象は射 $y \to x$ である。
$\mathcal{F}_x$ における射 $(y \to x) \to (z \to x)$ は可換図式
$$
\xymatrix{
y \ar[rr] \ar[dr] & & z \ar[dl] \\
& x &
}
$$
である。忘却関手 $\mathcal{F}_x \to \mathcal{F}$ がある。
関手 $p_x : \mathcal{F}_x \to \mathcal{C}_\Lambda$ を合成
$\mathcal{F}_x \to \mathcal{F} \xrightarrow{p} \mathcal{C}_\Lambda$
として定義する。このとき $p_x : \mathcal{F}_x \to
\mathcal{C}_\Lambda$ は前変形圏である（証明は省略する）。
このようにして，$\mathcal{C}_\Lambda$ 上で群亜群をファイバーとする任意の
余ファイバー圏から，各 $x \in \Ob(\mathcal{F}(k))$ における前変形圏へ
移ることができる。
\end{remark}


\section{形式対象と完備化圏}
\label{formal-defos-section-formal-objects}

\noindent
本節では，$\mathcal{C}_\Lambda$ 上で群亜群をファイバーとする余ファイバー圏と，
$\widehat{\mathcal{C}}_\Lambda$ 上で群亜群をファイバーとする余ファイバー圏との
間を双方向に移る方法を論じる。

\begin{definition}
\label{formal-defos-definition-formal-objects}
$\mathcal{F}$ を $\mathcal{C}_\Lambda$ 上で群亜群をファイバーとする
余ファイバー圏とする。$\mathcal{F}$ の{\it 形式対象の圏
$\widehat{\mathcal{F}}$}とは，次の対象と射をもつ圏である。
\begin{enumerate}
\item $\mathcal{F}$ の{\it 形式対象 $\xi = (R, \xi_n, f_n)$}は，
$\widehat{\mathcal{C}}_\Lambda$ の対象 $R$，ならびに
$n \in \mathbf{N}$ で添字付けられた
$\mathcal{F}(R/\mathfrak m_R^n)$ の対象 $\xi_n$ と，射影
$R/\mathfrak m_R^{n + 1} \to R/\mathfrak m_R^n$ 上にある射
$f_n : \xi_{n + 1} \to \xi_n$ の族からなる。
\item $\xi = (R, \xi_n, f_n)$ と $\eta = (S, \eta_n, g_n)$ を
$\mathcal{F}$ の形式対象とする。{\it 形式対象の射
$a : \xi \to \eta$}は，$\widehat{\mathcal{C}}_\Lambda$ における写像
$a_0 : R \to S$ と，$R/\mathfrak m_R^n \to S/\mathfrak m_S^n$
上にある $\mathcal{F}$ の射の族 $a_n : \xi_n \to \eta_n$ であって，
各 $n$ に対して図式
$$
\xymatrix{
\xi_{n + 1} \ar[r]_{f_n} \ar[d]_{a_{n + 1}} & \xi_n \ar[d]^{a_n} \\
\eta_{n + 1} \ar[r]^{g_n} & \eta_n
}
$$
が可換となるものからなる。
\end{enumerate}
\end{definition}

\noindent
形式対象の圏には，対象 $(R, \xi_n, f_n)$ を $R$ に写し，射
$(R, \xi_n, f_n) \to (S, \eta_n, g_n)$ を写像 $R \to S$ に写す関手
$\widehat{p}: \widehat{\mathcal{F}} \to
\widehat{\mathcal{C}}_\Lambda$ が備わる。

\begin{lemma}
\label{formal-defos-lemma-completion-cofibred}
$p : \mathcal{F} \to \mathcal{C}_\Lambda$ を群亜群をファイバーとする
余ファイバー圏とする。このとき
$\widehat{p} : \widehat{\mathcal{F}} \to
\widehat{\mathcal{C}}_\Lambda$ は群亜群をファイバーとする
余ファイバー圏である。
\end{lemma}

\begin{proof}
$R \to S$ を $\widehat{\mathcal{C}}_\Lambda$ における環写像とし，
$(R, \xi_n, f_n)$ を $\widehat{\mathcal{F}}$ の対象とする。
各 $n$ に対し，$R/\mathfrak m_R^n \to S/\mathfrak m_S^n$ に沿う
$\xi_n$ の順像 $\xi_n \to \eta_n$ を選ぶ。各 $n$ に対し，
$S/\mathfrak m_S^{n + 1} \to S/\mathfrak m_S^n$ 上にある
$\mathcal{F}$ の射 $g_n : \eta_{n + 1} \to \eta_n$ であって，図式
$$
\xymatrix{
\xi_{n + 1} \ar[d] \ar[r]_{f_n} & \xi_n \ar[d] \\
\eta_{n + 1} \ar[r]^{g_n} & \eta_n
}
$$
を可換にするものが一意に存在する。
% P10-E0432: preserve the frozen attribution to the first axiom; correction is sidecar-only.
（群亜群をファイバーとする余ファイバー圏の第一公理による。）
したがって $R \to S$ 上にある射
$(R, \xi_n, f_n) \to (S, \eta_n, g_n)$ を得る。すなわち，
群亜群をファイバーとする余ファイバー圏の第一公理が
$\widehat{\mathcal{F}}$ に対して成り立つ。第二公理を見るため，射
$a : (R, \xi_n, f_n) \to (S, \eta_n, g_n)$ と
$b : (R, \xi_n, f_n) \to (T, \theta_n, h_n)$ が
$\widehat{\mathcal{F}}$ にあり，さらに
$c_0 \circ a_0 = b_0$ を満たす射 $c_0 : S \to T$ が
$\widehat{\mathcal{C}}_\Lambda$ にあると仮定する。
$\mathcal{F}$ に対する，群亜群をファイバーとする余ファイバー圏の第二公理により，
$S/\mathfrak m_S^n \to T/\mathfrak m_T^n$ 上にあり
$c_n \circ a_n = b_n$ を満たす写像 $c_n : \eta_n \to \theta_n$ が一意に得られる。
$c = (c_n)_{n \geq 0}$ と置けば，求める
$\widehat{\mathcal{F}}$ の射
$c : (S, \eta_n, g_n) \to (T, \theta_n, h_n)$ を得る
（$h_n \circ c_{n + 1} = c_n \circ g_n$ の確認は省略する）。
\end{proof}

\begin{definition}
\label{formal-defos-definition-completion}
$p : \mathcal{F} \to \mathcal{C}_\Lambda$ を群亜群をファイバーとする
余ファイバー圏とする。群亜群をファイバーとする余ファイバー圏
$\widehat{p} : \widehat{\mathcal  F} \to
\widehat{\mathcal{C}}_\Lambda$ を $\mathcal{F}$ の{\it 完備化}という。
\end{definition}

\noindent
$\mathcal{F}$ が $\mathcal C_\Lambda$ 上で群亜群をファイバーとする
余ファイバー圏であるとき，$R \in
\Ob(\widehat{\mathcal{C}}_\Lambda)$ に対する
$\widehat{\mathcal{F}}(R)$ を，$R$ の極大イデアルの冪による
フィルトレーションを用いて定義した。しかし
$\mathcal{I} = (I_n)$ が $\mathfrak{m}_R$-進位相を誘導する，
$R$ のイデアルによるフィルトレーションであるとする。
$\widehat{\mathcal{F}}_\mathcal{I}(R)$ を次の対象と射をもつ圏として定義する：
\begin{enumerate}
\item 対象は，$\mathcal{F}(R/I_n)$ の対象 $\xi_n$ と，射影
$R/I_{n + 1} \to R/I_n$ 上にある射
$f_n : \xi_{n + 1} \to \xi_n$ の族
$(\xi_n, f_n)_{n \in \mathbf{N}}$ である。
\item 射 $a : (\xi_n, f_n) \to (\eta_n, g_n)$ は，
$\mathcal{F}(R/I_n)$ における射の族
$a_n : \xi_n \to \eta_n$ であって，各 $n$ に対して図式
$$
\xymatrix{
\xi_{n + 1} \ar[r]^{f_n} \ar[d]_{a_{n + 1}} & \xi_n \ar[d]^{a_n} \\
\eta_{n + 1} \ar[r]^{g_n} & \eta_n
}
$$
が可換となるものからなる。
\end{enumerate}

\begin{lemma}
\label{formal-defos-lemma-formal-objects-different-filtration}
上の状況において，$\widehat{\mathcal{F}}_\mathcal{I}(R)$ は圏
$\widehat{\mathcal{F}}(R)$ と同値である。
\end{lemma}

\begin{proof}
同値
$\widehat{\mathcal{F}}_\mathcal{I}(R) \to
\widehat{\mathcal{F}}(R)$ を次のように定義できる。
各 $n$ に対し，$m(n)$ を $I_m \subset \mathfrak m_R^n$ となる最小の
$m$ とする。$\widehat{\mathcal{F}}_\mathcal{I}(R)$ の対象
$(\xi_n, f_n)$ が与えられたとき，$\eta_n$ を
$R/I_{m(n)} \to R/\mathfrak m_R^n$ に沿う $\xi_{m(n)}$ の順像とする。
$g_n : \eta_{n + 1} \to \eta_n$ を，
$R/\mathfrak m_R^{n + 1} \to R/\mathfrak m_R^n$ 上にあり，図式
$$
\xymatrix{
\xi_{m(n + 1)} \ar[rrr]_{f_{m(n)} \circ \ldots \circ f_{m(n + 1) - 1}} \ar[d]
& & & \xi_{m(n)} \ar[d] \\
\eta_{n + 1} \ar[rrr]^{g_n} & & & \eta_n
}
$$
を可換にする一意な $\mathcal{F}$ の射とする
（存在と一意性は余ファイバー圏の公理により保証される）。
関手 $\widehat{\mathcal{F}}_\mathcal{I}(R) \to
\widehat{\mathcal{F}}(R)$ は $(\xi_n, f_n)$ を
$(R, \eta_n, g_n)$ に写す。これが実際に圏の同値であることの確認は省略する。
\end{proof}

\begin{remark}
\label{formal-defos-remark-different-sequence-ideals}
$p : \mathcal{F} \to \mathcal{C}_\Lambda$ を群亜群をファイバーとする
余ファイバー圏とする。各 $R \in
\Ob(\widehat{\mathcal{C}}_\Lambda)$ に対し，$R$ のイデアルによる
フィルトレーション $\mathcal{I}_R$ が与えられていると仮定する。
すべての $R$ に対して $\mathcal{I}_R$ が $R$ 上の
$\mathfrak m_R$-進位相を誘導するならば，
$\widehat{\mathcal{F}}$ の定義を模倣して圏
$\widehat{\mathcal{F}}_\mathcal{I}$ を定義できる。この圏には，
群亜群をファイバーとする余ファイバー圏にする射
$\widehat{p}_\mathcal{I} : \widehat{\mathcal{F}}_\mathcal{I} \to
\widehat{\mathcal{C}}_\Lambda$ が備わり，上で定義した
$\widehat{\mathcal{F}}_{\mathcal{I}_R}(R)$ と
$\widehat{\mathcal{F}}_\mathcal{I}(R)$ は同型である。
群亜群をファイバーとする余ファイバー圏
$\widehat{\mathcal{F}}_\mathcal{I}$ と $\widehat{\mathcal{F}}$ は，
各対象 $R \in \Ob(\widehat{\mathcal{C}}_\Lambda)$ 上で補題
\ref{formal-defos-lemma-formal-objects-different-filtration} の同値を用いることにより同値である。
\end{remark}

\begin{remark}
\label{formal-defos-remark-completion-functor}
$F: \mathcal{C}_\Lambda \to \textit{Sets}$ を関手とする。
関手を集合をファイバーとする余ファイバー圏と同一視すると，$F$ の完備化は関手
$\widehat{F} : \widehat{\mathcal{C}}_\Lambda \to \textit{Sets}$
であり，$\widehat{F}(S) = \lim F(S/\mathfrak{m}_S^{n})$ で与えられる。
これは Schlessinger の論文 \cite{Sch} における定義と一致する。
\end{remark}

\begin{remark}
\label{formal-defos-remark-restrict-completion}
$\mathcal{F}$ を $\mathcal{C}_\Lambda$ 上で群亜群をファイバーとする
余ファイバー圏とする。標準的同値
$$
can :
\widehat{\mathcal{F}}|_{\mathcal{C}_\Lambda}
\longrightarrow
\mathcal{F}.
$$
が存在すると主張する。実際，$A \in \Ob(\mathcal{C}_\Lambda)$ とし，
$(A, \xi_n, f_n)$ を
$\widehat{\mathcal{F}}|_{\mathcal{C}_\Lambda}(A)$ の対象とする。
$A$ は Artin なので，$\mathfrak m_A^m = 0$ を満たす最小の
$m \in \mathbf{N}$ が存在する。このとき $can$ は
$(A, \xi_n, f_n)$ を $\xi_m$ に写す。この関手は
Categories, Lemma \ref{categories-lemma-equivalence-fibred-categories}
により，群亜群をファイバーとする余ファイバー圏の同値である。
実際，補題 \ref{formal-defos-lemma-formal-objects-different-filtration} と，局所 Artin 環
$A$ 上の $\mathfrak m_A$-進位相が零イデアルから生じるという事実により，
すべてのファイバー圏上で同値となる。関手 $can$ の準逆を通じて，
$\mathcal{F}$ を $\widehat{\mathcal{F}}$ の充満部分圏と同一視することが多い。
\end{remark}

\begin{remark}
\label{formal-defos-remark-completion-morphism}
$\varphi : \mathcal{F} \to \mathcal{G}$ を，
$\mathcal{C}_\Lambda$ 上で群亜群をファイバーとする余ファイバー圏の射とする。
このとき $\widehat{\mathcal{C}}_\Lambda$ 上で群亜群をファイバーとする
余ファイバー圏の誘導射
$\widehat{\varphi}: \widehat{\mathcal{F}} \to
\widehat{\mathcal{G}}$ が存在する。これは
$\widehat{\mathcal{F}}$ の対象 $\xi = (R, \xi_n, f_n)$ を
$(R, \varphi(\xi_n), \varphi(f_n))$ に写し，
$\widehat{\mathcal{F}}$ の対象 $\xi$ と $\eta$ の間の射
$(a_0 : R \to S, a_n : \xi_n \to \eta_n)$ を
$(a_0 : R \to S, \varphi(a_n) : \varphi(\xi_n) \to
\varphi(\eta_n))$ に写す。最後に，$t : \varphi \to \varphi'$ が，
群亜群をファイバーとする余ファイバー圏の $1$-射
$\varphi, \varphi': \mathcal{F} \to \mathcal{G}$ の間の $2$-射ならば，
$2$-射 $\widehat{t} : \widehat{\varphi} \to \widehat{\varphi}'$ を得る。
具体的に，上の $\xi = (R, \xi_n, f_n)$ に対して
% P10-E0433: preserve the frozen ill-typed component at this locus.
$\widehat{t}_\xi = (t_{\varphi(\xi_n)})$ と置く。
したがって完備化は $2$-圏の間の関手
$$
\widehat{~} :
\text{Cof}(\mathcal{C}_\Lambda)
\longrightarrow
\text{Cof}(\widehat{\mathcal{C}}_\Lambda)
$$
を定める。すなわち，$\mathcal{C}_\Lambda$ 上で群亜群をファイバーとする
余ファイバー圏の $2$-圏から，$\widehat{\mathcal{C}}_\Lambda$ 上で
群亜群をファイバーとする余ファイバー圏の $2$-圏への関手である。
\end{remark}

\begin{remark}
\label{formal-defos-remark-completion-restriction-adjoint}
Remark \ref{formal-defos-remark-completion-morphism} の完備化関手と，Remarks
\ref{formal-defos-remarks-cofibered-groupoids}
(\ref{formal-defos-item-definition-restricting-base-category}) の制限関手
$|_{\mathcal{C}_\Lambda} : \text{Cof}(\widehat{\mathcal{C}}_\Lambda)
\to \text{Cof}(\mathcal{C}_\Lambda)$ は，次の厳密な意味で
「2-随伴」であると主張する。
$\mathcal{F} \in \Ob(\text{Cof}(\mathcal{C}_\Lambda))$ および
$\mathcal{G} \in \Ob(\text{Cof}(\widehat{\mathcal{C}}_\Lambda))$
とする。このとき圏の同値
$$
\Phi :
\Mor_{\mathcal{C}_\Lambda}(
\mathcal{G}|_{\mathcal{C}_\Lambda}, \mathcal{F})
\longrightarrow
\Mor_{\widehat{\mathcal{C}}_\Lambda}(\mathcal{G}, \widehat{\mathcal{F}})
$$
が存在する。この同値を記述するため，標準射
$\mathcal{G} \to \widehat{\mathcal{G}|_{\mathcal{C}_\Lambda}}$ および
$\widehat{\mathcal{F}}|_{\mathcal{C}_\Lambda} \to \mathcal{F}$ を
次のように定義する。
\begin{enumerate}
\item
% P10-E0434: preserve the frozen extra closing parenthesis.
$R \in \Ob(\widehat{\mathcal{C}}_\Lambda))$ とし，$\xi$ をファイバー圏
$\mathcal{G}(R)$ の対象とする。各 $n \in \mathbf{N}$ に対し，
$R/\mathfrak m_R^n$ への $\xi$ の順像 $\xi \to \xi_n$ を選び，
$f_n : \xi_{n + 1} \to \xi_n$ を誘導される射とする。
このとき $\mathcal{G} \to
\widehat{\mathcal{G}|_{\mathcal{C}_\Lambda}}$ は
$\xi$ を $(R, \xi_n, f_n)$ に写す。
\item これは Remark \ref{formal-defos-remark-restrict-completion} の同値
$can : \widehat{\mathcal{F}}|_{\mathcal{C}_\Lambda} \to
\mathcal{F}$ である。
\end{enumerate}
以上により，同値
$\Phi : \Mor_{\mathcal{C}_\Lambda}(
\mathcal{G}|_{\mathcal{C}_\Lambda}, \mathcal{F})  \to
\Mor_{\widehat{\mathcal{C}}_\Lambda}(\mathcal{G},
\widehat{\mathcal{F}})$
は射 $\varphi : \mathcal{G}|_{\mathcal{C}_\Lambda} \to
\mathcal{F}$ を
$$
\mathcal{G} \to \widehat{\mathcal{G}|_{\mathcal{C}_\Lambda}}
\xrightarrow{\widehat{\varphi}} \widehat{\mathcal{F}}
$$
に写す。$\Phi$ の準逆
$\Psi :
\Mor_{\widehat{\mathcal{C}}_\Lambda}(
\mathcal{G}, \widehat{\mathcal{F}}) \to
\Mor_{\mathcal{C}_\Lambda}(
\mathcal{G}|_{\mathcal{C}_\Lambda}, \mathcal{F})$
があり，$\psi : \mathcal{G} \to \widehat{\mathcal{F}}$ を
$$
\mathcal{G}|_{\mathcal{C}_\Lambda} \xrightarrow{\psi|_{\mathcal{C}_\Lambda}}
\widehat{\mathcal{F}}|_{\mathcal{C}_\Lambda} \to \mathcal{F}.
$$
に写す。$\Phi$ と $\Psi$ が準逆であることの確認は省略する。
$\Phi$ の関手性も扱わない（それは，どうしても避けたい 3-圏の領域へ
踏み込むことになるからである）。
\end{remark}

\begin{remark}
\label{formal-defos-remark-completion-restriction-cofset-adjoint}
圏 $\mathcal{C}$ に対し，$\mathcal{C}$ 上の余ファイバー集合の圏を
$\text{CofSet}(\mathcal{C})$ と書く。これは関手の圏
$\mathcal{C} \to \textit{Sets}$ と同型な $1$-圏である。
% P10-E0435: Japanese supplies the preposition missing from the frozen English sentence.
Remarks \ref{formal-defos-remarks-cofibered-groupoids}
(\ref{formal-defos-item-convention-cofibered-sets}) を参照せよ。
完備化関手と制限関手は，同じ記号で表す関手
$\widehat{~} : \text{CofSet}(\mathcal{C}_\Lambda) \to
\text{CofSet}(\widehat{\mathcal{C}}_\Lambda)$ および
$|_{\mathcal{C}_\Lambda} : \text{CofSet}(\widehat{\mathcal{C}}_\Lambda) \to
\text{CofSet}(\mathcal{C}_\Lambda)$ に制限される。
余ファイバー集合の圏上の関手として，完備化と制限は通常の 1-圏的意味で随伴する。
すなわち Remark \ref{formal-defos-remark-completion-restriction-adjoint} と同じ構成が，
$F \in \Ob(\text{CofSet}(\mathcal{C}_\Lambda))$ および
$G \in \Ob(\text{CofSet}(\widehat{\mathcal{C}}_\Lambda))$ に対して
関手的全単射
$$
\Mor_{\mathcal{C}_\Lambda}(G|_{\mathcal{C}_\Lambda}, F)
\longrightarrow
\Mor_{\widehat{\mathcal{C}}_\Lambda}(G, \widehat{F})
$$
を定める。ここでも写像
$\widehat{F}|_{\mathcal{C}_\Lambda} \to F$ は同型である。
\end{remark}

\begin{remark}
\label{formal-defos-remark-restrict-complete-continuous-functor}
$G : \widehat{\mathcal{C}}_\Lambda \to \textit{Sets}$ を
極限と可換する関手とする。このとき Remark
\ref{formal-defos-remark-completion-restriction-adjoint} で述べた写像
$G \to \widehat{G|_{\mathcal{C}_\Lambda}}$ は同型である。
実際，$S$ が $\widehat{\mathcal{C}}_\Lambda$ の対象ならば，標準的全単射
$$
\widehat{G|_{\mathcal{C}_\Lambda}}(S) =
\lim_n G(S/\mathfrak{m}_S^n) =
G(\lim_n S/\mathfrak{m}_S^n) = G(S).
$$
をもつ。とくに $R$ が $\widehat{\mathcal{C}}_\Lambda$ の対象ならば，
表現可能関手 $\underline{R}$ は極限の定義により極限と可換するので，
$\underline{R} = \widehat{\underline{R}|_{\mathcal{C}_\Lambda}}$ である。
\end{remark}

\begin{remark}
\label{formal-defos-remark-formal-objects-yoneda}
$R$ を $\widehat{\mathcal{C}}_\Lambda$ の対象とする。これは Remarks
\ref{formal-defos-remarks-cofibered-groupoids} (\ref{formal-defos-item-definition-yoneda}) で述べた関手
$\underline{R}: \widehat{\mathcal{C}}_\Lambda \to \textit{Sets}$
を定める。通常どおり，この関手を付随する余ファイバー集合と同一視する。
$\mathcal{F}$ が $\mathcal{C}_\Lambda$ 上の余ファイバー圏ならば，
圏の同値
\begin{equation}
\label{formal-defos-equation-formal-objects-maps}
\Mor_{\mathcal{C}_\Lambda}(
\underline{R}|_{\mathcal{C}_\Lambda}, \mathcal{F})
\longrightarrow
\widehat{\mathcal{F}}(R).
\end{equation}
が存在する。これは合成
$$
\Mor_{\mathcal{C}_\Lambda}(
\underline{R}|_{\mathcal{C}_\Lambda}, \mathcal{F})
\xrightarrow{\Phi}
\Mor_{\widehat{\mathcal{C}}_\Lambda}(
\underline{R}, \widehat{\mathcal{F}})
\xrightarrow{\sim}
\widehat{\mathcal{F}}(R)
$$
で与えられる。ここで $\Phi$ は Remark
\ref{formal-defos-remark-completion-restriction-adjoint} のものであり，第二の同値は
2-Yoneda 補題（Categories, Lemma
\ref{categories-lemma-yoneda-2category} の余ファイバー版）から得られる。
明示的には，この同値は射
$\varphi : \underline{R}|_{\mathcal{C}_\Lambda} \to \mathcal{F}$ を
$\widehat{\mathcal{F}}(R)$ の形式対象
$(R, \varphi(R \to R/\mathfrak{m}_R^n), \varphi(f_n))$ に写す。
ここで $f_n : R/\mathfrak m_R^{n + 1} \to R/\mathfrak m_R^n$
は射影である。

\medskip\noindent
$\mathcal{F}$ に対して順像が選択されていると仮定する。
任意の $\xi \in \Ob(\widehat{\mathcal{F}}(R))$ が与えられたとき，
(\ref{formal-defos-equation-formal-objects-maps}) により $\xi$ に写る明示的な射
$\underline{\xi} : \underline{R}|_{\mathcal{C}_\Lambda} \to
\mathcal{F}$ を構成する。$\xi = (R, \xi_n, f_n)$ と書く。
$\underline{R}|_{\mathcal{C}_\Lambda}$ の対象 $\alpha$ は，
$A$ が Artin であるような $\widehat{\mathcal{C}}_\Lambda$ の射
$\alpha : R \to A$ と同じものである。
$\mathfrak m_A^m = 0$ を満たす最小の $m \in \mathbf{N}$ を取る。
このとき $\alpha$ は一意な
$\alpha_m : R/\mathfrak m_R^m \to A$ を経由するので，
$\underline{\xi}(\alpha) = \alpha_{m, *}\xi_m$ と置ける。
射上での $\underline{\xi}$ の記述，および
(\ref{formal-defos-equation-formal-objects-maps}) により
$\underline{\xi}$ が $\xi$ に写ることの証明は省略する。

\medskip\noindent
$\widehat{\mathcal{F}}$ に対して順像が選択されていると仮定する。
この場合，Categories, Lemma \ref{categories-lemma-yoneda-2category}
の証明は，2-Yoneda 同値の明示的な準逆
$$
\iota :
\widehat{\mathcal{F}}(R) \longrightarrow
\Mor_{\widehat{\mathcal{C}}_\Lambda}(
\underline{R}, \widehat{\mathcal{F}})
$$
を与える。これは $\xi$ を，
% P10-E0436: preserve the frozen uncompleted base-category subscript.
$f \in \underline{R}(S) = \Mor_{\mathcal{C}_\Lambda}(R, S)$
を $f_*\xi$ に写す射
$\iota(\xi) : \underline{R} \to \widehat{\mathcal{F}}$ に写す。
したがって (\ref{formal-defos-equation-formal-objects-maps}) の準逆は
$$
\widehat{\mathcal{F}}(R)
\xrightarrow{\iota}
\Mor_{\widehat{\mathcal{C}}_\Lambda}(
\underline{R}, \widehat{\mathcal{F}})
\xrightarrow{\Psi}
\Mor_{\mathcal{C}_\Lambda}(
\underline{R}|_{\mathcal{C}_\Lambda}, \mathcal{F})
$$
である。ここで $\Psi$ は Remark
\ref{formal-defos-remark-completion-restriction-adjoint} のものである。
$\xi \in \Ob(\widehat{\mathcal{F}}(R))$ が与えられると，
$\Psi(\iota(\xi)) \cong \underline{\xi}$ である。
ここで $\underline{\xi}$ は前段落のものである。実際，圏の同値
(\ref{formal-defos-equation-formal-objects-maps}) の下で両者とも $\xi$ に写る。
$\underline{R} = \widehat{\underline{R}|_{\mathcal{C}_\Lambda}}$
（Remark \ref{formal-defos-remark-restrict-complete-continuous-functor} を参照）を用い，
$\Phi$ と $\Psi$ の定義を展開すれば，$\iota(\xi)$ は
$\underline{\xi}$ の完備化と同型であると結論される。
\end{remark}

\begin{remark}
\label{formal-defos-remark-formal-objects-yoneda-map}
$\mathcal{F}$ を $\mathcal{C}_\Lambda$ 上で群亜群をファイバーとする
余ファイバー圏とする。$\xi = (R, \xi_n, f_n)$ と
$\eta = (S, \eta_n, g_n)$ を $\mathcal{F}$ の形式対象とする。
$a = (a_n) : \xi \to \eta$ を形式対象の射，すなわち
$\widehat{\mathcal{F}}$ の射とする。
$f = \widehat{p}(a) = a_0 : R \to S$ を
$\widehat{\mathcal{C}}_\Lambda$ における $a$ の射影とする。
このとき $2$-可換図式
$$
\xymatrix{
\underline{R}|_{\mathcal{C}_\Lambda} \ar[rd]_{\underline{\xi}} & &
\underline{S}|_{\mathcal{C}_\Lambda} \ar[ll]^f \ar[ld]^{\underline{\eta}} \\
& \mathcal{F}
}
$$
を得る。ここで $\underline{\xi}$ と $\underline{\eta}$ は Remark
\ref{formal-defos-remark-formal-objects-yoneda} で構成した射である。
これを見るため，$\alpha : S \to A$ を
$\underline{S}|_{\mathcal{C}_\Lambda}$ の対象とする
（loc.\ cit. を参照）。
$\mathfrak m_A^m = 0$ を満たす最小の $m \in \mathbf{N}$ を取る。
下の二射の合成が $\alpha$ となる可換図式
$$
\xymatrix{
R \ar[d]^f \ar[r] & R/\mathfrak m_R^m \ar[d]_{f_m} \ar[rd]^{\beta_m} \\
S \ar[r] & S/\mathfrak m_S^m \ar[r]^{\alpha_m} & A
}
$$
を得る。このとき
$\underline{\eta}(\alpha) = \alpha_{m, *}\eta_m$ および
$\underline{\xi}(\alpha \circ f) = \beta_{m, *}\xi_m$ である。
射 $a_m : \xi_m \to \eta_m$ は $f_m$ 上にあるので，
$\text{id}_A$ 上にある標準射
$$
\underline{\xi}(\alpha \circ f) = \beta_{m, *}\xi_m
\longrightarrow
\underline{\eta}(\alpha) = \alpha_{m, *}\eta_m
$$
であって，図式
$$
\xymatrix{
\xi_m \ar[r] \ar[d]^{a_m} &  \beta_{m, *}\xi_m \ar[d] \\
\eta_m \ar[r] &  \alpha_{m, *}\eta_m
}
$$
を可換にするものを，群亜群をファイバーとする余ファイバー圏の公理により得る。
これは，この remark の最初の図式の 2-可換性を与える関手の変換
$\underline{\xi} \circ f \to \underline{\eta}$ を定める。
\end{remark}

\begin{remark}
\label{formal-defos-remark-spell-out-formal-object}
Remark \ref{formal-defos-remark-formal-objects-yoneda} によれば，
$\mathcal{F}$ の形式対象 $\xi$ を与えることは，pro-表現可能関手
$U : \mathcal{C}_\Lambda \to \textit{Sets}$ と射
$U \to \mathcal{F}$ を与えることに同値である。
\end{remark}

\section{滑らかな射}
\label{formal-defos-section-smooth-morphisms}

\noindent
本節では，$\mathcal{C}_\Lambda$ 上で群亜群をファイバーとする
余ファイバー圏の滑らかな射を論じる。

\begin{definition}
\label{formal-defos-definition-smooth-morphism}
$\varphi : \mathcal{F} \to \mathcal{G}$ を，
$\mathcal{C}_\Lambda$ 上で群亜群をファイバーとする余ファイバー圏の射とする。
$\varphi$ が次の条件を満たすとき，これを{\it 滑らか}という：
$B \to A$ を $\mathcal{C}_\Lambda$ における全射環写像とする。
$y \in \Ob(\mathcal{G}(B)), x \in \Ob(\mathcal{F}(A))$ とし，
$y \to \varphi(x)$ を $B \to A$ 上にある射とする。このとき
$x' \in \Ob(\mathcal{F}(B))$ と，$B \to A$ 上にある射
$x' \to x$，および $\text{id}: B \to B$ 上にある射
$\varphi(x') \to y$ が存在して，図式
$$
\xymatrix{
\varphi(x') \ar[r] \ar[dr] & y \ar[d] \\
& \varphi(x)
}
$$
を可換にする。
\end{definition}

\begin{lemma}
\label{formal-defos-lemma-smoothness-small-extensions}
$\varphi : \mathcal{F} \to \mathcal{G}$ を，
$\mathcal{C}_\Lambda$ 上で群亜群をファイバーとする余ファイバー圏の射とする。
Definition \ref{formal-defos-definition-smooth-morphism} の条件が小拡大
$B \to A$ に対してのみ成り立つと仮定しても，$\varphi$ は滑らかである。
\end{lemma}

\begin{proof}
$B \to A$ を $\mathcal{C}_\Lambda$ における全射環写像とする。
$y \in \Ob(\mathcal{G}(B))$, $x \in \Ob(\mathcal{F}(A))$ とし，
$y \to \varphi(x)$ を $B \to A$ 上にある射とする。
Lemma \ref{formal-defos-lemma-factor-small-extension} により，$B \to A$ を小拡大
$B = B_n \to B_{n-1} \to \ldots \to B_0 = A$ の合成に分解できる。
$n$ に関する帰納法で論じる。$n = 1$ なら仮定により成り立つ。
$n > 1$ ならば，$f : B = B_n \to B_{n - 1}$ および
$g : B_{n - 1} \to B_0 = A$ と書く。$f$ に沿う $y$ の順像
$y \to f_* y$ を選べば，射 $y \to \varphi(x)$ は
$y \to f_* y \to \varphi(x)$ と分解する。帰納法の仮定により，
$g : B_{n - 1} \to A$ 上にある $x_{n - 1} \to x$ と，
$\text{id} : B_{n - 1} \to B_{n - 1}$ 上にある
$a : \varphi(x_{n - 1}) \to f_*y$ であって，図式
$$
\xymatrix{
\varphi(x_{n - 1}) \ar[r]_-a \ar[dr] & f_*y \ar[d] \\
& \varphi(x)
}
$$
を可換にするものを見いだせる。$y \to f_*y$ と
$a^{-1} : f_*y \to \varphi(x_{n - 1})$ の合成
$y \to \varphi(x_{n - 1})$ に仮定を適用できる。
これにより，$B_n \to B_{n - 1}$ 上にある
$x_n \to x_{n - 1}$ と，$\text{id} : B_n \to B_n$ 上にある
$\varphi(x_n) \to y$ であって，図式
$$
\xymatrix{
\varphi(x_n) \ar[r] \ar[d] & y \ar[d] \\
\varphi(x_{n - 1}) \ar[r]^-a \ar[dr] & f_*y \ar[d] \\
& \varphi(x)
}
$$
を可換にするものを得る。このとき合成
$x_n \to x_{n - 1} \to x$ と $\varphi(x_n) \to y$ が，
滑らかさの定義で要求される射である。
\end{proof}

\begin{remark}
\label{formal-defos-remark-smoothness-2-categorical}
$\varphi : \mathcal{F} \to \mathcal{G}$ を，
$\mathcal{C}_\Lambda$ 上で群亜群をファイバーとする余ファイバー圏の射とし，
$B \to A$ を $\mathcal{C}_\Lambda$ における環写像とする。
ファイバー圏 $\mathcal{F}(B)$ と $\mathcal{G}(B)$ の対象について
$B \to A$ に沿う順像を選択すると，$2$-可換図式
$$
\xymatrix{
\mathcal{F}(B) \ar[r]^{\varphi} \ar[d] & \mathcal{G}(B) \ar[d] \\
\mathcal{F}(A) \ar[r]^{\varphi}        & \mathcal{G}(A) .
}
$$
に入る関手 $\mathcal{F}(B) \to \mathcal{F}(A)$ および
$\mathcal{G}(B) \to \mathcal{G}(A)$ が定まる。
したがって誘導関手
$\mathcal{F}(B) \to \mathcal{F}(A)
\times_{\mathcal{G}(A)} \mathcal{G}(B)$ がある。
定義を展開すると，$\varphi : \mathcal{F} \to \mathcal{G}$ が
滑らかであるための必要十分条件は，$B \to A$ が全射であるときは常に
この誘導関手が本質的全射となることである（あるいは同値なこととして，
Lemma \ref{formal-defos-lemma-smoothness-small-extensions} により，
$B \to A$ が小拡大であるときは常にそうなることである）。
\end{remark}

\begin{remark}
\label{formal-defos-remark-compare-smooth-schlessinger}
Remark \ref{formal-defos-remark-smoothness-2-categorical} における滑らかな射の特徴付けは，
Schlessinger による関手の滑らかな射の概念に類似している。
cf.\ \cite[Definition 2.2.]{Sch} を参照せよ。
実際，$\mathcal{F}$ と $\mathcal{G}$ が集合をファイバーとする余ファイバー圏
であるとき，本稿の概念は Schlessinger の概念と同値である。
すなわち，この場合，対応する関手を
$F, G : \mathcal{C}_\Lambda \to \textit{Sets}$ とする。Remarks
\ref{formal-defos-remarks-cofibered-groupoids}
(\ref{formal-defos-item-convention-cofibered-sets}) を参照せよ。
このとき $F \to G$ が滑らかであるための必要十分条件は，
$\mathcal{C}_\Lambda$ における任意の全射環写像 $B \to A$ に対して，
写像 $F(B) \to F(A) \times_{G(A)} G(B)$ が全射となることである。
\end{remark}

\begin{remark}
\label{formal-defos-remark-smooth-to-iso-classes}
$\mathcal{F}$ を $\mathcal{C}_\Lambda$ 上で群亜群をファイバーとする
余ファイバー圏とする。このとき射
$\mathcal{F} \to \overline{\mathcal{F}}$ は滑らかである。
実際，$f : B \to A$ を $\mathcal{C}_\Lambda$ における環写像とする。
$x \in \Ob(\mathcal{F}(A))$ とし，
$\overline{y} \in \overline{\mathcal{F}}(B)$ を，
$\overline{f_*y} = \overline{x}$ となる
$y \in \Ob(\mathcal{F}(B))$ の同型類とする。このとき単に
$x' = y$，$B \to A$ 上の暗黙の射 $x' = y \to x$，および等式
$\overline{x'} = \overline{y}$ を，Definition
\ref{formal-defos-definition-smooth-morphism} で提示された問題の解として取ればよい。
\end{remark}

\noindent
% P10-E0437: Japanese supplies the locative preposition missing from the frozen English sentence.
$R \to S$ が $\widehat{\mathcal{C}}_\Lambda$ における環写像であるとき，
関手 $\underline{S}, \underline{R}: \widehat{\mathcal{C}}_\Lambda
\to \textit{Sets}$ の間に誘導射
$\underline{S} \to \underline{R}$ が存在する。この状況では，制限
$\underline{S}|_{\mathcal{C}_\Lambda} \to
\underline{R}|_{\mathcal{C}_\Lambda}$ の滑らかさは馴染み深い概念である：

\begin{lemma}
\label{formal-defos-lemma-smooth-morphism-power-series}
$R \to S$ を $\widehat{\mathcal{C}}_\Lambda$ における環写像とする。
このとき誘導射
$\underline{S}|_{\mathcal{C}_\Lambda} \to
\underline{R}|_{\mathcal{C}_\Lambda}$ が滑らかであるための必要十分条件は，
$S$ が $R$ 上の冪級数環であることである。
\end{lemma}

\begin{proof}
$S$ が $R$ 上の冪級数環であると仮定する。
$S = R[[x_1, \ldots, x_n]]$ と書く。
$\underline{S}|_{\mathcal{C}_\Lambda} \to
\underline{R}|_{\mathcal{C}_\Lambda}$ の滑らかさは次を意味する
（Remark \ref{formal-defos-remark-compare-smooth-schlessinger} を参照）：
$B \to A$ を $\mathcal{C}_\Lambda$ における全射環写像，
$R \to B$ と $S \to A$ を環写像とし，実線からなる図式
$$
\xymatrix{
S \ar[r] \ar@{..>}[rd] & A \\
R \ar[u] \ar[r] & B \ar[u]
}
$$
が可換であるとする。このとき図式を可換にする点線の射が存在する。
（Algebra, Definition \ref{algebra-definition-formally-smooth}
との類似に注意せよ。）
点線の射を構成するため，$A$ における像が $x_i$ の $A$ における像に
等しい元 $b_i \in B$ を選ぶ。$x_i$ は $\mathfrak m_A$ の元へ写るので，
$b_i \in \mathfrak m_B$ であることに注意せよ。
したがって $x_i$ を $b_i$ に写す一意な $R$-代数写像
$R[[x_1, \ldots, x_n]] \to B$ が存在し，これを点線の射にできる。

\medskip\noindent
逆に，
$\underline{S}|_{\mathcal{C}_\Lambda} \to
\underline{R}|_{\mathcal{C}_\Lambda}$ が滑らかであると仮定する。
$x_1, \ldots, x_n \in S$ を，その像が $R$ 上の $S$ の相対余接空間
$\mathfrak m_S/(\mathfrak m_R S + \mathfrak m_S^2)$ の基底をなす元とする。
$T = R[[X_1, \ldots, X_n]]$ と置く。次の二つの同型
$$
S/(\mathfrak m_R S + \mathfrak m_S^2) \cong
R/\mathfrak m_R[x_1, \ldots, x_n]/(x_ix_j)
$$
および
$$
T/(\mathfrak m_R T + \mathfrak m_T^2) \cong
R/\mathfrak m_R[X_1, \ldots, X_n]/(X_iX_j).
$$
に注意せよ。$x_i$ の類を $X_i$ の類へ写すことによって得られる局所
$R$-代数同型
$S/(\mathfrak m_R S + \mathfrak m_S^2) \to
T/(\mathfrak m_R T + \mathfrak m_T^2)$ を取る。
$f_1 : S \to T/(\mathfrak m_R T + \mathfrak m_T^2)$ を合成
$S \to S/(\mathfrak m_R S + \mathfrak m_S^2)
\to T/(\mathfrak m_R T + \mathfrak m_T^2)$ とする。
$\underline{S}|_{\mathcal{C}_\Lambda} \to
\underline{R}|_{\mathcal{C}_\Lambda}$ が滑らかであるという仮定により，
$f_1$ を写像 $f_2 : S \to T/\mathfrak{m}_T^2$ へ持ち上げ，次に
$f_3 : S \to T/\mathfrak{m}_T^3$ へ，さらに同様にしてすべての
$n \geq 1$ について持ち上げられる。したがって局所 $R$-代数の誘導写像
$f : S \to T = \lim T/\mathfrak m_T^n$ を得る。
$f_1$ の選び方により，$f$ は相対余接空間の同型
$\mathfrak m_S/(\mathfrak m_R S + \mathfrak m_S^2) \to
\mathfrak m_T/(\mathfrak m_R T + \mathfrak m_T^2)$ を誘導する。
したがって Lemma \ref{formal-defos-lemma-surjective-cotangent-space} により
$f$ は全射である（ここで $f$ を $\widehat{\mathcal{C}}_R$ における
写像とみなす）。$f$ による $X_i \in T$ の逆像
$y_i \in S$ を選ぶ。$T$ は $R$ 上の冪級数環なので，$X_i$ を $y_i$ に
写す局所 $R$-代数準同型 $s : T \to S$ が存在する。
構成により $f \circ s = \text{id}$ である。したがって $s$ は単射である。
しかし $f$ が相対余接空間上の同型を誘導するので $s$ もそうであり，
再び Lemma \ref{formal-defos-lemma-surjective-cotangent-space} により全射でもある。
ゆえに $s$ と $f$ は同型である。
\end{proof}

\noindent
滑らかな射は次の関手的性質を満たす。

\begin{lemma}
\label{formal-defos-lemma-smooth-properties}
$\varphi : \mathcal{F} \to \mathcal{G}$ および
$\psi : \mathcal{G} \to \mathcal{H}$ を，
$\mathcal{C}_\Lambda$ 上で群亜群をファイバーとする余ファイバー圏の射とする。
\begin{enumerate}
\item $\varphi$ と $\psi$ が滑らかならば，$\psi \circ \varphi$ は滑らかである。
\item $\varphi$ が本質的全射であり，$\psi \circ \varphi$ が滑らかならば，
$\psi$ は滑らかである。
\item $\mathcal{G}' \to \mathcal{G}$ が群亜群をファイバーとする
余ファイバー圏の射であり，$\varphi$ が滑らかならば，
$\mathcal{F} \times_\mathcal{G} \mathcal{G}' \to \mathcal{G}'$ は滑らかである。
\end{enumerate}
\end{lemma}

\begin{proof}
主張 (1) と (2) は定義から直ちに従う。(3) の証明は省略する。
ヒント：Remark \ref{formal-defos-remark-smoothness-2-categorical} で与えた滑らかさの
定式化，および $\mathcal{F} \times_\mathcal{G} \mathcal{G}'$ が
$2$-ファイバー積であることを用いる。Remarks
\ref{formal-defos-remarks-cofibered-groupoids} (\ref{formal-defos-item-fibre-product}) を参照せよ。
\end{proof}

\begin{lemma}
\label{formal-defos-lemma-smooth-morphism-essentially-surjective}
$\varphi : \mathcal{F} \to \mathcal{G}$ を，
$\mathcal{C}_\Lambda$ 上で群亜群をファイバーとする余ファイバー圏の
滑らかな射とする。$\varphi : \mathcal{F}(k) \to \mathcal{G}(k)$ が
本質的全射であると仮定する。このとき
$\varphi : \mathcal{F} \to \mathcal{G}$ および
$\widehat{\varphi} : \widehat{\mathcal{F}} \to
\widehat{\mathcal{G}}$ は本質的全射である。
\end{lemma}

\begin{proof}
$y$ を $A \in \Ob(\mathcal{C}_\Lambda)$ 上にある
$\mathcal{G}$ の対象とする。$A \to k$ に沿う $y$ の順像を
$y \to y_0$ とする。
$\varphi : \mathcal{F}(k) \to \mathcal{G}(k)$ の本質的全射性の仮定により，
$k$ 上にある $\mathcal{F}$ の対象 $x_0$ と同型
$y_0 \to \varphi(x_0)$ が存在する。$\varphi$ の滑らかさにより，
$A$ 上の $\mathcal{F}$ の対象 $x$ であって，像 $\varphi(x)$ が
$y$ と同型になるものが存在する。したがって
$\varphi : \mathcal{F} \to \mathcal{G}$ は本質的全射である。

\medskip\noindent
$\eta = (R, \eta_n, g_n)$ を $\widehat{\mathcal{G}}$ の対象とする。
同型 $\eta \to \varphi(\xi)$ をもつ
$\widehat{\mathcal{F}}$ の対象 $\xi$ を構成する。
$\varphi : \mathcal{F}(k) \to \mathcal{G}(k)$ の本質的全射性の仮定により，
ある $\xi_1 \in \Ob(\mathcal{F}(k))$ に対して
$\mathcal{G}(k)$ の射 $\eta_1 \to \varphi(\xi_1)$ が存在する。
射 $\eta_2 \xrightarrow{g_1} \eta_1 \to \varphi(\xi_1)$ は全射環写像
$R/\mathfrak m_R^2 \to k$ 上にあるので，$\varphi$ の滑らかさにより，
$\xi_2 \in \Ob(\mathcal{F}(R/\mathfrak m_R^2))$ と，
$R/\mathfrak m_R^2 \to k$ 上にある射
$f_1: \xi_2 \to \xi_1$，および射
$\eta_2 \to \varphi(\xi_2)$ であって，図式
$$
\xymatrix{
\varphi(\xi_2)  \ar[r]^{\varphi(f_1)} &  \varphi(\xi_{1})   \\
\eta_2   \ar[u] \ar[r]^{g_1}  & \eta_1  \ar[u] \\
}
$$
を可換にするものが存在する。同様に続けると，
$\widehat{\mathcal{F}}$ の対象 $\xi = (R, \xi_n, f_n)$ と，
$\widehat{\mathcal{G}}(R)$ における射
$\eta \to \varphi(\xi) = (R, \varphi(\xi_n), \varphi(f_n))$
を構成できる。
\end{proof}

\noindent
後に，pro-表現可能関手から前変形圏 $\mathcal{F}$ への滑らかな射を
構成することに関心をもつ。Remark
\ref{formal-defos-remark-formal-objects-yoneda} の議論により，これらの射は
$\mathcal{F}$ のある種の形式対象に対応する。より正確には，
これらはいわゆる $\mathcal{F}$ の versal 形式対象である。

\begin{definition}
\label{formal-defos-definition-versal}
$\mathcal{F}$ を群亜群をファイバーとする余ファイバー圏とする。
$\xi$ を，$R \in \Ob(\widehat{\mathcal{C}}_\Lambda)$ 上にある
$\mathcal{F}$ の形式対象とする。Remark
\ref{formal-defos-remark-formal-objects-yoneda} の対応する射
$\underline{\xi}: \underline{R}|_{\mathcal{C}_\Lambda} \to
\mathcal{F}$ が滑らかであるとき，$\xi$ を{\it versal} という。
\end{definition}

\begin{remark}
\label{formal-defos-remark-versal-object}
$\mathcal{F}$ を $\mathcal C_\Lambda$ 上で群亜群をファイバーとする
余ファイバー圏とし，$\xi$ を $\mathcal{F}$ の形式対象とする。
滑らかさの定義から，$\xi$ が versal であることは次の条件と同値である：
$$
\xymatrix{
& y \ar[d] \\
\xi \ar[r] & x
}
$$
が $\widehat{\mathcal{F}}$ における図式であり，$y \to x$ が
Artin 環の全射 $B \to A$ 上にあるとする（これを小拡大と仮定してよい）。
このとき，図式
$$
\xymatrix{
& y \ar[d] \\
\xi \ar[r] \ar[ur] & x
}
$$
を可換にする射 $\xi \to y$ が存在する。
とくに，$\xi$ が versal であるという条件は，Remark
\ref{formal-defos-remark-formal-objects-yoneda} における
$\underline{\xi} : \underline{R}|_{\mathcal{C}_\Lambda} \to
\mathcal{F}$ の構成で行った順像の選択に依存しない。
\end{remark}

\begin{lemma}
\label{formal-defos-lemma-versal-object-quasi-initial}
$\mathcal{F}$ を前変形圏とする。
$\xi$ を $\mathcal{F}$ の versal 形式対象とする。
% P10-E0438: preserve the frozen double-level wording.
$\widehat{\mathcal{F}}$ の任意の形式対象 $\eta$ に対し，射
$\xi \to \eta$ が存在する。
\end{lemma}

\begin{proof}
仮定により射
$\underline{\xi} : \underline{R}|_{\mathcal{C}_\Lambda} \to
\mathcal{F}$ は滑らかである。Remark
\ref{formal-defos-remark-formal-objects-yoneda} により
$\iota(\xi) : \underline{R} \to \widehat{\mathcal{F}}$ は
$\underline{\xi}$ の完備化である。Lemma
\ref{formal-defos-lemma-smooth-morphism-essentially-surjective} により，
$\iota(\xi)(f) = \eta$ となる $\underline{R}$ の対象 $f$ が存在する。
このとき $f$ は $\widehat{\mathcal{C}}_\Lambda$ における環写像
$f : R \to S$ である。また $\iota(\xi)(f) = \eta$ とは
$f_*\xi \cong \eta$ を意味し，これはまさに $f$ 上にある射
$\xi \to \eta$ が存在することを意味する。
\end{proof}



\section{滑らかな圏または非障害的な圏}
\label{formal-defos-section-smooth}

\noindent
$p : \mathcal{F} \to \mathcal{C}_\Lambda$ を群亜群をファイバーとする
余ファイバー圏とする。恒等関手を用いて，$\mathcal{C}_\Lambda$ を
$\mathcal{C}_\Lambda$ 上で群亜群をファイバーとする余ファイバー圏と
みなせる。このようにして
$p : \mathcal{F} \longrightarrow \mathcal{C}_\Lambda$ は，
$\mathcal{C}_\Lambda$ 上で群亜群をファイバーとする余ファイバー圏の射となる。

\begin{definition}
\label{formal-defos-definition-cofibered-groupoid-projection-smooth}
$p : \mathcal{F} \to \mathcal{C}_\Lambda$ を群亜群をファイバーとする
余ファイバー圏とする。その構造射 $p$ が Definition
\ref{formal-defos-definition-smooth-morphism} の意味で滑らかであるとき，
$\mathcal{F}$ は{\it 滑らか}または{\it 非障害的}であるという。
\end{definition}

\noindent
これは $\mathcal{C}_\Lambda$ 上で群亜群をファイバーとする余ファイバー圏の
滑らかさの「絶対的」概念であるが，より正確には
$\mathcal{F}$ は $\Lambda$ 上滑らかであるというべきであろう。
「{\it $\mathcal{F}$ は非障害的である}」という表現には注意が必要である：
$\mathcal{F}$ が滑らかであっても，非零の障害空間をもつ障害理論を
$\mathcal{F}$ が備えることがある。

\begin{remark}
\label{formal-defos-remark-smooth-on-top}
$\mathcal{F}$ を前変形圏とし，前変形圏 $\mathcal{U}$ からの滑らかな射
$\varphi : \mathcal U \to \mathcal{F}$ をもつと仮定する。
Lemma \ref{formal-defos-lemma-smooth-morphism-essentially-surjective} により
$\varphi$ は本質的全射である。したがって Lemma
\ref{formal-defos-lemma-smooth-properties} により，
$p: \mathcal{F} \to \mathcal{C}_\Lambda$ が滑らかであることと，合成
$\mathcal U \xrightarrow{\varphi} \mathcal{F} \xrightarrow{p}
\mathcal{C}_\Lambda$ が滑らかであることは同値である。
すなわち，\ $\mathcal{F}$ が滑らかであることと
$\mathcal{U}$ が滑らかであることは同値である。
\end{remark}

\begin{lemma}
\label{formal-defos-lemma-smooth}
$R \in \Ob(\widehat{\mathcal{C}}_\Lambda)$ とする。次は同値である：
\begin{enumerate}
\item $\underline{R}|_{\mathcal{C}_\Lambda}$ は滑らかである。
\item $\Lambda \to R$ は $\mathfrak m_R$-進位相について形式的に滑らかである。
\item $\Lambda \to R$ は平坦であり，
$R \otimes_\Lambda k'$ は $k'$ 上幾何学的正則である。
\item $\Lambda \to R$ は平坦であり，
$k' \to R \otimes_\Lambda k'$ は $\mathfrak m_R$-進位相について
形式的に滑らかである。
\end{enumerate}
古典的な場合，これらは次とも同値である：
\begin{enumerate}
\item[(5)] ある $n$ に対して
$R$ は $\Lambda[[x_1, \ldots, x_n]]$ と同型である。
\end{enumerate}
\end{lemma}

\begin{proof}
$p : \underline{R}|_{\mathcal{C}_\Lambda} \to
\mathcal{C}_\Lambda$ の滑らかさは，$\mathcal{C}_\Lambda$ における
全射 $B \to A$ と写像 $R \to A$ が与えられたとき，図式
$$
\xymatrix{
R \ar[r] \ar@{-->}[rd] & A \\
\Lambda \ar[r] \ar[u] & B \ar[u]
}
$$
の点線の射を見いだせることを意味する。
$\Lambda \to R$ が $\mathfrak m_R$-進位相について形式的に滑らかならば，
これは確かに成り立つ。More on Algebra, Definitions
\ref{more-algebra-definition-formally-smooth-adic} および
\ref{more-algebra-definition-formally-smooth} を参照せよ。
逆に，これが成り立つならば，More on Algebra, Lemma
\ref{more-algebra-lemma-fs-local} により
$\Lambda \to R$ は $\mathfrak m_R$-進位相について形式的に滑らかである。
したがって (1) と (2) は同値である。

\medskip\noindent
(2)，(3)，(4) の同値性は More on Algebra, Proposition
\ref{more-algebra-proposition-fs-flat-fibre-fs} である。
(5) との同値性は，たとえば Lemma
\ref{formal-defos-lemma-smooth-morphism-power-series} と，古典的な場合には
$\mathcal{C}_\Lambda$ が
$\underline{\Lambda}|_{\mathcal{C}_\Lambda}$ と同じであるという事実から従う。
\end{proof}

\begin{lemma}
\label{formal-defos-lemma-smooth-power-series-classical}
$\mathcal{F}$ を前変形圏とする。$\xi$ を，
$R \in \Ob(\widehat{\mathcal{C}}_\Lambda)$ 上にある
$\mathcal{F}$ の versal 形式対象とする。次は同値である：
\begin{enumerate}
\item $\mathcal{F}$ は非障害的である。
\item $\Lambda \to R$ は $\mathfrak m_R$-進位相について形式的に滑らかである。
\end{enumerate}
古典的な場合，これらは次とも同値である：
\begin{enumerate}
\item[(3)] ある $n$ に対して
$R \cong \Lambda[[x_1, \ldots, x_n]]$ である。\
% P10-E0439: preserve the frozen stray control-space token; correction is sidecar-only.
\end{enumerate}
\end{lemma}

\begin{proof}
(1) が成り立つ，すなわち $\mathcal{F}$ が非障害的であるとすると，合成
$$
\underline{R}|_{\mathcal{C}_\Lambda}
\xrightarrow{\underline{\xi}}
\mathcal{F}
\to
\mathcal{C}_\Lambda
$$
は滑らかである。Lemma \ref{formal-defos-lemma-smooth-properties} を参照せよ。
したがって Lemma \ref{formal-defos-lemma-smooth} により (2) が成り立つ。
逆に (2) が成り立つならばこの合成は滑らかであり，さらに Lemma
\ref{formal-defos-lemma-versal-object-quasi-initial} により第一の射は本質的全射である。
したがって Lemma \ref{formal-defos-lemma-smooth-properties} により第二の射は滑らかであり，
これは定義により $\mathcal{F}$ が非障害的であることを意味する。
古典的な場合の (3) との同値性は Lemma \ref{formal-defos-lemma-smooth} から従う。
\end{proof}

\begin{lemma}
\label{formal-defos-lemma-exists-smooth}
Lemma \ref{formal-defos-lemma-smooth} の同値な条件を満たし，さらに
$H_1(L_{k/\Lambda}) = \mathfrak m_R/\mathfrak m_R^2$ および
$\Omega_{R/\Lambda} \otimes_R k = \Omega_{k/\Lambda}$ を満たす
$R \in \Ob(\widehat{\mathcal{C}}_\Lambda)$ が存在する。
\end{lemma}

\begin{proof}
古典的な場合は $R = \Lambda$ と取る。より一般に，剰余体拡大
$k/k'$ が分離的ならば，$\Lambda$ の完備化 $\Lambda^\wedge$ の
一意な有限 \'etale 拡大 $\Lambda^\wedge \to R$ が存在し
（Algebra, Lemmas \ref{algebra-lemma-complete-henselian} および
\ref{algebra-lemma-henselian-cat-finite-etale}），剰余体上で拡大
$k/k'$ を誘導する。

\medskip\noindent
一般の場合は次のように進める。滑らかな $\Lambda$-代数 $P$ と
$\Lambda$-代数全射 $P \to k$ を選ぶ
（たとえば $P$ を多項式代数と取ればよい）。
$P \to k$ の核を $\mathfrak m_P$ と書く。
(\ref{formal-defos-equation-sequence-extended}) および Algebra, Lemma
\ref{algebra-lemma-exact-sequence-NL} を参照すると，
Jacobi--Zariski 完全列は完全列
$$
0 \to H_1(\NL_{k/\Lambda}) \to
\mathfrak m_P/\mathfrak m_P^2 \to
\Omega_{P/\Lambda} \otimes_P k \to
\Omega_{k/\Lambda} \to 0
$$
である。
% P10-E0440: preserve the frozen wrong smoothness base P/k.
左端が $0$ であるのは，$P/k$ が滑らかであり，したがって
$\NL_{P/\Lambda}$ が次数 $0$ に置かれた有限射影加群と準同型であり，
ゆえに $H_1(\NL_{P/\Lambda} \otimes_P k) = 0$ だからである。
$f \in \mathfrak m_P$ が
$\Omega_{P/\Lambda} \otimes_P k$ の非零元へ写ると仮定する。
$P' = P/(f)$ と置けば，$\Lambda$-代数全射 $P' \to k$ がある。
$P'$ は $\mathfrak m_{P'}$ において滑らかであることに注意せよ。
これは More on Morphisms, Lemma
\ref{more-morphisms-lemma-slice-smooth-given-element} から従う。
そこで $P$ を $P'$ の主局所化で置き換えると，
$\dim(\mathfrak m_P/\mathfrak m_P^2)$ が減少する。
有限回繰り返すと，写像
$\mathfrak m_P/\mathfrak m_P^2 \to
\Omega_{P/\Lambda} \otimes_P k$ が零であると仮定でき，完全列は同型
$H_1(L_{k/\Lambda}) = \mathfrak m_P/\mathfrak m_P^2$ および
$\Omega_{P/\Lambda} \otimes_P k = \Omega_{k/\Lambda}$ に分かれる。

\medskip\noindent
$R$ を $P$ の $\mathfrak m_P$-進完備化とする。このとき
$R$ は $\widehat{\mathcal{C}}_\Lambda$ の対象である。
実際，これは完備局所 Noether 環であり
（Algebra, Lemma \ref{algebra-lemma-completion-Noetherian-Noetherian}），
その剰余体は $k$ と同一視される。$R$ が求めるものであると主張する。

\medskip\noindent
まず，写像 $P \to R$ は同型
$\mathfrak m_P/\mathfrak m_P^2 = \mathfrak m_R/\mathfrak m_R^2$ および
$\Omega_{P/\Lambda} \otimes_P k =
\Omega_{R/\Lambda} \otimes_R k$ を誘導することに注意せよ。
これは $\mathfrak m_P/\mathfrak m_P^2$ と
$\Omega_{P/\Lambda} \otimes_P k$ のいずれも
$\Lambda$-代数 $P/\mathfrak m_P^2$ のみに依存し，$R$ についても同様で，
$P/\mathfrak m_P^2 = R/\mathfrak m_R^2$ だからである。
Algebra, Lemma \ref{algebra-lemma-differential-mod-power-ideal} を参照せよ。
Jacobi--Zariski 完全列の関手性
(\ref{formal-defos-equation-sequence-functorial}) を用いると，
$P$ について同じことが成り立つので，
$H_1(L_{k/\Lambda}) = \mathfrak m_R/\mathfrak m_R^2$ および
$\Omega_{R/\Lambda} \otimes_R k = \Omega_{k/\Lambda}$ を得る。

\medskip\noindent
最後に，$\Lambda \to P$ は滑らかなので，Algebra, Proposition
\ref{algebra-proposition-smooth-formally-smooth} により
$\Lambda \to P$ は形式的に滑らかである。次に More on Algebra, Lemma
\ref{more-algebra-lemma-formally-smooth} により
$\Lambda \to P$ は $\mathfrak m_P$-進位相について形式的に滑らかである。
この性質は More on Algebra, Lemma
\ref{more-algebra-lemma-formally-smooth-completion} により完備化 $R$ に
継承され，証明が完了する。実際には，
$\underline{R}|_{\mathcal{C}_\Lambda}$ が滑らかなときは常に，
$R$ は $\Lambda$ 上の滑らかな代数の完備化と同型になることが分かるが，
これは用いない。
\end{proof}

\begin{example}
\label{formal-defos-example-smooth-explicit}
Lemma \ref{formal-defos-lemma-exists-smooth} における $R$ の，より明示的な例を与える。
$p$ を素数，$n \in \mathbf{N}$ とする。
$\Lambda = \mathbf{F}_p(t_1, t_2, \ldots, t_n)$ および
$k = \mathbf{F}_p(x_1, \ldots, x_n)$ とし，写像
$\Lambda \to k$ を $t_i \mapsto x_i^p$ で与える。このとき
$$
R = \Lambda[x_1, \ldots, x_n]^\wedge_{(x_1^p - t_1, \ldots, x_n^p - t_n)}
$$
と取れる。この例ではこれより「よく」することはできない。
すなわち，$\mathcal{C}_\Lambda$ を
$\widehat{\mathcal{C}}_\Lambda$ のより小さい滑らかな対象で近似することは
できない（写像
$\Omega_{R/\Lambda} \otimes_R k \to \Omega_{k/\Lambda}$ が全射なので，
$R$ の次元は少なくとも $n$ でなければならないと論じられる）。
この現象は後にさらに詳しく論じる。
\end{example}

\section{Schlessinger の条件}
\label{formal-defos-section-schlessinger-conditions}

\noindent
以下ではしばしば，圏 $\mathcal{C}_\Lambda$ における環のファイバー積
$A_1 \times_A A_2$ を考える。Example \ref{formal-defos-example-fibre-product} で見たように，
このようなファイバー積は必ずしも $\mathcal{C}_\Lambda$ の対象とは限らない。
しかし以下で現れるほとんどすべての場合には，二つの写像 $A_i \to A$ の
一方が全射であり，Lemma \ref{formal-defos-lemma-fiber-product-CLambda} により
$A_1 \times_A A_2$ は $\mathcal{C}_\Lambda$ の対象となる。
以下，この結果は断りなく用いる。

\medskip\noindent
$k$ 上の双対数環を $k[\epsilon]$ と書く。より一般に，$k$-ベクトル空間
$V$ に対し，台となるベクトル空間が $k \oplus V$ であり，積が
$(a, v) \cdot (a', v') = (aa', av' + a'v)$ で与えられる $k$-代数を
$k[V]$ と書く。$V = k$ のとき，$k[V]$ は $k$ 上の双対数環である。
任意の有限次元 $k$-ベクトル空間 $V$ に対し，環 $k[V]$ は
$\mathcal{C}_\Lambda$ に属する。

\begin{definition}
\label{formal-defos-definition-S1-S2}
$\mathcal C_\Lambda$ 上で群亜群をファイバーとする余ファイバー圏
$\mathcal{F}$ を取る。$\mathcal{F}$ に対する{\it 条件 (S1) および (S2)}を
次のように定義する：
\begin{enumerate}
\item[(S1)] $A_2 \to A$ が全射である $\mathcal{C}_\Lambda$ の図式の上にある，
$\mathcal{F}$ の任意の図式
$$
\vcenter{
\xymatrix{
           & x_2 \ar[d] \\
x_1 \ar[r] & x
}
}
\quad\text{次の図式の上にある}\quad
\vcenter{
\xymatrix{
           & A_2 \ar[d] \\
A_1 \ar[r] & A
}
}
$$
は，可換図式
$$
\vcenter{
\xymatrix{
y \ar[r] \ar[d] & x_2 \ar[d] \\
x_1 \ar[r]      & x
}
}
\quad\text{次の図式の上にある}\quad
\vcenter{
\xymatrix{
A_1 \times_A A_2 \ar[r] \ar[d] & A_2 \ar[d] \\
A_1 \ar[r]      & A.
}
}
$$
へ拡張できる。
\item[(S2)]
$\mathcal{C}_\Lambda$ の図式
$$
\xymatrix{
          & k[\epsilon] \ar[d] \\
A  \ar[r] & k.
}
$$
の上にある $\mathcal{F}$ の図式について，(S1) の条件が成り立つ。
さらに，$\mathcal{F}$ に二つの可換図式
$$
\vcenter{
\xymatrix{
y \ar[r]_c \ar[d]_a & x_\epsilon \ar[d]^e \\
x \ar[r]^d          & x_0
}
}
\quad\text{および}\quad
\vcenter{
\xymatrix{
y' \ar[r]_{c'} \ar[d]_{a'} & x_\epsilon \ar[d]^e \\
x \ar[r]^d                 & x_0
}
}
\quad\text{が次の図式の上にあるとする}\quad
\vcenter{
\xymatrix{
A \times_k k[\epsilon] \ar[r] \ar[d] & k[\epsilon] \ar[d] \\
A  \ar[r] & k
}
}
$$
このとき，$a = a' \circ b$ を満たす射
$b : y \to y'$ が $\mathcal{F}(A \times_k k[\epsilon])$ に存在する。
\end{enumerate}
\end{definition}

\noindent
条件 (S1) と (S2) の意味は，ファイバー圏を用いて部分的に説明できる。
$f_1 : A_1 \to A$ と $f_2 : A_2 \to A$ を $\mathcal{C}_\Lambda$ の
環写像とし，$f_2$ は全射であるとする。射影を
$p_i : A_1 \times_A A_2 \to A_i$ と書く。
$\mathcal{F}$ の順像を一つずつ選んであると仮定する。
すると環の可換図式は，ファイバー圏の $2$-可換図式
$$
\xymatrix{
\mathcal{F}(A_1 \times_A A_2) \ar[r]_-{p_{2, *}} \ar[d]_{p_{1, *}} &
\mathcal{F}(A_2) \ar[d]^{f_{2, *}} \\
\mathcal{F}(A_1) \ar[r]^{f_{1, *}} & \mathcal{F}(A)
}
$$
へ移り，したがって圏の $2$-ファイバー積への関手
\begin{equation}
\label{formal-defos-equation-compare}
\mathcal{F}(A_1 \times_A A_2) \to
\mathcal{F}(A_1) \times_{\mathcal{F}(A)} \mathcal{F}(A_2)
\end{equation}
を与える。条件 (S1) はこの関手が本質的全射であることを要求する。
% P10-E0441: preserve the frozen indefinite-article defect sidecar-only.
条件 (S2) の前半は，$f_2$ が写像 $k[\epsilon] \to k$ に等しい場合に，
この関手が本質的全射であることを要求する。さらにこの場合，
(S2) の後半は，標的で同型になる二対象が始域でも同型であることを含意する
（ただし，この主張と{\it 同値ではない}）。定義の形で条件を述べる利点は，
何も選択する必要がないことである。

\begin{lemma}
\label{formal-defos-lemma-S1-small-extensions}
$\mathcal C_\Lambda$ 上で群亜群をファイバーとする余ファイバー圏
$\mathcal{F}$ を取る。$A_2 \to A$ が小拡大である場合に限って
(S1) の条件が成り立つと仮定すれば，$\mathcal{F}$ は (S1) を満たす。
\end{lemma}

\begin{proof}
証明は省略する。ヒント：Lemma \ref{formal-defos-lemma-factor-small-extension} を適用し，
Lemma \ref{formal-defos-lemma-smoothness-small-extensions} の証明と同様の帰納法を用いる。
\end{proof}

\begin{remark}
\label{formal-defos-remark-compare-S1-S2-schlessinger}
$\mathcal{F}$ が集合をファイバーとする場合，条件 (S1) と (S2) は
Schlessinger の論文 \cite{Sch} における条件 (H1) と (H2) にちょうど一致する。
すなわち，関手 $F: \mathcal{C}_\Lambda \to \textit{Sets}$ に対し，
条件 (S1) と (S2) は次を述べる：
\begin{enumerate}
\item [(S1)] $A_1 \to A$ と $A_2 \to A$ が $\mathcal{C}_\Lambda$ の
写像で，$A_2 \to A$ が全射ならば，誘導写像
$F(A_1 \times_A A_2) \to F(A_1) \times_{F(A)} F(A_2)$ は全射である。
\item [(S2)] $A \to k$ が $\mathcal{C}_\Lambda$ の写像ならば，誘導写像
$F(A \times_k k[\epsilon]) \to F(A) \times_{F(k)} F(k[\epsilon])$
は全単射である。
\end{enumerate}
写像 $F(A \times_k k[\epsilon]) \to
F(A) \times_{F(k)} F(k[\epsilon])$ の単射性は，条件 (S2) の後半と，
すべての射が恒等射であることから従う。
\end{remark}

\begin{lemma}
\label{formal-defos-lemma-S2-extensions}
$\mathcal{F}$ を $\mathcal{C}_\Lambda$ 上で群亜群をファイバーとする
余ファイバー圏とする。$\mathcal{F}$ が (S2) を満たすならば，任意の有限次元
$k$-ベクトル空間 $V$ に対し，$k[\epsilon]$ を $k[V]$ に置き換えても
(S2) の条件は成り立つ。
\end{lemma}

\begin{proof}
$\mathcal{F}$ が集合をファイバーとする場合，すなわち関手
$F : \mathcal{C}_\Lambda \to \textit{Sets}$ に対応する場合，これは
Remark \ref{formal-defos-remark-compare-S1-S2-schlessinger} における $F$ に対する
(S2) の記述と，
$k[V] \cong k[\epsilon] \times_k \ldots \times_k k[\epsilon]$
（因子は $\dim_k V$ 個）であることから従う。本章の残りで用いるのは
関手の場合である。

\medskip\noindent
一般の場合を $\dim(V)$ に関する帰納法で証明する。$\dim(V) = 1$ ならば
$k[V] \cong k[\epsilon]$ であり，主張は仮定から従う。
$\dim(V) > 1$ ならば $V = V' \oplus k\epsilon$ と書く。図式
$$
\vcenter{
\xymatrix{
& x_V \ar[d] \\
x \ar[r] & x_0
}
}
\quad\text{次の図式の上にある}\quad
\vcenter{
\xymatrix{
& k[V] \ar[d] \\
A \ar[r] & k
}
}
$$
を取る。$k[V] \to k[V']$ 上にある射 $x_V \to x_{V'}$ と，
$k[V] \to k[\epsilon]$ 上にある射 $x_V \to x_\epsilon$ を選ぶ。
射 $x_V \to x_0$ は $x_V \to x_{V'} \to x_0$ と
$x_V \to x_\epsilon \to x_0$ の両方に分解することに注意せよ。
帰納法の仮定により，図式
$$
\vcenter{
\xymatrix{
y' \ar[d] \ar[r] & x_{V'} \ar[d] \\
x \ar[r] & x_0
}
}
\quad\text{次の図式の上にある}\quad
\vcenter{
\xymatrix{
A \times_k k[V'] \ar[d] \ar[r] & k[V'] \ar[d] \\
A \ar[r] & k
}
}
$$
を見いだせる。これにより可換図式
$$
\vcenter{
\xymatrix{
& x_\epsilon \ar[d] \\
y' \ar[r] & x_0
}
}
\quad\text{次の図式の上にある}\quad
\vcenter{
\xymatrix{
& k[\epsilon] \ar[d] \\
A \times_k k[V'] \ar[r] & k
}
}
$$
を得る。したがって (S2) により，可換図式
$$
\vcenter{
\xymatrix{
y \ar[d] \ar[r] & x_\epsilon \ar[d] \\
y' \ar[r] & x_0
}
}
\quad\text{次の図式の上にある}\quad
\vcenter{
\xymatrix{
(A \times_k k[V']) \times_k k[\epsilon] \ar[d] \ar[r] & k[\epsilon] \ar[d] \\
A \times_k k[V'] \ar[r] & k
}
}
$$
を得る。ここで
$(A \times_k k[V']) \times_k k[\epsilon] = A \times_k k[V' \oplus k\epsilon]
= A \times_k k[V]$ である。$y$ が所要の可換図式に入ることを示す。
そのため，$A \times_k k[V] \to k[V]$ 上にある射 $y \to y_V$ を取る。
群亜群をファイバーとする余ファイバー圏の公理により，射
$y \to y' \to x_{V'}$ と $y \to x_\epsilon$ は射 $y \to y_V$ を通して
分解できる。したがって $y_V$ と $x_V$ の双方が可換図式
$$
\vcenter{
\xymatrix{
y_V \ar[r] \ar[d] & x_\epsilon \ar[d] \\
x_{V'} \ar[r]          & x_0
}
}
\quad\text{および}\quad
\vcenter{
\xymatrix{
x_V \ar[r] \ar[d] & x_\epsilon \ar[d] \\
x_{V'} \ar[r]                 & x_0
}
}
$$
に入り，ゆえに (S2) の後半により，$y_V \to x_{V'}$ と
$x_V \to x_{V'}$ に両立する同型 $y_V \to x_V$ が存在する。
とくにこれは $x_0$ への写像とも両立する。したがって合成
$y \to y_V \to x_V$ は所要の可換図式
$$
\vcenter{
\xymatrix{
y \ar[r] \ar[d] & x_V \ar[d] \\
x \ar[r] & x_0
}
}
\quad\text{次の図式の上にある}\quad
\vcenter{
\xymatrix{
A \times_k k[V] \ar[d] \ar[r] & k[V] \ar[d] \\
A \ar[r] & k
}
}
$$
に入る。このようにして，$k[\epsilon]$ を $k[V]$ に置き換えても
$(S2)$ の前半が成り立つことが分かる。

\medskip\noindent
後半を証明するため，二つの可換図式
$$
\vcenter{
\xymatrix{
y \ar[r] \ar[d] & x_V \ar[d] \\
x \ar[r] & x_0
}
}
\quad\text{および}\quad
\vcenter{
\xymatrix{
y' \ar[r] \ar[d] & x_V \ar[d] \\
x \ar[r] & x_0
}
}
\quad\text{が次の図式の上にあるとする}\quad
\vcenter{
\xymatrix{
A \times_k k[V] \ar[d] \ar[r] & k[V] \ar[d] \\
A \ar[r] & k
}
}
$$
証明の最初の段落で導入した射 $x_V \to x_{V'} \to x_0$ と
$x_V \to x_\epsilon \to x_0$ を用いる。
$A \times_k k[V] \to A \times_k k[V']$ 上にある射
$y \to y_{V'}$ と $y' \to y'_{V'}$ を選ぶ。
余ファイバー圏の公理により，可換図式
$$
\vcenter{
\xymatrix{
y_{V'} \ar[r] \ar[d] & x_{V'} \ar[d] \\
x \ar[r] & x_0
}
}
\quad\text{および}\quad
\vcenter{
\xymatrix{
y'_{V'} \ar[r] \ar[d] & x_{V'} \ar[d] \\
x \ar[r] & x_0
}
}
\quad\text{が次の図式の上にある}\quad
\vcenter{
\xymatrix{
A \times_k k[V'] \ar[d] \ar[r] & k[V'] \ar[d] \\
A \ar[r] & k
}
}
$$
を見いだせる。帰納法の仮定により，射 $y_{V'} \to x$ と
$y'_{V'} \to x$ に両立する同型 $b : y_{V'} \to y'_{V'}$ を得る。
とくに，これは $x_0$ への射に両立する。すると可換図式
$$
\vcenter{
\xymatrix{
y \ar[r] \ar[d] & x_\epsilon \ar[d] \\
y'_{V'} \ar[r] & x_0
}
}
\quad\text{および}\quad
\vcenter{
\xymatrix{
y' \ar[r] \ar[d] & x_\epsilon \ar[d] \\
y'_{V'} \ar[r] & x_0
}
}
\quad\text{が次の図式の上にある}\quad
\vcenter{
\xymatrix{
A \times_k k[\epsilon] \ar[d] \ar[r] & k[\epsilon] \ar[d] \\
A \ar[r] & k
}
}
$$
を得る。ここで射 $y \to y'_{V'}$ は合成
$y \to y_{V'} \xrightarrow{b} y'_{V'}$ であり，射
$y \to x_\epsilon$ と $y' \to x_\epsilon$ はそれぞれ写像
$y \to x_V$ と $y' \to x_V$ と射 $x_V \to x_\epsilon$ の合成である。
したがって (S2) の後半は，$y'_{V'}$ への写像に両立する同型
$y \to y'$ の存在を保証する。とくに，これは $x$ への写像に両立する
（$b$ が $x$ への写像に両立していたからである）。
\end{proof}

\begin{lemma}
\label{formal-defos-lemma-S1-S2-associated-functor}
$\mathcal{F}$ を $\mathcal{C}_\Lambda$ 上で群亜群をファイバーとする
余ファイバー圏とする。
\begin{enumerate}
\item $\mathcal{F}$ が (S1) を満たすならば，
$\overline{\mathcal{F}}$ も (S1) を満たす。
\item $\mathcal{F}$ が (S2) を満たし，さらに次の条件の少なくとも一つが
成り立つならば，$\overline{\mathcal{F}}$ も (S2) を満たす：
\begin{enumerate}
\item $\mathcal{F}$ は前変形圏である，
\item 圏 $\mathcal{F}(k)$ は集合またはセトイドである，または
\item $k[\epsilon] \to k$ 上にある $\mathcal{F}$ の任意の射
$x_\epsilon \to x_0$ に対し，順像写像
$\text{Aut}_{k[\epsilon]}(x_\epsilon) \to \text{Aut}_k(x_0)$
が全射である。
\end{enumerate}
\end{enumerate}
\end{lemma}

\begin{proof}
$\mathcal{F}$ が (S1) をもつと仮定する。
$f_i : A_i \to A$ を $\mathcal{C}_\Lambda$ の環写像とし，
$f_2$ は全射であるとする。$x_i \in \mathcal{F}(A_i)$ を，順像
$f_{1, *}(x_1)$ と $f_{2, *}(x_2)$ が同型になるように取る。
% P10-E0442: preserve the frozen malformed naming construction sidecar-only.
すると，これらの両方に同型な $A$ 上の $\mathcal{F}$ の対象を $x$ と書け，
(S1) に現れる図式を得る。したがって $A_1 \times_A A_2$ 上の
$\mathcal{F}$ の対象 $y$ であって，$A_1$ および $A_2$ への順像が
それぞれ\ $x_1$ および\ $x_2$ に同型なものを得る。このようにして
(S1) が $\overline{\mathcal{F}}$ に対して成り立つことが分かる。

\medskip\noindent
$\mathcal{F}$ が (S2) をもつと仮定する。
$\overline{\mathcal{F}}$ に対する (S2) の前半は，上の議論と同様に従う。
$\overline{\mathcal{F}}$ に対する (S2) の後半は，
$\mathcal{C}_\Lambda$ の任意の環 $A$ に対して写像
$$
\overline{\mathcal{F}}(A \times_k k[\epsilon]) \to
\overline{\mathcal{F}}(A)
\times_{\overline{\mathcal{F}}(k)} \overline{\mathcal{F}}(k[\epsilon])
$$
が単射であることを意味する。
$y, y' \in \mathcal{F}(A \times_k k[\epsilon])$ とする。
余ファイバー圏の公理を用いて，可換図式
$$
\vcenter{
\xymatrix{
y \ar[r]_c \ar[d]_a & x_\epsilon \ar[d]^e \\
x \ar[r]^d          & x_0
}
}
\quad\text{および}\quad
\vcenter{
\xymatrix{
y' \ar[r]_{c'} \ar[d]_{a'} & x'_\epsilon \ar[d]^{e'} \\
x' \ar[r]^{d'}                 & x'_0
}
}
\quad\text{が次の図式の上にあるように選べる}\quad
\vcenter{
\xymatrix{
A \times_k k[\epsilon] \ar[d] \ar[r] & k[\epsilon] \ar[d] \\
A \ar[r] & k
}
}
$$
$\mathcal{F}(A)$ の同型 $\alpha : x \to x'$ と
$\mathcal{F}(k[\epsilon])$ の同型
$\beta : x_\epsilon \to x'_\epsilon$ が存在すると仮定する。
これはまた，$\alpha$ に両立する同型
$\gamma : x_0 \to x'_0$ が存在することも意味する。
$\overline{\mathcal{F}}$ に対する (S2) を証明するには，
$\mathcal{F}(A \times_k k[\epsilon])$ に同型 $y \to y'$ が存在することを
示さなければならない。$\mathcal{F}$ に対する (S2) により，
同型 $\alpha$，$\beta$，$\gamma$ を図式
$$
\xymatrix{
x \ar[d]^\alpha \ar[r] & x_0 \ar[d]^\gamma &
x_\epsilon \ar[d]^\beta \ar[l]^e \\
x' \ar[r] & x'_0 & x'_\epsilon \ar[l]_{e'}
}
$$
が可換となるように選べれば，そのような射は存在する
（そのとき，先の表示図式で $x$ を $x'$ に，$x_\epsilon$ を
$x'_\epsilon$ に置き換えられるからである）。左側の正方形は
$\gamma$ の選び方により可換である。ある第二の写像
$\gamma' : x_0 \to x'_0$ により，$e' \circ \beta$ を
$\gamma' \circ e$ と分解できる。そこで
$\gamma = \gamma'$ となるようにできるかが問題となる。
$\mathcal{F}(k)$ が集合またはセトイドならば，これは明らかである。
さらに，$\text{Aut}_{k[\epsilon]}(x_\epsilon) \to
\text{Aut}_k(x_0)$ が全射ならば，像が
$\gamma^{-1} \circ \gamma'$ となる $x_\epsilon$ の自己同型を
前合成することで $\beta$ の選択を調整でき，所要の結果を得る。
\end{proof}

\begin{lemma}
\label{formal-defos-lemma-S1-S2-localize}
$\mathcal{F}$ を $\mathcal{C}_\Lambda$ 上で群亜群をファイバーとする
余ファイバー圏とする。$x_0 \in \Ob(\mathcal{F}(k))$ とする。
$\mathcal{F}_{x_0}$ を Remark \ref{formal-defos-remark-localize-cofibered-groupoid} で
構成した，$\mathcal{C}_\Lambda$ 上で群亜群をファイバーとする
余ファイバー圏とする。
\begin{enumerate}
\item $\mathcal{F}$ が (S1) を満たすならば，$\mathcal{F}_{x_0}$ も
(S1) を満たす。
\item $\mathcal{F}$ が (S2) を満たすならば，$\mathcal{F}_{x_0}$ も
(S2) を満たす。
\end{enumerate}
\end{lemma}

\begin{proof}
$\mathcal{F}_{x_0}$ における Definition \ref{formal-defos-definition-S1-S2} の形の
任意の図式は $\mathcal{F}$ の図式を与え，この図式に対する
$\mathcal{F}$ における条件 (S1) または (S2) の出力は，そのまま
$\mathcal{F}_{x_0}$ における出力とみなせる。
\end{proof}

\begin{lemma}
\label{formal-defos-lemma-lifting-section}
$p: \mathcal{F} \to \mathcal{C}_\Lambda$ を，群亜群をファイバーとする
余ファイバー圏とする。$\mathcal{F}$ の図式
$$
\vcenter{
\xymatrix{
y \ar[r] \ar[d]_a & x_\epsilon \ar[d]_e \\
x \ar[r]^d        & x_0
}
}
\quad\text{次の図式の上にある}\quad
\vcenter{
\xymatrix{
A \times_k k[\epsilon] \ar[r] \ar[d] & k[\epsilon] \ar[d] \\
A \ar[r] & k.
}
}
$$
を考える。ただし右側は $\mathcal{C}_\Lambda$ の図式である。
$\mathcal{F}$ が (S2) を満たすと仮定する。このとき
$a \circ s = \text{id}_x$ を満たす射 $s : x \to y$ が存在することと，
$e \circ s_\epsilon = d$ を満たす射
$s_\epsilon : x \to x_\epsilon$ が存在することは同値である。
\end{lemma}

\begin{proof}
「ならば」は明らかである。逆に，$e \circ s_\epsilon = d$ を満たす射
$s_\epsilon : x \to x_\epsilon$ が存在すると仮定する。
$p(s_\epsilon) : A \to k[\epsilon]$ は写像 $A \to k$ と両立する環写像で
あることに注意せよ。したがって
$$
\sigma = (\text{id}_A, p(s_\epsilon)) : A \to A \times_k k[\epsilon].
$$
を得る。順像 $x \to \sigma_*x$ を選ぶ。構成により，$s_\epsilon$ は
$x \to \sigma_*x \to x_\epsilon$ と分解できる。さらに $\sigma$ は
$A \times_k k[\epsilon] \to A$ の切断なので，$x \to \sigma_*x \to x$
が $\text{id}_x$ となる射 $\sigma_*x \to x$ を得る。
$e \circ s_\epsilon = d$ なので，図式
$$
\xymatrix{
\sigma_*x \ar[r] \ar[d] & x_\epsilon \ar[d]_e \\
x \ar[r]^d        & x_0
}
$$
は可換である。したがって (S2) により，$\sigma_*x \to y \to x$ が
与えられた写像 $\sigma_*x \to x$ となる射 $\sigma_*x \to y$ を得る。
% P10-E0443: preserve the frozen wrong map name a rather than s.
ここで問題の解として，$a : x \to y$ を合成
$x \to \sigma_*x \to y$ に等しく取ればよい。
\end{proof}

\begin{lemma}
\label{formal-defos-lemma-lifting-along-small-extension}
前変形圏 $\mathcal{F}$ の可換図式
$$
\vcenter{
\xymatrix{
y \ar[r] \ar[d] & x_2 \ar[d]^{a_2} \\
x_1 \ar[r]^{a_1}        & x
}
}
\quad\text{次の図式の上にある}
\vcenter{
\xymatrix{
A_1 \times_A A_2 \ar[r] \ar[d] & A_2 \ar[d]^{f_2} \\
A_1 \ar[r]^{f_1} & A
}
}
$$
を考える。ただし右側は $\mathcal{C}_\Lambda$ の図式であり，
$f_2 : A_2 \to A$ は小拡大である。
% P10-E0444: preserve the frozen ill-typed factorization sidecar-only.
$f_2 = f_1 \circ h$ を満たす写像 $h : A_1 \to A_2$ が存在すると仮定する。
$I = \Ker(f_2)$ とする。環写像
$$
g : A_1 \times_A A_2 \longrightarrow k[I] = k \oplus I, \quad
(u, v) \longmapsto \overline{u} \oplus (v - h(u))
$$
を考える。順像 $y \to g_*y$ を選び，$\mathcal{F}$ が (S2) を満たすと
仮定する。射 $x_1 \to g_*y$ が存在するならば，
$a_1 =  a_2 \circ b$ を満たす射 $b: x_1 \to x_2$ が存在する。
\end{lemma}

\begin{proof}
$\text{id}_{A_1} \times g : A_1 \times_A A_2 \to A_1 \times_k k[I]$
は同型であり，$k[I] \cong k[\epsilon]$ であることに注意せよ。
したがって図式
$$
\vcenter{
\xymatrix{
y \ar[r] \ar[d] & g_*y \ar[d] \\
x_1 \ar[r]        & x_0
}
}
\quad\text{次の図式の上にある}\quad
\vcenter{
\xymatrix{
A_1 \times_k k[\epsilon] \ar[r] \ar[d] & k[\epsilon] \ar[d] \\
A_1 \ar[r] & k.
}
}
$$
を得る。ここで $x_0$ は $k$ 上にある $\mathcal{F}$ の対象である
（$\mathcal{F}$ の各対象から $x_0$ への射は一意である。
Definition \ref{formal-defos-definition-predeformation-category} の後の議論を参照せよ）。
射 $x_1 \to g_*y$ があれば，Lemma \ref{formal-defos-lemma-lifting-section} により
写像 $y \to x_1$ の切断 $s : x_1 \to y$ を得る。
これと写像 $y \to x_2$ を合成して $b : x_1 \to x_2$ を得る。
補題の図式が可換であり，$s$ が切断であることから，これは
$a_1 =  a_2 \circ b$ を満たす。
\end{proof}

\section{関手の接空間}
\label{formal-defos-section-tangent-spaces-functors}

\noindent
$R$ を環とする。$R$-加群の圏を $\text{Mod}_R$，有限生成
$R$-加群の圏を $\text{Mod}^{fg}_R$ と書く。

\begin{definition}
\label{formal-defos-definition-linear}
$L: \text{Mod}^{fg}_R \to \text{Mod}_R$，それぞれ\ $L: \text{Mod}_R \to \text{Mod}_R$ を関手とする。
$\text{Mod}^{fg}_R$，それぞれ\ $\text{Mod}_R$ の任意の二対象
$M, N$ に対して写像
$$
L : \Hom_R(M, N) \longrightarrow \Hom_R(L(M), L(N))
$$
が $R$-加群の写像であるとき，$L$ は {\it $R$-線形}であるという。
\end{definition}

\begin{remark}
\label{formal-defos-remark-linear-enriched-over-modules}
$\text{Mod}_R$ 上豊穣化された圏の間の任意の関手に対して
$R$-線形性の概念を定義できる。ここではとくに関手
$L: \text{Mod}^{fg}_R \to \text{Mod}_R$ と
$L: \text{Mod}_R \to \text{Mod}_R$ に対して定義したが，
これはこれまで必要となったのがこの二つの場合だからである。
\end{remark}

\begin{remark}
\label{formal-defos-remark-linear-functor}
$L: \text{Mod}^{fg}_R \to \text{Mod}_R$ が $R$-線形関手ならば，
$L$ は有限積を保ち，零加群を零加群へ送る。Homology, Lemma
\ref{homology-lemma-additive-additive} を参照せよ。
一方，関手 $\text{Mod}^{fg}_R \to \textit{Sets}$ が有限積を保ち，
零加群を一点集合へ送るならば，これは $R$-線形関手への一意な持ち上げをもつ。
Lemma \ref{formal-defos-lemma-linear-functor} を参照せよ。
\end{remark}

\begin{lemma}
\label{formal-defos-lemma-linear-functor}
$L: \text{Mod}^{fg}_R \to \textit{Sets}$，それぞれ\ $L: \text{Mod}_R \to \textit{Sets}$ を関手とする。
$L(0)$ が一点集合で，$L$ が有限積を保つと仮定する。このとき
$$
\vcenter{
\xymatrix{
& \text{Mod}_R \ar[dr]^{\text{忘却}} &   \\
\text{Mod}^{fg}_R  \ar[ur]^{\widetilde{L}} \ar[rr]^{L} &  &
\textit{Sets}
}
}
\quad\text{それぞれ}\quad
\vcenter{
\xymatrix{
& \text{Mod}_R \ar[dr]^{\text{忘却}} &   \\
\text{Mod}_R  \ar[ur]^{\widetilde{L}} \ar[rr]^{L} &  &
\textit{Sets}
}
}
$$
が可換となるような，一意な $R$-線形関手
$\widetilde{L} : \text{Mod}^{fg}_R \to \text{Mod}_R$，それぞれ\ $\widetilde{L} : \text{Mod}_R \to \text{Mod}_R$ が存在する。
\end{lemma}

\begin{proof}
$L: \text{Mod}^{fg}_R \to \textit{Sets}$ の場合のみを証明する。
$M$ を有限生成 $R$-加群とする。集合 $L(M)$ に次の $R$-加群構造を
入れたものとして $\widetilde{L}(M)$ を定義する。

\medskip\noindent
乗法：$r \in R$ とする。$L(M)$ 上の $r$ 倍写像を，乗法写像
$r \cdot : M \to M$ により誘導される写像 $L(M) \to L(M)$ と定義する。

\medskip\noindent
加法：和写像 $M \times M \to M: (m_1, m_2) \mapsto m_1 + m_2$ は
写像 $L(M \times M) \to L(M)$ を誘導する。仮定により
$L(M \times M)$ は $L(M) \times L(M)$ と標準的に同型である。
$L(M)$ 上の加法を写像
$L(M) \times L(M) \cong L(M \times M) \to L(M)$ により定義する。

\medskip\noindent
零元：一意な写像 $0 \to M$ が存在する。$L(M)$ の零元を
$L(0) \to L(M)$ の像と定義する。

\medskip\noindent
% P10-E0445: preserve the frozen missing noun sidecar-only.
これにより $R$-加群 $\widetilde{L}(M)$ が定まり，しかも $M$ に関して
$R$-線形に関手的なものとして一意であることの検証は省略する。
\end{proof}

\begin{lemma}
\label{formal-defos-lemma-morphism-linear-functors}
$L_1, L_2: \text{Mod}^{fg}_R \to \textit{Sets}$ を，$0$ を一点集合へ送り，
有限積を保つ関手とする。$t : L_1 \to L_2$ を関手の射とする。
このとき $t$ は，Lemma \ref{formal-defos-lemma-linear-functor} が保証する関手の間の射
$\widetilde{t} : \widetilde{L}_1 \to \widetilde{L}_2$ を誘導し，これは各
$M \in \Ob(\text{Mod}^{fg}_R)$ に対して単に
$\widetilde{t}_M = t_M: \widetilde{L}_1(M) \to \widetilde{L}_2(M)$
で与えられる。言い換えると，
$t_M: \widetilde{L}_1(M) \to \widetilde{L}_2(M)$ は $R$-加群の写像である。
\end{lemma}

\begin{proof}
省略する。
\end{proof}

\noindent
$R = K$ が体である場合，$K$-線形関手
$L : \text{Mod}^{fg}_K \to \text{Mod}_K$ はその値 $L(K)$ により決定される。

\begin{lemma}
\label{formal-defos-lemma-linear-functor-over-field}
$K$ を体とし，$L: \text{Mod}^{fg}_K \to \text{Mod}_K$ を
$K$-線形関手とする。このとき $L$ は関手
$L(K) \otimes_K - : \text{Mod}^{fg}_K \to \text{Mod}_K$ と同型である。
\end{lemma}

\begin{proof}
$V \in \Ob(\text{Mod}^{fg}_K)$ に対し，同型
$L(K) \otimes_K V \to L(V)$ は単純テンソル上で
$x \otimes v \mapsto L(f_v)(x)$ により与えられる。ここで
$f_v: K \to V$ は $1 \mapsto v$ と送る $K$-線形写像である。
% P10-E0446: preserve the frozen malformed multiplication-by-K wording.
$V = K$ のとき，これは $K$ による乗法で与えられる同型
$L(K) \otimes_K K \to L(K)$ である。一般の $V$ に対しては，
$V = K$ の場合と，$L$ が有限積と可換であること
（Remark \ref{formal-defos-remark-linear-functor}）から同型となる。
\end{proof}

\noindent
$R$ を環，$M$ を $R$-加群とする。台となる $R$-加群が $R \oplus M$，
積が $(r, m) \cdot (r', m') = (rr', rm' + r'm)$ で与えられる
$R$-代数を $R[M]$ とする。$M = R$ のとき，これは $R$ 上の双対数環であり，
$R[\epsilon]$ と書く。

\medskip\noindent
次に $S$ を環とし，$R$ が $S$-代数であると仮定する。
対応 $M \mapsto R[M]$ は関手 $\text{Mod}_R \to S\text{-Alg}/R$ を定める。
ここで $S\text{-Alg}/R$ は $R$ 上の $S$-代数の圏を表す。
$S\text{-Alg}/R$ は有限積をもつことに注意せよ。すなわち
$A_1 \to R$ と $A_2 \to R$ が二対象ならば，
$A_1 \times_R A_2$ がその積である。

\begin{lemma}
\label{formal-defos-lemma-preserves-products}
$R$ を $S$-代数とする。このとき，上で述べた関手
$\text{Mod}_R \to S\text{-Alg}/R$ は有限積を保つ。
\end{lemma}

\begin{proof}
これは単に，$M$ と $N$ が $R$-加群ならば，写像
$R[M \times N] \to R[M] \times_R R[N]$ が $S\text{-Alg}/R$ の同型で
あるという主張である。
\end{proof}

\begin{lemma}
\label{formal-defos-lemma-tangent-space-functor}
$R$ を $S$-代数とし，$\mathcal{C}$ を $S\text{-Alg}/R$ の厳密充満部分圏で，
すべての $M \in \Ob(\text{Mod}^{fg}_R)$ に対する $R[M]$ を含むものとする。
$F: \mathcal{C} \to \textit{Sets}$ を関手とする。$F(R)$ が一点集合であり，
任意の $M, N \in \Ob(\text{Mod}^{fg}_R)$ に対する誘導写像
$$
F(R[M] \times_R R[N]) \to F(R[M]) \times F(R[N])
$$
が全単射であると仮定する。このとき，任意の
$M \in \Ob(\text{Mod}^{fg}_R)$ に対し $F(R[M])$ は自然な
$R$-加群構造をもつ。
\end{lemma}

\begin{proof}
$R \cong R[0]$ および $R[M] \times_R R[N] \cong R[M \times N]$ なので，
$\mathcal{C}$ に関する仮定により $R$ と $R[M] \times_R R[N]$ は
$\mathcal{C}$ の対象である。したがって $F$ に対する条件は意味をもつ。
Lemma \ref{formal-defos-lemma-preserves-products} の関手
$\text{Mod}_R \to S\text{-Alg}/R$ は，$\mathcal{C}$ に関する仮定により
関手 $\text{Mod}^{fg}_R \to \mathcal{C}$ に制限される。
$L$ を合成 $\text{Mod}^{fg}_R \to \mathcal{C} \to \textit{Sets}$，
すなわち $L(M) = F(R[M])$ とする。
Lemma \ref{formal-defos-lemma-preserves-products} と $F$ に関する仮定により，
$L$ は有限積を保つ。したがって Lemma \ref{formal-defos-lemma-linear-functor} により，
任意の $M \in \Ob(\text{Mod}^{fg}_R)$ に対して
$L(M) = F(R[M])$ は自然な $R$-加群構造をもつ。
\end{proof}

\begin{definition}
\label{formal-defos-definition-tangent-space-over-R}
$\mathcal{C}$ を Lemma \ref{formal-defos-lemma-tangent-space-functor} における圏とする。
$F : \mathcal{C} \to \textit{Sets}$ を，$F(R)$ が一点集合となる関手とする。
{\it $F$ の接空間 $TF$} とは $F(R[\epsilon])$ のことである。
\end{definition}

\noindent
$F : \mathcal{C} \to \textit{Sets}$ が
Lemma \ref{formal-defos-lemma-tangent-space-functor} の仮定を満たすとき，
接空間 $TF$ は自然な $R$-加群構造をもつ。

\begin{example}
\label{formal-defos-example-tangent-space-functor}
$\mathcal{C}_\Lambda$ は有限次元ベクトル空間 $V$ に対するすべての
$k[V]$ を含むので，Definition \ref{formal-defos-definition-tangent-space-over-R} は
$S = \Lambda$，$R = k$，$\mathcal{C} = \mathcal{C}_\Lambda$，および
$F : \mathcal{C}_\Lambda \to \textit{Sets}$ が前変形関手である場合に
適用できる。接空間は $TF = F(k[\epsilon])$ である。
\end{example}

\begin{example}
\label{formal-defos-example-tangent-space-prorepresentable-functor}
Example \ref{formal-defos-example-tangent-space-functor} の接空間を，
$F : \mathcal{C}_\Lambda \to \textit{Sets}$ が pro-表現可能関手，すなわち
ある $S \in \Ob(\widehat{\mathcal{C}}_\Lambda)$ に対する
$F = \underline{S}|_{\mathcal{C}_\Lambda}$ である場合に計算しよう。
すると $F$ は任意の極限と可換し，したがって
Lemma \ref{formal-defos-lemma-tangent-space-functor} の仮定を満たす。計算すると
$$
% P10-E0447: preserve both frozen wrong ambient-category subscripts.
TF = F(k[\epsilon]) = \Mor_{\mathcal{C}_\Lambda}(S, k[\epsilon])
= \text{Der}_\Lambda(S, k)
$$
を得る。より一般に，有限次元 $k$-ベクトル空間 $V$ に対し
$$
F(k[V]) = \Mor_{\mathcal{C}_\Lambda}(S, k[V]) = \text{Der}_\Lambda(S, V).
$$
を得る。明示的には，増大写像 $q : S \to k$ および $k[V] \to k$ と
両立する $\Lambda$-代数写像 $f : S \to k[V]$ は，
$s \mapsto f(s) - q(s)$ で定まる導分 $D$ に対応する。逆に，
$\Lambda$-導分 $D : S \to V$ は，規則 $f(s) = q(s) + D(s)$ で定まる
$\mathcal{C}_\Lambda$ の $f : S \to k[V]$ に対応する。
% P10-E0448: preserve the frozen plural-modifier defect sidecar-only.
これらの同一視は関手的なので，$TF$ と
$\text{Der}_\Lambda(S, k)$ 上の $k$-ベクトル空間構造は対応する
（Lemma \ref{formal-defos-lemma-morphism-linear-functors} を参照せよ）。
Lemma \ref{formal-defos-lemma-derivations-finite} により，$\dim_k TF$ は有限である。
\end{example}

\begin{example}
\label{formal-defos-example-tangent-space-classical-prorepresentable-functor}
Example \ref{formal-defos-example-tangent-space-prorepresentable-functor} の計算は
古典的な場合には簡単になる。すなわちこの場合，関手
$F = \underline{S}|_{\mathcal{C}_\Lambda}$ の接空間は単に
$\Lambda$ 上の $S$ の相対余接空間であり，式では $TF = T_{S/\Lambda}$ で
ある。実際，これは体拡大 $k/k'$ が分離的である場合にもより一般に成り立つ。
Exercises, Exercise \ref{exercises-exercise-tangent-space-Zariski} を参照せよ。
\end{example}

\begin{lemma}
\label{formal-defos-lemma-morphism-tangent-spaces}
$F, G: \mathcal{C} \to \textit{Sets}$ を
Lemma \ref{formal-defos-lemma-tangent-space-functor} の仮定を満たす関手とする。
$t : F \to G$ を関手の射とする。任意の
$M \in \Ob(\text{Mod}^{fg}_R)$ に対し，写像
$t_{R[M]}: F(R[M]) \to G(R[M])$ は $R$-加群の写像である。
ここで $F(R[M])$ と $G(R[M])$ には
Lemma \ref{formal-defos-lemma-tangent-space-functor} による $R$-加群構造を入れる。
とくに，$t_{R[\epsilon]} : TF \to TG$ は $R$-加群の写像である。
\end{lemma}

\begin{proof}
Lemma \ref{formal-defos-lemma-morphism-linear-functors} から従う。
\end{proof}

\begin{example}
\label{formal-defos-example-tangent-space-map-prorepresentable-functor}
$f : R \to S$ を $\widehat{\mathcal{C}}_\Lambda$ の環写像とする。
$F = \underline{R}|_{\mathcal{C}_\Lambda}$ および
$G = \underline{S}|_{\mathcal{C}_\Lambda}$ と置く。環写像 $f$ は
関手の変換 $G \to F$ を誘導する。Lemma \ref{formal-defos-lemma-morphism-tangent-spaces}
により，$k$-線形写像 $TG \to TF$ を得る。これは写像
$$
TG = \text{Der}_\Lambda(S, k) \longrightarrow \text{Der}_\Lambda(R, k) = TF
$$
である。これは Example \ref{formal-defos-example-tangent-space-prorepresentable-functor}
における標準的な同一視
$F(k[V]) = \text{Der}_\Lambda(R, V)$ および
$G(k[V]) = \text{Der}_\Lambda(S, V)$ と，接空間上の写像を計算する規則から従う。
\end{example}

\begin{lemma}
\label{formal-defos-lemma-tangent-space-tensor}
$F: \mathcal{C} \to \textit{Sets}$ を
Lemma \ref{formal-defos-lemma-tangent-space-functor} の仮定を満たす関手とする。
$R = K$ が体であると仮定する。このとき，任意の有限次元
$K$-ベクトル空間 $V$ に対し $F(K[V]) \cong TF \otimes_K V$ である。
\end{lemma}

\begin{proof}
Lemma \ref{formal-defos-lemma-linear-functor-over-field} から従う。
\end{proof}

\section{前変形圏の接空間}
\label{formal-defos-section-tangent-spaces}

\noindent
一般の Definition \ref{formal-defos-definition-tangent-space-over-R} を用いて
前変形関手の接空間を定義する。これは
Example \ref{formal-defos-example-tangent-space-functor} で詳しく述べた。
同型類の付随関手を考えることにより，前変形圏にも適用できる。

\begin{definition}
\label{formal-defos-definition-tangent-space}
$\mathcal{F}$ を前変形圏とする。{\it $\mathcal{F}$ の接空間
$T \mathcal{F}$} とは，ファイバー圏 $\mathcal F(k[\epsilon])$ の対象の
同型類からなる集合 $\overline{\mathcal{F}}(k[\epsilon])$ のことである。
\end{definition}

\noindent
したがって $T \mathcal{F}$ は，付随関手
$\overline{\mathcal{F}}: \mathcal{C}_\Lambda \to \textit{Sets}$ の
接空間にほかならない。$\mathcal{F}$ が (S2) を満たすとき，あるいは実際には
$\overline{\mathcal{F}}$ が満たすだけで，これは自然なベクトル空間構造をもつ。

\begin{lemma}
\label{formal-defos-lemma-tangent-space-vector-space}
$\mathcal{F}$ を，$\overline{\mathcal{F}}$ が (S2) を満たす前変形圏とする
\footnote{例えば $\mathcal{F}$ が (S2) を満たす場合。
Lemma \ref{formal-defos-lemma-S1-S2-associated-functor} を参照せよ。}。
このとき $T \mathcal{F}$ は自然な $k$-ベクトル空間構造をもつ。
任意の有限次元ベクトル空間 $V$ に対し，$V$ に関手的に
$\overline{\mathcal{F}}(k[V]) = T\mathcal{F} \otimes_k V$ が成り立つ。
\end{lemma}

\begin{proof}
$F = \overline{\mathcal{F}} : \mathcal{C}_\Lambda \to \textit{Sets}$ と書く。
これは前変形関手であり，$F$ は (S2) を満たす。
Lemma \ref{formal-defos-lemma-S2-extensions} と，
% P10-E0449: preserve the frozen opaque translation wording sidecar-only.
Remark \ref{formal-defos-remark-compare-S1-S2-schlessinger} の翻訳により，
任意の有限次元ベクトル空間 $V$ と任意の
$A \in \Ob(\mathcal{C}_\Lambda)$ に対して写像
$$
F(A \times_k k[V]) \longrightarrow F(A) \times F(k[V])
$$
は全単射である。とくに $A = k[W]$ ならば，
$F(k[W] \times_k k[V]) = F(k[W]) \times F(k[V])$ である。
言い換えると，Lemma \ref{formal-defos-lemma-tangent-space-functor} の仮定が成り立ち，
$TF = T \mathcal{F}$ は自然な $k$-ベクトル空間構造をもつ。
最後の主張は Lemma \ref{formal-defos-lemma-tangent-space-tensor} から従う。
\end{proof}

\noindent
前変形圏の射は接空間上の写像を誘導する。

\begin{definition}
\label{formal-defos-definition-differential}
% P10-E0450: preserve the frozen missing preposition sidecar-only.
$\varphi : \mathcal{F} \to \mathcal{G}$ を前変形圏の射とする。
{\it $\varphi$ の微分
$d \varphi : T \mathcal{F} \to T \mathcal{G}$} とは，関手の射
$\overline{\varphi}: \overline{\mathcal{F}} \to \overline{\mathcal{G}}$
を $A = k[\epsilon]$ で評価して得られる写像のことである。
\end{definition}

\begin{lemma}
\label{formal-defos-lemma-k-linear-differential}
$\varphi : \mathcal{F} \to \mathcal{G}$ を前変形圏の射とする。
$\overline{\mathcal{F}}$ と $\overline{\mathcal{G}}$ がともに (S2) を満たすと
仮定する。このとき $d \varphi : T \mathcal{F} \to T \mathcal{G}$ は
$k$-線形である。
\end{lemma}

\begin{proof}
Lemma \ref{formal-defos-lemma-tangent-space-vector-space} の証明で，
$\overline{\mathcal{F}}$ と $\overline{\mathcal{G}}$ が
Lemma \ref{formal-defos-lemma-tangent-space-functor} の仮定を満たすことを見た。
したがって補題は Lemma \ref{formal-defos-lemma-morphism-tangent-spaces} から従う。
\end{proof}

\begin{remark}
\label{formal-defos-remark-tangent-space-cofibered-groupoid}
接空間と微分の概念は，次のように，群亜群をファイバーとする任意の
余ファイバー圏へ大域化できる。$\mathcal{F}$ を $\mathcal{C}_\Lambda$ 上で
群亜群をファイバーとする余ファイバー圏とし，
$x \in \Ob(\mathcal{F}(k))$ とする。
Remark \ref{formal-defos-remark-localize-cofibered-groupoid} と同様に，前変形圏
$\mathcal{F}_x$ を得る。
$$
T_x\mathcal{F} = T\mathcal{F}_x
$$
を，{\it $x$ における $\mathcal{F}$ の接空間}と定義する。
$\varphi : \mathcal{F} \to \mathcal{G}$ が $\mathcal{C}_\Lambda$ 上で
群亜群をファイバーとする余ファイバー圏の射で，
$x \in \Ob(\mathcal{F}(k))$ ならば，誘導される射
$\varphi_x: \mathcal{F}_x \to \mathcal{G}_{\varphi(x)}$ が存在する。
{\it $x$ における $\varphi$ の微分
$d_x \varphi : T_x \mathcal{F} \to T_{\varphi(x)} \mathcal{G}$} を写像
$d \varphi_x: T \mathcal{F}_x \to T \mathcal{G}_{\varphi(x)}$ と定義する。
$\mathcal{F}$ と $\mathcal{G}$ がともに (S2) を満たすならば，
これらすべての接空間は自然な $k$-ベクトル空間構造をもち，すべての微分
$d_x \varphi : T_x \mathcal{F} \to T_{\varphi(x)} \mathcal{G}$ は
$k$-線形である（Lemmas \ref{formal-defos-lemma-S1-S2-localize} と
\ref{formal-defos-lemma-k-linear-differential} を用いよ）。
\end{remark}

\noindent
以下の観察は，古典的な場合，または $k/k'$ が分離体拡大である場合には
重要でない。実際，そのとき $\text{Der}_\Lambda(k, k)$ と
$\text{Der}_\Lambda(k, V)$ は零である。標準的な同一視
$$
\Mor_{\mathcal{C}_\Lambda}(k, k[\epsilon]) =
\text{Der}_\Lambda(k, k).
$$
が存在する。すなわち，$D \in \text{Der}_\Lambda(k, k)$ に対し，
$f_D : k \to k[\epsilon]$ を $a \mapsto a + D(a)\epsilon$ で定まる写像とする。
より一般に，$k$ 上の有限次元ベクトル空間 $V$ に対し
$$
\Mor_{\mathcal{C}_\Lambda}(k, k[V]) =
\text{Der}_\Lambda(k, V)
$$
であり，導分 $D$ に対応する写像にも同じ記号 $f_D$ を用いる。また
$$
% P10-E0451: preserve the frozen reversed Hom direction.
\Mor_{\mathcal{C}_\Lambda}(k[W], k[V]) =
\Hom_k(V, W) \oplus \text{Der}_\Lambda(k, V)
$$
である。ここで $(\varphi, D)$ は写像
$f_{\varphi, D} : a + w \mapsto a + \varphi(w) + D(a)$ に対応する。
% P10-E0452: preserve the frozen arrow rather than mapsto.
導分 $D : k \to V$ により定まる $k[V]$ の自己同型を，ときに
$f_{1, D} : a + v \to a + v + D(a)$ と書く。
$f_{1, D} \circ f_{1, D'} = f_{1, D + D'}$ であることに注意せよ。

\medskip\noindent
$\mathcal{F}$ を $\mathcal{C}_\Lambda$ 上の前変形圏とし，
$x_0 \in \Ob(\mathcal{F}(k))$ とする。上記により，
$D \mapsto f_{D, *}(x_0)$ で定まる標準写像
$$
\gamma_V :
\text{Der}_\Lambda(k, V)
\longrightarrow
\overline{\mathcal{F}}(k[V])
$$
が存在する。さらに，$(D, x) \mapsto f_{1, D, *}(x)$ で定まる作用
$$
a_V : \text{Der}_\Lambda(k, V) \times \overline{\mathcal{F}}(k[V])
\longrightarrow
\overline{\mathcal{F}}(k[V])
$$
が存在する。この二写像は両立する。すなわち，これらの写像の合成を
計算すれば $f_{1, D, *}f_{D', *}x_0 = f_{D + D', *}x_0$ である。
$x_0$ は同型を除いて一意なので，写像 $\gamma_V$ と $a_V$ は
$x_0$ の選択に依存しないことに注意せよ。

\begin{lemma}
\label{formal-defos-lemma-action-linear}
$\mathcal{F}$ を $\mathcal{C}_\Lambda$ 上の前変形圏とする。
$\overline{\mathcal{F}}$ が (S2) をもつならば，写像 $\gamma_V$ は
$k$-線形であり，$a_V(D, x) = x + \gamma_V(D)$ が成り立つ。
\end{lemma}

\begin{proof}
Lemma \ref{formal-defos-lemma-tangent-space-vector-space} の証明で，関手
$V \mapsto \overline{\mathcal{F}}(k[V])$ が $0$ を一点集合へ，積を積へ
送ることを見た。関手 $V \mapsto \text{Der}_\Lambda(k, V)$ も同様である。
したがって Lemma \ref{formal-defos-lemma-morphism-linear-functors} により
$\gamma_V$ は線形である。$D : k \to V$ を $\Lambda$-導分とする。
$D_1 : k \to V^{\oplus 2}$ を $a \mapsto (D(a), 0)$ と等しく置く。
すると
$$
\xymatrix{
k[V \times V] \ar[r]_{+} \ar[d]^{f_{1, D_1}} & k[V] \ar[d]^{f_{1, D}} \\
k[V \times V] \ar[r]^{+} & k[V]
}
$$
は可換である。定義を展開し，
% P10-E0453: preserve the frozen undefined overline F(V) and a_D notation.
$\overline{F}(V \times V) = \overline{F}(V) \times \overline{F}(V)$ を用いると，
これはすべての $x_1, x_2 \in \overline{F}(V)$ に対して
$a_D(x_1) + x_2 = a_D(x_1 + x_2)$ であることを意味する。
したがって，$0 \in \overline{F}(V)$ が零ベクトルであるとき，
$a_V(D, 0) = 0 + \gamma_V(D)$ を示せば十分である。
定義により，これは元 $f_{0, *}(x_0)$ である。
$f_D = f_{1, D} \circ f_0$ なので，所要の結果が従う。
\end{proof}

\noindent
上の構成の特別な場合として，任意の前変形圏 $\mathcal{F}$ に対して
定義される写像
\begin{equation}
\label{formal-defos-equation-map}
\gamma : \text{Der}_\Lambda(k, k) \longrightarrow T\mathcal{F}
\end{equation}
および作用
\begin{equation}
\label{formal-defos-equation-action}
a : \text{Der}_\Lambda(k, k) \times T\mathcal{F} \longrightarrow T\mathcal{F}
\end{equation}
がある。$\varphi : \mathcal{F} \to \mathcal{G}$ が前変形圏の射ならば，
可換図式
$$
\vcenter{
\xymatrix{
\text{Der}_\Lambda(k, k) \ar[r]_-\gamma \ar[rd]_\gamma &
T\mathcal{F} \ar[d]_{d\varphi} \\
& T\mathcal{G}
}
}
\quad\text{および}\quad
\vcenter{
\xymatrix{
\text{Der}_\Lambda(k, k) \times T\mathcal{F} \ar[r]_-a
\ar[d]_{1 \times d\varphi} &
T\mathcal{F} \ar[d]^{d\varphi} \\
\text{Der}_\Lambda(k, k) \times T\mathcal{G} \ar[r]^-a &
T\mathcal{G}
}
}
$$
を得ることに注意せよ。

\section{Versal 形式対象}
\label{formal-defos-section-versal-objects}

\noindent
versal 形式対象の存在は，$\mathcal{F}$ が性質 (S1) をもつことを強制する。

\begin{lemma}
\label{formal-defos-lemma-versal-object-S1}
$\mathcal{F}$ を前変形圏とする。$\mathcal{F}$ が versal 形式対象をもつと
仮定する。このとき $\mathcal{F}$ は (S1) を満たす。
\end{lemma}

\begin{proof}
$\xi$ を $\mathcal{F}$ の versal 形式対象とする。図式
$$
\xymatrix{
           & x_2 \ar[d] \\
x_1 \ar[r] & x
}
$$
を $\mathcal{F}$ の図式とし，$x_2 \to x$ は全射環写像の上にあるとする。
自然な射
$\widehat{\mathcal{F}}|_{\mathcal{C}_\Lambda} \xrightarrow{\sim} \mathcal{F}$
は同値なので（Remark \ref{formal-defos-remark-restrict-completion} を参照せよ），
この図式を $\widehat{\mathcal{F}}$ の図式とみなすこともできる。
Lemma \ref{formal-defos-lemma-versal-object-quasi-initial} により射 $\xi \to x_1$ が存在し，
したがって Remark \ref{formal-defos-remark-versal-object} により，図式
$$
\xymatrix{
\xi \ar[r] \ar[d]          & x_2 \ar[d] \\
x_1 \ar[r] & x
}
$$
を可換にする射 $\xi \to x_2$ も得る。$x_1 \to x$ と $x_2 \to x$ が
環写像 $A_1 \to A$ と $A_2 \to A$ の上にあるならば，
$\xi$ の $A_1 \times_A A_2$ への順像を取ることで，(S1) が要求する
対象 $y$ を得る。
\end{proof}

\noindent
余ファイバー圏が (S1) と (S2) を満たす場合，versal 形式対象を
次のように特徴付けられる。

\begin{lemma}
\label{formal-defos-lemma-versal-criterion}
$\mathcal{F}$ を (S1) と (S2) を満たす前変形圏とする。
$\xi$ を $\mathcal{F}$ の形式対象とし，
Remark \ref{formal-defos-remark-formal-objects-yoneda} の対応する射を
$\underline{\xi} : \underline{R}|_{\mathcal{C}_\Lambda} \to \mathcal{F}$
とする。このとき，$\xi$ が versal であることと，次の二条件は同値である：
\begin{enumerate}
\item 接空間上の写像
$d\underline{\xi} : T\underline{R}|_{\mathcal{C}_\Lambda} \to T\mathcal{F}$
は全射である；
\item $\widehat{\mathcal{F}}$ の図式
$$
\vcenter{
\xymatrix{
            &  y \ar[d] \\
\xi \ar[r]  &  x
}
}
\quad\text{次の図式の上にある}\quad
\vcenter{
\xymatrix{
         &   B  \ar[d]^{f} \\
R \ar[r] &   A
}
}
$$
が $\widehat{\mathcal{C}}_\Lambda$ の図式の上にあり，
$B \to A$ が Artin 環の小拡大であるとする。このとき，
$$
\xymatrix{
         &   B  \ar[d]^{f} \\
R \ar[ur] \ar[r] &   A
}
$$
を可換にする環写像 $R \to B$ が存在する。
\end{enumerate}
\end{lemma}

\begin{proof}
$\xi$ が versal ならば，(1) は
Lemma \ref{formal-defos-lemma-smooth-morphism-essentially-surjective} から，(2) は
Remark \ref{formal-defos-remark-versal-object} から従う。
(1) と (2) が成り立つと仮定する。
Remark \ref{formal-defos-remark-versal-object} により，(2) の形の
$\widehat{\mathcal{F}}$ の図式が与えられたとき，
$$
\xymatrix{
            &  y \ar[d] \\
\xi \ar[ur] \ar[r]  &  x
}
$$
を可換にする $\xi \to y$ が存在することを示せばよい。
$b : R \to B$ を (2) が保証する写像とする。$y' = b_*\xi$ と書き，
与えられた射 $\xi \to x$ の，$R \to B \to A$ 上にある分解
$\xi \to y' \to x$ を選ぶ。(S1) により，可換図式
$$
\vcenter{
\xymatrix{
z  \ar[r] \ar[d]          &  y \ar[d] \\
y' \ar[r]  &  x
}
}
\quad\text{次の図式の上にある}\quad
\vcenter{
\xymatrix{
B \times_A B \ar[d] \ar[r] &   B  \ar[d]^{f} \\
B \ar[r]^{f} &   A .
}
}
$$
を得る。$I = \Ker(f)$ と置く。環写像
$\overline{g} : B \times_A B \to k[I]$ を
$(u, v) \mapsto \overline{u} \oplus (v - u)$ と定める。
Lemma \ref{formal-defos-lemma-lifting-along-small-extension} と比較せよ。\ (1) により，ある環写像 $i : R \to k[\epsilon]$ の上にある射
$\xi \to \overline{g}_*z$ が存在する。Artin 商 $b_1 : R \to B_1$ を，
$b : R \to B$ と $i : R \to k[\epsilon]$ がともに $R \to B_1$ を通して
分解するように選ぶ。すなわち，$h : B_1 \to B$ と
$i' : B_1 \to k[\epsilon]$ を得る。順像 $y_1 = b_{1, *}\xi$，
$\xi \to y'$ の $R \to B_1 \to B$ 上にある分解
$\xi \to y_1 \to y'$，および $\xi \to \overline{g}_*z$ の
$R \to B_1 \to k[\epsilon]$ 上にある分解
$\xi \to y_1 \to \overline{g}_*z$ を選ぶ。
(S1) をもう一度適用すると，
$$
\vcenter{
\xymatrix{
z_1  \ar[r] \ar[d] & z \ar[r] \ar[d] & y \ar[d] \\
y_1 \ar[r] & y' \ar[r] &  x
}
}
\quad\text{次の図式の上にある}\quad
\vcenter{
\xymatrix{
B_1 \times_A B \ar[d] \ar[r] & B \times_A B \ar[r] \ar[d] & B \ar[d]^{f} \\
B_1 \ar[r]  & B \ar[r] &  A .
}
}
$$
を得る。Lemma \ref{formal-defos-lemma-lifting-along-small-extension} における写像
$g : B_1 \times_A B \to k[I]$（$h$ を用いて定義する）は，
$B_1 \times_A B \to B \times_A B$ と上の写像 $\overline{g}$ の合成である。
構成により射 $y_1 \to g_*z_1 \cong \overline{g}_*z$ が存在する。
したがって Lemma \ref{formal-defos-lemma-lifting-along-small-extension} を
上の図式の外側の長方形に適用して射 $y_1 \to y$ を得る。
これに $\xi \to y_1$ を前合成すれば，所要の射 $\xi \to y$ を得る。
\end{proof}

\noindent
$\mathcal{F}$ が性質 (S1) をもつならば，「持ち上げが存在する最大の商」が
存在する。正確な主張は次である。

\begin{lemma}
\label{formal-defos-lemma-largest-closed-where-lift}
$\mathcal{F}$ を $\mathcal{C}_\Lambda$ 上で群亜群をファイバーとする
(S1) をもつ余ファイバー圏とする。$B \to A$ を
$\mathcal{C}_\Lambda$ の全射とし，その核 $I$ は $\mathfrak m_B$ で
零化されるとする。$x \in \mathcal{F}(A)$ とする。イデアルの集合
$$
\mathcal{J} = \{ J \subset I \mid
% P10-E0455: preserve the frozen article defect sidecar-only.
\text{次の射が存在する： }y \to x\text{ は次の写像の上にある： }B/J \to A\}
$$
は最小元をもつ。
\end{lemma}

\begin{proof}
$I \in \mathcal{J}$ なので，$\mathcal{J}$ は空でない。
また，$J \in \mathcal{J}$ かつ $J \subset J' \subset I$ ならば，
対象 $y$ を $B/J'$ 上の対象 $y'$ へ順像できるので
$J' \in \mathcal{J}$ である。表示した集合の元 $J$ と $K$ を取る。
$J \cap K \in \mathcal{J}$ を示せば補題が従う。
$I$ は $k$-ベクトル空間なので，$J \cap K = J' \cap K$ かつ
$J' + K = I$ を満たすイデアル $J \subset J' \subset I$ を取れる。
上記により $J$ を $J'$ で置き換え，$J + K = I$ と仮定してよい。
この場合
$$
% P10-E0454: preserve the frozen wrong A-quotient symbols.
A/(J \cap K) = A/J \times_{A/I} A/K.
$$
したがって，仮定した $A/J$ と $A/K$ 上の元の存在から，(S1) により，
$x$ へ写る元 $z \in \mathcal{F}(A/(J \cap K))$ の存在が従う。
\end{proof}

\noindent
次の結果は後で改良する。

\begin{lemma}
\label{formal-defos-lemma-versal-object-existence}
$\mathcal{F}$ を $\mathcal{C}_\Lambda$ 上で群亜群をファイバーとする
余ファイバー圏とする。次の条件が成り立つと仮定する：
\begin{enumerate}
\item $\mathcal{F}$ は前変形圏である。
\item $\mathcal{F}$ は (S1) を満たす。
\item $\mathcal{F}$ は (S2) を満たす。
\item $\dim_k T\mathcal{F}$ は有限である。
\end{enumerate}
このとき $\mathcal{F}$ は versal 形式対象をもつ。
\end{lemma}

\begin{proof}
(1)，(2)，(3)，(4) が成り立つと仮定する。
$\underline{R}|_{\mathcal{C}_\Lambda}$ が滑らかとなる対象
$R \in \Ob(\widehat{\mathcal{C}}_\Lambda)$ を選ぶ。
Lemma \ref{formal-defos-lemma-exists-smooth} を参照せよ。
$r = \dim_k T\mathcal{F}$ とし，$S = R[[X_1, \ldots, X_r]]$ と置く。

\medskip\noindent
$n \geq 2$ に対して，対
$(J_n, f_{n - 1} : \xi_n \to \xi_{n - 1})$ を帰納的に構成する。
% P10-E0456: preserve the frozen wrong article sidecar-only.
ここで $J_n \subset S$ はイデアルの減少列であり，
$f_{n - 1} : \xi_n \to \xi_{n - 1}$ は射影
$S/J_n \to S/J_{n - 1}$ 上にある $\mathcal{F}$ の射である。

\medskip\noindent
第 1 段階。$J_1 = \mathfrak m_S$ と置く。
$\xi_1$ を $k = S/J_1 = S/\mathfrak m_S$ 上の $\mathcal{F}$ の
（一意な同型を除いて）一意な対象とする。
% P10-E0457: the frozen source lacks terminal punctuation here and below.

\medskip\noindent
第 2 段階。$J_2 = \mathfrak m_S^2 + \mathfrak{m}_R S$ と置く。このとき
$S/J_2 = k[V]$，ただし $V = kX_1 \oplus \ldots \oplus kX_r$ である。
$\overline{\mathcal{F}}$ に対する (S2) により，全単射
$$
\overline{\mathcal{F}}(S/J_2)
\longrightarrow
T\mathcal{F} \otimes_k V,
$$
を得る。Lemmas \ref{formal-defos-lemma-S1-S2-associated-functor} と
\ref{formal-defos-lemma-tangent-space-vector-space} を参照せよ。
$T\mathcal{F}$ の基底 $\theta_1, \ldots, \theta_r$ を選び，
% P10-E0458: preserve the frozen object/class equality and membership notation.
$\xi_2 = \sum \theta_i \otimes X_i \in \Ob(\mathcal{F}(S/J_2))$ と置く。
この選択の要点は，
$$
d\xi_2 :
\Mor_{\mathcal{C}_\Lambda}(S/J_2, k[\epsilon])
\longrightarrow
T\mathcal{F}
$$
が全射となることである。$f_1 : \xi_2 \to \xi_1$ を一意な射とする。

\medskip\noindent
帰納段階。ある $n \geq 2$ に対して
$(J_n, f_{n - 1} : \xi_n \to \xi_{n - 1})$ が構成済みであると仮定する。
次を満たすイデアル $J \subset S$ の集合には最小元 $J_{n + 1}$ が存在する：
(a) $\mathfrak m_S J_n \subset J \subset J_n$，かつ
(b) $S/J \to S/J_n$ 上にある射
$\xi_{n + 1} \to \xi_n$ が存在する。
Lemma \ref{formal-defos-lemma-largest-closed-where-lift} を参照せよ。
$f_n : \xi_{n + 1} \to \xi_n$ を，$S/J_{n + 1} \to S/J_n$ 上にある
$\mathcal{F}$ の任意の射とする。

\medskip\noindent
$J = \bigcap J_n$，$\overline{S} = S/J$，
$\overline{J}_n = J_n/J$ と置く。
Lemma \ref{formal-defos-lemma-m-adic-topology} により，イデアル列
$(\overline{J}_n)$ は $\overline{S}$ 上の
$\mathfrak m_{\overline{S}}$-進位相を誘導する。
$(\xi_n, f_n)$ は $\widehat{\mathcal{F}}_\mathcal{I}(\overline{S})$ の
対象である。ここで $\mathcal{I}$ は $\overline{S}$ のフィルトレーション
$(\overline{J}_n)$ である。したがって $(\xi_n, f_n)$ は
$\widehat{\mathcal{F}}(\overline{S})$ の対象 $\xi$ を誘導する。
% P10-E0459: preserve the frozen lowercase citation start sidecar-only.
Lemma \ref{formal-defos-lemma-formal-objects-different-filtration} を参照せよ。

\medskip\noindent
$\xi$ が versal であることを証明する。Versality のためには，
Lemma \ref{formal-defos-lemma-versal-criterion} の条件 (1) と (2) を確認すれば十分である。
条件 (1) は上の第 2 段階における $\xi_2$ の選択から従う。
$\widehat{\mathcal{F}}$ の図式
$$
\vcenter{
\xymatrix{
            &  y \ar[d] \\
\xi \ar[r]  &  x
}
}
\quad\text{次の図式の上にある}\quad
\vcenter{
\xymatrix{
         &   B  \ar[d]^{f} \\
\overline{S} \ar[r] &   A
}
}
$$
が与えられたとする。右側は $\widehat{\mathcal{C}}_\Lambda$ の図式で，
$f: B \to A$ は Artin 環の小拡大である。
右側の図式に入る写像 $\overline{S} \to B$ の存在を示さなければならない。
$\overline{S} \to A$ が $\overline{S} \to S/J_n$ を通して分解するような
$n$ を選ぶ。これは，上で見たように列 $(\overline{J}_n)$ が
$\mathfrak m_{\overline{S}}$-進位相を誘導するため可能である。
$\xi$ の $\overline{S} \to S/J_n$ に沿う順像は $\xi_n$ である。
$\xi \to x$ を $\xi \to \xi_n \to x$ と分解してよいので，
$\mathcal{F}$ の図式
$$
\vcenter{
\xymatrix{
            &  y \ar[d] \\
\xi_n \ar[r]  &  x
}
}
\quad\text{次の図式の上にある}\quad
\vcenter{
\xymatrix{
         &   B  \ar[d]^{f} \\
S/J_n \ar[r] &   A .
}
}
$$
を得る。Lemma \ref{formal-defos-lemma-versal-criterion} の条件 (2) を確認するには，
図式
$$
\xymatrix{
S/J_{n + 1} \ar[d] \ar@{-->}[r] & B \ar[d]^{f} \\
S/J_n   \ar[r] & A
}
$$
を完成させれば十分であり，これは同値に，図式
$$
\xymatrix{
  &  S/J_n \times_A B \ar[d]^{p_1} \\
S/J_{n + 1} \ar@{-->}[ur] \ar[r] &  S/J_n.
}
$$
を完成させることである。$p_1$ が切断をもてば証明は終わる。
そうでなければ，Lemma \ref{formal-defos-lemma-fiber-product-CLambda} (2) により
$p_1$ は小拡大であり，したがって
Lemma \ref{formal-defos-lemma-essential-surjection} (4) により $p_1$ は本質的全射である。
$S = R[[X_1, \ldots, X_r]]$ であり，
$\underline{R}|_{\mathcal{C}_\Lambda}$ が滑らかとなるように $R$ を
選んだことを思い出そう。したがって，写像
$R \to S \to S/J_n \to A$ を持ち上げる写像 $h : R \to B$ が存在する。
% P10-E0460: preserve the frozen reuse of h for the power-series extension.
冪級数環の普遍性により，与えられた写像
$S \to S/J_n \to A$ を持ち上げる $R$-代数写像
$h : S = R[[X_1, \ldots, X_r]] \to B$ が存在する。
これは，図式
$$
\xymatrix{
S \ar[d] \ar[r]_-g  &  S/J_n \times_A B \ar[d]^{p_1} \\
S/J_{n + 1} \ar@{-->}[ur] \ar[r] &  S/J_n
}
$$
の実線の正方形を可換にする写像 $g: S \to S/J_n \times_A B$ を誘導する。
$p_1$ は本質的全射なので，$g$ は全射である。
$S$ のイデアル $K = \Ker(g)$ が，上の帰納段階における $J_{n + 1}$ の
構成の条件 (a) と (b) を満たすと主張する。実際，$K \subset J_n$ は
明らかであり，$p_1$ は小拡大なので $\mathfrak m_SJ_n \subset K$ である。
これで (a) が示された。(S1) を
$$
\xymatrix{
            &  y \ar[d] \\
\xi_n \ar[r]  &  x,
}
$$
に適用すると，$\xi_n$ の $S/K \cong S/J_n \times_A B$ への持ち上げが
存在するので，(b) が成り立つ。$J_{n + 1}$ は性質 (a) と (b) をもつ
最小のイデアルなので，$J_{n + 1} \subset K$ である。したがって所要の写像
$S/J_{n+1} \to S/K \cong S/J_n \times_A B$ が存在する。
\end{proof}

\begin{remark}
\label{formal-defos-remark-compare-schlessinger-H3-pre}
$F : \mathcal{C}_\Lambda \to \textit{Sets}$ を，(S1) と (S2) を満たす
前変形関手とする。条件 $\dim_k TF < \infty$ は Schlessinger の論文の
条件 (H3) にちょうど一致する。(S1) と (S2) は条件 (H1) と (H2) に
対応することを思い出そう。
Remark \ref{formal-defos-remark-compare-S1-S2-schlessinger} を参照せよ。
したがって Lemma \ref{formal-defos-lemma-versal-object-existence} は，前変形関手に対して
$$
(H1) + (H2) + (H3)
\Rightarrow
\text{versal 形式対象が存在する}
$$
ことを述べる。Hull との関係は
Remark \ref{formal-defos-remark-compare-schlessinger-H3} で明らかにする。
\end{remark}

\section{極小 versal 形式対象}
\label{formal-defos-section-minimal-versal}

\noindent
versal 形式対象の（一意でなさ）一意性を理解するため，少し準備を行う。
前変形圏が versal 形式対象をもつならば，極小 versal 形式対象をもち，
そのような任意の二対象は同型であることが分かる。さらに，基底環を
冪級数拡大で置き換える違いを除けば，すべての versal 形式対象は
「ほぼ」同じである。

\medskip\noindent
$\mathcal{F}$ を $\mathcal{C}_\Lambda$ 上で群亜群をファイバーとする
余ファイバー圏とする。$A \in \Ob(\mathcal{C}_\Lambda)$ 上にある
$\mathcal{F}$ の各対象 $x$ に対し，圏 $\mathcal{S}_x$ を考える。
その対象は
$$
\Ob(\mathcal{S}_x) =
\{x' \to x \mid x' \to x\text{ は次の包含の上にある： }A' \subset A\}
$$
であり，射は $x$ 上の射である。$\mathcal{C}_\Lambda$ の
$f : B \to A$ 上にある $\mathcal{F}$ の各 $y \to x$ に対し，関手
$f_* : \mathcal{S}_y \to \mathcal{S}_x$ を次のように定義する。
% P10-E0461: preserve the frozen duplicated-verb construction sidecar-only.
$B' \subset B$ 上にある $y' \to y$ が与えられたとき，
$A' = f(B')$ と置き，$y'$ の順像を，$B' \to f(B')$ 上にある
$y' \to x'$ とする。群亜群をファイバーとする余ファイバー圏の公理により，
$$
\xymatrix{
y' \ar[d] \ar[r] & x' \ar[d] \\
y \ar[r] & x
}
$$
を可換にする，$f(B') \to A$ 上にある一意な射 $x' \to x$ を得る。
すると $x' \to x$ は $\mathcal{S}_x$ の対象である。
$\mathcal{S}_x$ の対象 $x' \to x$ が{\it 極小}であるとは，
$\mathcal{S}_x$ の任意の射
$(x'_1 \to x) \to (x' \to x)$ が同型であること，すなわち $x'$ と
$x'_1$ が $A$ の同じ部分環上で定義されることをいう。
$A$ は $\Lambda$-加群として有限長なので，極小対象は常に存在する。

\begin{lemma}
\label{formal-defos-lemma-smallest-where-descends}
$\mathcal{F}$ を $\mathcal{C}_\Lambda$ 上で群亜群をファイバーとする
(S1) をもつ余ファイバー圏とする。
\begin{enumerate}
\item $\mathcal{F}$ の $y \to x$ に対し，$\mathcal{S}_y$ の極小対象は
$\mathcal{S}_x$ の極小対象へ写る。
\item $\mathcal{F}$ の $y \to x$ が $\mathcal{C}_\Lambda$ の全射
$f : B \to A$ 上にあるとする。このとき $\mathcal{S}_x$ の各極小対象は，
$\mathcal{S}_y$ のある極小対象の像である。
\end{enumerate}
\end{lemma}

\begin{proof}
(1) の証明。$y \to x$ は $f : B \to A$ 上にあるとする。
$B' \subset B$ 上にある $y' \to y$ を，$\mathcal{S}_y$ の極小対象とする。
上の $f_*$ の構成に現れる図式を
$$
\vcenter{
\xymatrix{
y' \ar[d] \ar[r] & x' \ar[d] \\
y \ar[r] & x
}
}
\quad\text{次の図式の上にある}\quad
\vcenter{
\xymatrix{
B' \ar[d] \ar[r] & f(B') \ar[d] \\
B \ar[r] & A
}
}
$$
とする。$(x'' \to x) \to (x' \to x)$ が $\mathcal{S}_x$ の射で，
$x'' \to x'$ は $A'' \subset f(B')$ 上にあると仮定する。
(S1) により，$B' \times_{f(B')} A'' \to B'$ 上にある
$y'' \to y'$ が存在する。$y' \to y$ は極小なので，
$B' \times_{f(B')} A'' \to B'$ は同型である。これは
$A'' = f(B')$，すなわち $x' \to x$ が極小であることを含意する。

\medskip\noindent
(2) の証明。$f : B \to A$ は全射で，$y \to x$ は $f$ 上にあるとする。
$A' \subset A$ 上にある $\mathcal{S}_x$ の極小対象 $x' \to x$ を取る。
(S1) により，$B' = f^{-1}(A') = B \times_A A' \to B$ 上にある
$y' \to y$ が存在し，その $\mathcal{S}_x$ における像は $x' \to x$ で
ある。ゆえに $f_*(y' \to y) = x' \to x$ である。
$\mathcal{S}_y$ の射 $(y'' \to y) \to (y' \to y)$ を，
$y'' \to y$ が極小対象となるように選ぶ
（補題の上の長さに関する注意により可能である）。すると
$f_*(y'' \to y)$ は $x' \to x$ へ写る $\mathcal{S}_x$ の対象であり
（$f_*$ の関手性による），したがって $x' \to x$ の極小性により
$x' \to x$ と同型である。
\end{proof}

\begin{lemma}
\label{formal-defos-lemma-smallest-where-descends-versal}
$\mathcal{F}$ を $\mathcal{C}_\Lambda$ 上で群亜群をファイバーとする
(S1) をもつ余ファイバー圏とする。$\xi$ を $R$ 上にある
$\mathcal{F}$ の versal 形式対象とする。次の極小性をもつ，
$R' \subset R$ 上にある射 $\xi' \to \xi$ が存在する：
\begin{enumerate}
\item $A \in \Ob(\mathcal{C}_\Lambda)$ への任意の $f : R \to A$ に対し，
順像の図式
$$
\vcenter{
\xymatrix{
\xi' \ar[d] \ar[r] & x' \ar[d] \\
\xi \ar[r] & x
}
}
\quad\text{次の図式の上にある}\quad
\vcenter{
\xymatrix{
R' \ar[d] \ar[r] & f(R') \ar[d] \\
R \ar[r] & A
}
}
$$
は $\mathcal{S}_x$ の極小対象 $x' \to x$ を生む；
\item 形式対象の任意の射 $\xi'' \to \xi'$ に対し，対応する写像
$R'' \to R'$ は全射である。
\end{enumerate}
\end{lemma}

\begin{proof}
$\xi = (R, \xi_n, f_n)$ と書く。$R'_1 = k$ および
$\xi'_1 = \xi_1$ と置く。$m \leq n$ に対して
$R'_m \subset R/\mathfrak m_R^m$ 上にある
$\mathcal{S}_{\xi_m}$ の極小対象 $\xi'_m \to \xi_m$ と，
$m \leq n - 1$ に対して $f_m$ に両立する射
$f'_m : \xi'_{m + 1} \to \xi'_m$ が構成済みであると仮定する。
Lemma \ref{formal-defos-lemma-smallest-where-descends} (2) により，
$R'_{n + 1} \subset R/\mathfrak m_R^{n + 1}$ 上にある極小対象
$\xi'_{n + 1} \to \xi_{n + 1}$ で，その
$R'_n \subset R/\mathfrak m_R^n$ 上の像が $\xi'_n \to \xi_n$ となるものが
存在する。構成により可換図式
$$
\xymatrix{
\xi'_{n + 1} \ar[r]_{f'_n} \ar[d] & \xi'_n \ar[d] \\
\xi_{n + 1} \ar[r]^{f_n} & \xi_n
}
$$
を得る。さらに環写像 $R'_{n + 1} \to R'_n$ は全射である。
$R' = \lim_n R'_n$ と置く。このとき $R' \to R$ は単射である。

\medskip\noindent
しかし，$R'$ が Noether であることは先験的には明らかでない。
これを示すため，$\xi$ が versal であることを用いる。
実際，Lemma \ref{formal-defos-lemma-versal-object-quasi-initial} により，versal 性は
$\widehat{\mathcal{F}}$ の射 $\xi \to \xi'_n$ の存在を含意する。
対応する写像 $R \to R'_n$ は全射でなければならない
（$\xi'_n \to \xi_n$ が $\mathcal{S}_{\xi_n}$ で極小だからである）。
したがって余接空間の次元は有界であり，
Lemma \ref{formal-defos-lemma-limit-artinian} により $R'$ は Noether，すなわち
$\widehat{\mathcal{C}}_\Lambda$ の対象である。
Lemma \ref{formal-defos-lemma-formal-objects-different-filtration} と
Lemma \ref{formal-defos-lemma-limit-artinian} のフィルトレーションに関する結果により，
元の列 $\xi'_n$ は $R'$ 上の形式対象 $\xi'$ を定め，射
$\xi' \to \xi$ を得る。

\medskip\noindent
構成により，(1) は各 $n$ に対する $R \to R/\mathfrak m_R^n$ について
成り立つ。(1) に現れる各 $R \to A$ は
$R \to R/\mathfrak m_R^n \to A$ を通して分解するので，
% P10-E0462: preserve the frozen f(R) rather than f(R') conclusion.
$f(R) \subset A$ 上の $x' \to x$ に対する (1) は，
$R'_n \to R/\mathfrak m_R^n$ 上の $\xi'_n \to \xi_n$ の極小性から
Lemma \ref{formal-defos-lemma-smallest-where-descends} (1) により従う。

\medskip\noindent
(2) に現れる $R'' \to R'$ が全射でなければ，ある $n$ に対して
$R'' \to R' \to R'_n$ も全射ではなく，$\xi'_n \to \xi_n$ は極小では
ないことになり，矛盾する。この矛盾により (2) が示された。
\end{proof}

\begin{lemma}
\label{formal-defos-lemma-descends-versal}
$\mathcal{F}$ を $\mathcal{C}_\Lambda$ 上で群亜群をファイバーとする
(S1) をもつ余ファイバー圏とする。$\xi$ を $R$ 上にある
$\mathcal{F}$ の versal 形式対象とする。$\xi' \to \xi$ を，
Lemma \ref{formal-defos-lemma-smallest-where-descends-versal} で構成された
$R' \subset R$ 上にある形式対象の射とする。このとき
% P10-E0463: preserve the frozen r/n and X_r/x_i index mismatch.
$$
R \cong R'[[x_1, \ldots, x_r]]
$$
は $R'$ 上の冪級数環である。
さらに，$\xi'$ も versal 形式対象である。
\end{lemma}

\begin{proof}
Lemma \ref{formal-defos-lemma-versal-object-quasi-initial} により，射
$\xi \to \xi'$ が存在する。Lemma
\ref{formal-defos-lemma-smallest-where-descends-versal} により，対応する写像
$f : R \to R'$ は全射 $f|_{R'} : R' \to R'$ を誘導する。これは
Algebra, Lemma \ref{algebra-lemma-surjective-endo-noetherian-ring-is-iso}
により同型である。したがって $I = \Ker(f)$ は
$R = R' \oplus I$ を満たす $R$ のイデアルである。
$I/\mathfrak m_RI$ の基底をなす元
$x_1, \ldots, x_n \in I$ を取る。$X_i$ を $x_i$ へ写す写像
$S = R'[[X_1, \ldots, X_r]] \to R$ を考える。
各 $n \geq 1$ に対して Artinian $R'$-代数の全射
$B = S/\mathfrak m_S^n \to R/\mathfrak m_R^n = A$ を得る。
% P10-E0464: preserve the frozen missing closing parenthesis in the B-side object expression.
$y \in \Ob(\mathcal{F}(B)$，resp.\ $x \in \Ob(\mathcal{F}(A))$ を，
それぞれ $\xi'$ の $R' \to S \to B$，resp.\ $R' \to S \to A$ に沿う
順像と書く。$\xi$ は $\xi'$ の $R' \to R$ に沿う順像なので，
$x$ は $\xi$ の $R \to A$ に沿う順像でもあることに注意する。
したがって実線からなる図式
$$
\vcenter{
\xymatrix{
& y \ar[d] \\
\xi \ar[r] \ar@{..>}[ru] & x
}
}
\quad\text{次の図式の上にある}\quad
\vcenter{
\xymatrix{
& S/\mathfrak m_S^n \ar[d] \\
R \ar[r] \ar@{..>}[ru] & R/\mathfrak m_R^n
}
}
$$
を得る。$\xi$ は versal なので，
Remark \ref{formal-defos-remark-versal-object} を用いると，これらの図式に収まる
点線の矢印を得る。とくに，写像
$S/\mathfrak m_S^n \to R/\mathfrak m_R^n$ は切断
$h_n : R/\mathfrak m_R^n \to S/\mathfrak m_S^n$ をもつ。
Lemma \ref{formal-defos-lemma-power-series} から，$S \to R$ は同型である。

\medskip\noindent
$\xi$ は $\xi'$ の $R' \to R$ に沿う順像なので，
Remark \ref{formal-defos-remark-formal-objects-yoneda-map} から可換図式
$$
\xymatrix{
\underline{R}|_{\mathcal{C}_\Lambda} \ar[rr] \ar[rd]_{\underline{\xi}} & &
\underline{R'}|_{\mathcal{C}_\Lambda} \ar[ld]^{\underline{\xi'}} \\
& \mathcal{F}
}
$$
を得る。$R' \to R$ は左逆写像（すなわち $R \to R/I = R'$）をもつので，
$\underline{R}|_{\mathcal{C}_\Lambda} \to
\underline{R'}|_{\mathcal{C}_\Lambda}$ は本質的全射である。
ゆえに Lemma \ref{formal-defos-lemma-smooth-properties} により
$\underline{\xi'}$ は滑らかである，すなわち $\xi'$ は versal
形式対象である。
\end{proof}

\noindent
以上の補題を踏まえ，次の定義を与える。

\begin{definition}
\label{formal-defos-definition-minimal-versal}
$\mathcal{F}$ を前変形圏とする。$\mathcal{F}$ の versal 形式対象
$\xi$ が {\it minimal}\footnote{これは非標準的な用語かもしれない。
多くの著者はこの概念を接空間の性質と結び付ける。
この関係は Section \ref{formal-defos-section-miniversal-objects-existence} で示す。}
であるとは，形式対象の任意の射
$\xi' \to \xi$ に対して，台となる環の写像が全射であることをいう。
極小 versal 形式対象を {\it miniversal} と呼ぶこともある。
\end{definition}

\noindent
この節の議論は，この定義が妥当であることを示している。
まず，versal 形式対象の存在は $\mathcal{F}$ が (S1) をもつことを含意する。
次に，上の補題は極小 versal 形式対象が存在することを示す。
最後に，任意の二つの極小 versal 形式対象は同型である。
以下に，以上の結果を（詳しい証明とともに）まとめる。

\begin{lemma}
\label{formal-defos-lemma-minimal-versal}
$\mathcal{F}$ を versal 形式対象をもつ前変形圏とする。このとき
\begin{enumerate}
\item $\mathcal{F}$ は極小 versal 形式対象をもつ；
\item 極小 versal 形式対象は同型を除いて一意である；
\item 任意の versal 形式対象は，極小 versal 形式対象を
冪級数環拡大に沿って順送したものである。
\end{enumerate}
\end{lemma}

\begin{proof}
$\mathcal{F}$ が $R$ 上の versal 形式対象 $\xi$ をもつと仮定する。
Lemma \ref{formal-defos-lemma-versal-object-S1} により，これは (S1) を満たす。
Lemma \ref{formal-defos-lemma-smallest-where-descends-versal} で構成された射の一つを，
$R' \subset R$ 上の $\xi' \to \xi$ とする。
Lemma \ref{formal-defos-lemma-descends-versal} により $\xi'$ は versal であり，
したがって（構成により）極小 versal 形式対象である。これで (1) が示された。
また $R \cong R'[[x_1, \ldots, x_n]]$ であり，(3) も示された。

\medskip\noindent
$\xi_i/R_i$ を二つの極小 versal 形式対象とする。
Lemma \ref{formal-defos-lemma-versal-object-quasi-initial} により，射
$\xi_1 \to \xi_2$ および $\xi_2 \to \xi_1$ が存在する。
対応する環写像 $f : R_1 \to R_2$ および $g : R_2 \to R_1$ は，
極小性により全射である。したがって合成
$g \circ f : R_1 \to R_1$ および $f \circ g : R_2 \to R_2$ は，
Algebra, Lemma \ref{algebra-lemma-surjective-endo-noetherian-ring-is-iso}
により同型である。ゆえに $f$ と $g$ は同型であり，したがって射
$\xi_1 \to \xi_2$ および $\xi_2 \to \xi_1$ は同型である
（Lemma \ref{formal-defos-lemma-completion-cofibred} により
$\widehat{\mathcal{F}}$ は群亜群をファイバーとする余ファイバー圏だからである）。
これで (2) が示され，補題の証明が完了した。
\end{proof}

\section{Miniversal 形式対象と接空間}
\label{formal-defos-section-miniversal-objects-existence}

\noindent
Definition \ref{formal-defos-definition-minimal-versal} で導入した一般的な極小性の概念は，
接空間上の挙動から導けることがある。$\xi$ を前変形圏
$\mathcal{F}$ の形式対象とし，
$\underline{\xi} : \underline{R}|_{\mathcal{C}_\Lambda} \to \mathcal{F}$
を対応する射とする。このとき次の条件を考える：
% P10-E0465: preserve the frozen duplicated article sidecar-only.
\begin{equation}
\label{formal-defos-equation-bijective}
d\underline{\xi} : \text{Der}_\Lambda(R, k) \to T\mathcal{F}
\text{ は全単射である}
\end{equation}
および条件
\begin{equation}
\label{formal-defos-equation-bijective-orbits}
d\underline{\xi} : \text{Der}_\Lambda(R, k) \to T\mathcal{F}
\text{ は次の }\text{Der}_\Lambda(k, k)\text{-軌道上で全単射である。}
\end{equation}
ここでは Example \ref{formal-defos-example-tangent-space-prorepresentable-functor} の同一視
$T\underline{R}|_{\mathcal{C}_\Lambda} = \text{Der}_\Lambda(R, k)$ と，
接空間への導分の作用 (\ref{formal-defos-equation-action}) を用いている。
$k' \subset k$ が分離的ならば $\text{Der}_\Lambda(k, k) = 0$ であり，
二つの条件は同値である。条件 (S2) のもとでは，versal 形式対象が極小で
あることと，$\underline{\xi}$ が
(\ref{formal-defos-equation-bijective-orbits}) を満たすこととは同値である。
さらに，$\underline{\xi}$ が (\ref{formal-defos-equation-bijective}) を満たすならば，
$\mathcal{F}$ は (S2) を満たす。

\begin{lemma}
\label{formal-defos-lemma-miniversal-object-unique}
$\mathcal{F}$ を前変形圏とする。$\xi$ を
(\ref{formal-defos-equation-bijective-orbits}) を満たす $\mathcal{F}$ の
versal 形式対象とする。このとき $\xi$ は極小 versal 形式対象である。
とくに，そのような $\xi$ は同型を除いて一意である。
\end{lemma}

\begin{proof}
$\xi$ が極小でなければ，Lemma \ref{formal-defos-lemma-minimal-versal} により，
$R = R'[[x_1, \ldots, x_n]]$，$n > 0$ となる
$R' \to R$ 上の射 $\xi' \to \xi$ が存在する。
したがって $d\underline{\xi}$ は
$$
\text{Der}_\Lambda(R, k) \to
\text{Der}_\Lambda(R', k) \to T\mathcal{F}
$$
と分解する。ここで
$D : f \mapsto \partial/\partial x_1(f) \bmod \mathfrak m_R$ は
最初の矢印の核に属するが
$\text{Der}_\Lambda(k, k) \to \text{Der}_\Lambda(R, k)$ の像には
属さないので，(\ref{formal-defos-equation-bijective-orbits}) は成り立ちえない。
\end{proof}

\begin{lemma}
\label{formal-defos-lemma-miniversal-object-existence-1}
$\mathcal{F}$ を前変形圏とする。$\xi$ を
(\ref{formal-defos-equation-bijective}) を満たす $\mathcal{F}$ の versal 形式対象とする。
このとき
\begin{enumerate}
\item $\mathcal{F}$ は (S1) を満たす。
\item $\mathcal{F}$ は (S2) を満たす。
\item $\dim_k T\mathcal{F}$ は有限である。
\end{enumerate}
\end{lemma}

\begin{proof}
Condition (S1) は Lemma \ref{formal-defos-lemma-versal-object-S1} により成り立つ。
(S1) が成り立つので，(S2) の第一部分も成り立つ。
$$
\vcenter{
\xymatrix{
y \ar[r]_c \ar[d]_a & x_\epsilon \ar[d]^e \\
x \ar[r]^d          & x_0
}
}
\quad\text{および}\quad
\vcenter{
\xymatrix{
y' \ar[r]_{c'} \ar[d]_{a'} & x_\epsilon \ar[d]^e \\
x \ar[r]^d                 & x_0
}
}
\quad\text{次の図式の上にある}\quad
\vcenter{
\xymatrix{
A \times_k k[\epsilon] \ar[r] \ar[d] & k[\epsilon] \ar[d] \\
A  \ar[r] & k
}
}
$$
を (S2) の第二部分に現れる図式とする。上と同様に，可換図式
$$
\xymatrix{
\xi \ar[r]^{b'} \ar[d]_b          & y' \ar[d]^{a'} \\
y \ar[r]^{a} & x
}
$$
を与える射 $b : \xi \to y$ および $b' : \xi \to y'$ を見つけられる。
$p : \mathcal{F} \to \mathcal{C}_\Lambda$ を構造射とする。
$\widehat{p}(\xi) = R$，すなわち $\xi$ は
$R \in \Ob(\widehat{\mathcal{C}}_\Lambda)$ 上にあるとする。
$\xi$ の $p(c) \circ p(b)$ に沿う順像は $x_\epsilon$ であり，
$\xi$ の $p(c') \circ p(b')$ に沿う順像も $x_\epsilon$ である。
$\xi$ は (\ref{formal-defos-equation-bijective}) を満たすので，写像
$R \to k[\epsilon]$ として
$p(c) \circ p(b) = p(c') \circ p(b')$ である。したがって
$R \to A \times_k k[\epsilon]$ からの写像として $p(b) = p(b')$ である。
ゆえに $y$ と $y'$ はいずれもこの写像に沿う $\xi$ の順像と同型であり，
所望のように $b$ および $b'$ と両立する
$A \times_k k[\epsilon]$ 上の一意な射 $y \to y'$ を得る。

\medskip\noindent
最後に Example \ref{formal-defos-example-tangent-space-prorepresentable-functor} により
$\dim_k T\mathcal{F} = \dim_k T\underline{R}|_{\mathcal{C}_\Lambda}$
は有限である。
\end{proof}

\begin{example}
\label{formal-defos-example-do-not-get-S2}
(\ref{formal-defos-equation-bijective-orbits}) を満たす versal 形式対象をもつが
(S2) を満たさない前変形圏が存在する。簡単な例として，
$F = \underline{k[\epsilon]}/G$ と取る。ここで
$G \subset \text{Aut}_{\mathcal{C}_\Lambda}(k[\epsilon])$ は
非自明な有限部分群である。実際，写像
$\underline{k[\epsilon]} \to F$ は滑らかだが，$F$ の接空間は
自然な $k$-ベクトル空間構造をもたない
（有限群による $k$-ベクトル空間の商だからである）。
\end{example}

\begin{lemma}
\label{formal-defos-lemma-construct-bijective-orbits}
$\mathcal{F}$ を (S2) を満たし versal 形式対象をもつ前変形圏とする。
このとき，その極小 versal 形式対象は
(\ref{formal-defos-equation-bijective-orbits}) を満たす。
\end{lemma}

\begin{proof}
$\xi$ を $\mathcal{F}$ の極小 versal 形式対象とする
（Lemma \ref{formal-defos-lemma-minimal-versal} 参照）。
$\xi$ は $R \in \Ob(\widehat{\mathcal{C}}_\Lambda)$ 上にあるとする。
(\ref{formal-defos-equation-bijective-orbits}) を解釈するため，次のことに注意する。
$T\mathcal{F}$ は自然な $k$-ベクトル空間構造をもち
（Lemma \ref{formal-defos-lemma-tangent-space-vector-space} 参照），
$d\underline{\xi} : \text{Der}_\Lambda(R, k) \to T\mathcal{F}$ は線形であり
（Lemma \ref{formal-defos-lemma-k-linear-differential} 参照），
$\text{Der}_\Lambda(k, k)$ の作用は加法で与えられる
（Lemma \ref{formal-defos-lemma-action-linear} 参照）。図式
$$
\xymatrix{
& \Hom_k(\mathfrak m_R/\mathfrak m_R^2, k) \\
K \ar[r] & \text{Der}_\Lambda(R, k) \ar[r]^{d\underline{\xi}} \ar[u] &
T\mathcal{F} \\
& \text{Der}_\Lambda(k, k) \ar[u] \ar[ru]
}
$$
を考える。ベクトル空間 $K$ は $d\underline{\xi}$ の核である。
中央の列は列 (\ref{formal-defos-equation-sequence}) の双対なので，中央で完全である。
(\ref{formal-defos-equation-bijective-orbits}) が成り立たないならば，
$\Hom_k(\mathfrak m_R/\mathfrak m_R^2, k)$ で零へ写らない
非零元 $D \in K$ を見つけられる。
% P10-E0466: preserve the frozen wrong-article locus sidecar-only.
これは $D(t) = 1$ となる $t \in \mathfrak m_R$ が存在することを意味する。
$R' = \{a \in R \mid D(a) = 0\}$ と置く。$D$ は導分なので，
これは $R$ の部分環である。$D(t) = 1$ であるから，
$R' \to k$ は全射である
（Lemma \ref{formal-defos-lemma-essential-surjection} の証明と比較せよ）。
$\mathfrak m_{R'} = \Ker(D : \mathfrak m_R \to k)$ は $R$ のイデアルであり，
$\mathfrak m_R^2 \subset \mathfrak m_{R'}$ であることに注意する。したがって
$$
\mathfrak m_R/\mathfrak m_R^2 =
\mathfrak m_{R'}/\mathfrak m_R^2 + k\overline{t}
$$
であり，これは $\epsilon$ を $\overline{t}$ へ写す写像
$$
R'/\mathfrak m_R^2 \times_k k[\epsilon] \to R/\mathfrak m_R^2
$$
が同型であることを含意する。とくに写像
$R/\mathfrak m_R^2 \to R'/\mathfrak m_R^2$ が存在する。

\medskip\noindent
$R \to R/\mathfrak m_R^2$ 上にある射 $\xi \to y$ を取る。
$R/\mathfrak m_R^2 \to R'/\mathfrak m_R^2$ 上にある射
$y \to x$ を取る。
$R/\mathfrak m_R^2 \to k[\epsilon]$ 上にある射
$y \to x_\epsilon$ を取る。$x_0$ を $k$ 上の $\mathcal{F}$ の
（一意な同型を除いて）一意な対象とする。
群亜群をファイバーとする余ファイバー圏の公理により，可換図式
$$
\vcenter{
\xymatrix{
y \ar[r] \ar[d] & x_\epsilon \ar[d] \\
x \ar[r]   & x_0
}
}
\quad\text{次の図式の上にある}\quad
\vcenter{
\xymatrix{
R'/\mathfrak m_R^2 \times_k k[\epsilon] \ar[r] \ar[d] & k[\epsilon] \ar[d] \\
R'/\mathfrak m_R^2 \ar[r] & k.
}
}
$$
を得る。$D \in K$ なので，$x_\epsilon$ は
$0 \in \mathcal{F}(k[\epsilon])$ と同型である。すなわち
$x_\epsilon$ は $k \to k[\epsilon], a \mapsto a$ に沿う
$x_0$ の順像である。したがって Lemma \ref{formal-defos-lemma-lifting-section} により，
射 $x \to y$ が存在する。一方，
$\text{length}_\Lambda(R'/\mathfrak m_R^2) <
\text{length}_\Lambda(R/\mathfrak m_R^2)$
なので，対応する環写像
$R'/\mathfrak m_R^2 \to R/\mathfrak m_R^2$ は全射でない。
これは $\xi/R$ の極小性に矛盾する
（Lemma \ref{formal-defos-lemma-smallest-where-descends-versal} の (1) 参照）。
% P10-E0467: preserve the frozen colloquial closing sidecar-only.
この矛盾により，そのような $D$ は存在しないことが示された。
\end{proof}

\begin{theorem}
\label{formal-defos-theorem-miniversal-object-existence}
$\mathcal{F}$ を前変形圏とする。次の条件を考える。
\begin{enumerate}
\item $\mathcal{F}$ は (\ref{formal-defos-equation-bijective}) を満たす
極小 versal 形式対象をもつ；
\item $\mathcal{F}$ は (\ref{formal-defos-equation-bijective-orbits}) を満たす
極小 versal 形式対象をもつ；
\item 次の条件が成り立つ：
\begin{enumerate}
\item $\mathcal{F}$ は (S1) を満たす。
\item $\mathcal{F}$ は (S2) を満たす。
\item $\dim_k T\mathcal{F}$ は有限である。
\end{enumerate}
\end{enumerate}
常に
$$
(1) \Rightarrow (3) \Rightarrow (2).
$$
が成り立つ。$k' \subset k$ が分離的ならば，三条件はすべて同値である。
\end{theorem}

\begin{proof}
Lemma \ref{formal-defos-lemma-miniversal-object-existence-1} により
(1) $\Rightarrow$ (3) である。Lemmas
\ref{formal-defos-lemma-versal-object-existence} and
\ref{formal-defos-lemma-construct-bijective-orbits} により
(3) $\Rightarrow$ (2) である。$k' \subset k$ が分離的ならば
$\text{Der}_\Lambda(k, k) = 0$ なので，
(\ref{formal-defos-equation-bijective}) $=$ (\ref{formal-defos-equation-bijective-orbits})，
すなわち (1) と (2) は同じ条件である。

\medskip\noindent
古典的な場合の (3) $\Rightarrow$ (1) の別証明として，
Lemma \ref{formal-defos-lemma-versal-object-existence} の証明に少し言葉を補えば，
この場合には (\ref{formal-defos-equation-bijective}) を満たす versal 対象を
直ちに構成できることが分かる。
% P10-E0468: preserve the frozen self-referential lemma citation.
これは古典的な場合に Lemma \ref{formal-defos-lemma-versal-object-existence} を
用いずに済ませる。詳細は省略する。
\end{proof}

\begin{remark}
\label{formal-defos-remark-compare-schlessinger-H3}
$F : \mathcal{C}_\Lambda \to \textit{Sets}$ を，(S1)，(S2)，および
$\dim_k TF < \infty$ を満たす前変形関手とする。
これらの条件は Schlessinger の論文の条件 (H1)，(H2)，(H3) に
対応することを思い出そう
（Remark \ref{formal-defos-remark-compare-schlessinger-H3-pre} 参照）。
いま，古典的な場合（または $k' \subset k$ が分離的な場合）に，
Schlessinger に従って hull の概念を導入する。{\it hull} とは，
$d\underline{\xi} : T\underline{R}|_{\mathcal{C}_\Lambda} \to TF$ が
同型である，すなわち (\ref{formal-defos-equation-bijective}) が成り立つような
versal 形式対象 $\xi \in \widehat{F}(R)$ のことである。
したがって Theorem \ref{formal-defos-theorem-miniversal-object-existence} は，
古典的な場合に
$$
(H1) + (H2) + (H3)
\Rightarrow
\text{ hull が存在する}
$$
ことを述べている。言い換えれば，本定理は hull の存在に関する
Schlessinger の定理を復元する。
\end{remark}

\begin{remark}
\label{formal-defos-remark-compose-minimal-into-iso-classes}
$\mathcal{F}$ を前変形圏とする。
Remark \ref{formal-defos-remark-smooth-to-iso-classes} により
$\mathcal{F} \to \overline{\mathcal{F}}$ は滑らかである。
したがって，$\xi \in \widehat{\mathcal{F}}(R)$ が versal 形式対象ならば，
合成
$$
\underline{R}|_{\mathcal{C}_\Lambda} \longrightarrow
\mathcal{F} \longrightarrow \overline{\mathcal{F}}
$$
は滑らかであり（Lemma \ref{formal-defos-lemma-smooth-properties}），
$\overline{\mathcal{F}}$ における $\xi$ の像 $\overline{\xi}$ は
versal 形式対象であると結論できる。
(\ref{formal-defos-equation-bijective}) が成り立つならば，
$\overline{\xi}$ は同型
$T\underline{R}|_{\mathcal{C}_\Lambda} \to T\overline{\mathcal{F}}$
を誘導する。実際，$\mathcal{F} \to \overline{\mathcal{F}}$ は
接空間を同一視するからである。したがってこの場合，
$\overline{\xi}$ は $\overline{\mathcal{F}}$ の hull である
（Remark \ref{formal-defos-remark-compare-schlessinger-H3} 参照）。
Theorem \ref{formal-defos-theorem-miniversal-object-existence} により，
$k' \subset k$ が分離的で，$\mathcal{F}$ が (S1)，(S2)，および
$\dim_k T\mathcal{F} < \infty$ を満たす前変形圏ならば，
そのような $\xi$ を常に見つけられる。
\end{remark}

\begin{example}
\label{formal-defos-example-smooth-continued}
Lemma \ref{formal-defos-lemma-exists-smooth} では，
$\underline{R}|_{\mathcal{C}_\Lambda}$ が滑らかであり，かつ
$$
H_1(L_{k/\Lambda}) = \mathfrak m_R/\mathfrak m_R^2
\quad\text{および}\quad
\Omega_{R/\Lambda} \otimes_R k = \Omega_{k/\Lambda}
$$
を満たす $R \in \widehat{\mathcal{C}}_\Lambda$ を構成した。
上の定理を用いてこれを解釈し直そう。すなわち，
$\mathcal{F} = \mathcal{C}_\Lambda$ を恒等関手によって
それ自身の上で群亜群をファイバーとする余ファイバー圏とみなす。
このとき $\mathcal{F}$ は前変形圏であり，(S1) と (S2) を満たし，
$T\mathcal{F} = 0$ である。したがって $\mathcal{F}$ は
Theorem \ref{formal-defos-theorem-miniversal-object-existence} の条件 (3) を満たす。
定理により (2) が成り立つ。すなわち，ある
$S \in \widehat{\mathcal{C}}_\Lambda$ 上に，
(\ref{formal-defos-equation-bijective-orbits}) を満たす極小 versal 形式対象
$\xi \in \widehat{\mathcal{F}}(S)$ を見つけられる。
Lemma \ref{formal-defos-lemma-smooth} により，$\Lambda \to S$ は
$\mathfrak m_S$-進位相について形式的に滑らかである
（なぜなら
% P10-E0469: preserve the frozen underline-R source symbol.
$\underline{\xi} : \underline{R}|_{\mathcal{C}_\Lambda} \to
\mathcal{F} = \mathcal{C}_\Lambda$ は滑らかだからである）。
いま条件 (\ref{formal-defos-equation-bijective-orbits}) は，
$\text{Der}_\Lambda(S, k) \to 0$ が
$\text{Der}_\Lambda(k, k)$-軌道上で全単射であることを述べる。
これは単射 $\text{Der}_\Lambda(k, k) \to \text{Der}_\Lambda(S, k)$ が
全射でもあることを意味する。言い換えれば，
$\Omega_{S/\Lambda} \otimes_S k = \Omega_{k/\Lambda}$ である。
$\Lambda \to S$ は $\mathfrak m_S$-進位相について形式的に滑らかなので，
More on Algebra, Lemma \ref{more-algebra-lemma-formally-smooth-JZ} を
適用でき，完全列 (\ref{formal-defos-equation-sequence-extended}) は
二つの同一視
$$
H_1(L_{k/\Lambda}) = \mathfrak m_S/\mathfrak m_S^2
\quad\text{および}\quad
\Omega_{S/\Lambda} \otimes_S k = \Omega_{k/\Lambda}
$$
になる。議論を逆にたどると，
Lemma \ref{formal-defos-lemma-exists-smooth} で構成した $R$ は
極小 versal 対象を担うことが分かる。
極小 versal 対象の一意性
（Lemma \ref{formal-defos-lemma-minimal-versal}）により，
$R \cong S$，すなわち二つの構成は同じ答えを与えることも分かる。
\end{example}

\section{Rim--Schlessinger 条件と変形圏}
\label{formal-defos-section-RS-condition}

\noindent
$\mathcal{C}_\Lambda$ 上で群亜群をファイバーとする余ファイバー圏には，
実際に検証しやすく，先に導入した Schlessinger の性質 (S1) と (S2) を
含意する，非常に自然な性質がある。

\begin{definition}
\label{formal-defos-definition-RS}
$\mathcal{F}$ を $\mathcal C_\Lambda$ 上で群亜群をファイバーとする
余ファイバー圏とする。$\mathcal{F}$ が {\it 条件 (RS)} を満たすとは，
$\mathcal{F}$ の任意の図式
$$
\vcenter{
\xymatrix{
           & x_2 \ar[d] \\
x_1 \ar[r] & x
}
}
\quad\text{次の図式の上にある}\quad
\vcenter{
\xymatrix{
           & A_2 \ar[d] \\
A_1 \ar[r] & A
}
}
$$
が，$A_2 \to A$ が全射である $\mathcal{C}_\Lambda$ の図式の上にあるとき，
$\mathcal{F}$ にファイバー積 $x_1 \times_x x_2$ が存在し，図式
$$
\vcenter{
\xymatrix{
x_1 \times_x x_2 \ar[r] \ar[d] & x_2 \ar[d] \\
x_1 \ar[r]      & x
}
}
\quad\text{は次の図式の上にある}\quad
\vcenter{
\xymatrix{
A_1 \times_A A_2 \ar[r] \ar[d] & A_2 \ar[d] \\
A_1 \ar[r]      & A.
}
}
$$
の上にあることをいう。
\end{definition}

\begin{lemma}
\label{formal-defos-lemma-RS-fiber-square}
$\mathcal{F}$ を (RS) を満たす，$\mathcal{C}_\Lambda$ 上で
群亜群をファイバーとする余ファイバー圏とする。$\mathcal{F}$ の可換図式
$$
\vcenter{
\xymatrix{
y \ar[r] \ar[d] & x_2 \ar[d]   \\
x_1 \ar[r]      & x
}
}
\quad\text{次の図式の上にある}\quad
\vcenter{
\xymatrix{
A_1 \times_A A_2 \ar[r] \ar[d] & A_2 \ar[d] \\
A_1 \ar[r]      & A.
}
}
$$
が与えられ，$A_2 \to A$ が全射ならば，この図式はファイバー平方である。
\end{lemma}

\begin{proof}
$\mathcal{F}$ は (RS) を満たすので，次のファイバー積図式が存在する：
$$
\vcenter{
\xymatrix{
x_1 \times_x x_2 \ar[r] \ar[d] & x_2 \ar[d] \\
x_1 \ar[r]      & x
}
}
\quad\text{次の図式の上にある}\quad
\vcenter{
\xymatrix{
A_1 \times_A A_2 \ar[r] \ar[d] & A_2 \ar[d] \\
A_1 \ar[r]      & A.
}
}
$$
% P10-E0470: preserve the frozen A1-times-A-A1 identity.
誘導写像 $y \to x_1 \times_x x_2$ は
$\text{id} : A_1 \times_A A_1 \to A_1 \times_A A_1$ 上にあるので，
同型である。
\end{proof}

\begin{lemma}
\label{formal-defos-lemma-RS-small-extension}
$\mathcal{F}$ を $\mathcal C_\Lambda$ 上で群亜群をファイバーとする
余ファイバー圏とする。Definition \ref{formal-defos-definition-RS} の条件を
$A_2 \to A$ が小拡大の場合にだけ仮定しても，
$\mathcal{F}$ は (RS) を満たす。
\end{lemma}

\begin{proof}
Lemma \ref{formal-defos-lemma-factor-small-extension} を適用する。証明は
Lemma \ref{formal-defos-lemma-smoothness-small-extensions} の証明と同様である。
\end{proof}

\begin{lemma}
\label{formal-defos-lemma-RS-2-categorical}
$\mathcal{F}$ を $\mathcal{C}_\Lambda$ 上で群亜群をファイバーとする
余ファイバー圏とする。次は同値である。
\begin{enumerate}
\item $\mathcal{F}$ は (RS) を満たす；
\item $A_2 \to A$ が全射であるときはいつでも，関手
$\mathcal{F}(A_1 \times_A A_2) \to
\mathcal{F}(A_1) \times_{\mathcal{F}(A)} \mathcal{F}(A_2)$
は圏の同値である（(\ref{formal-defos-equation-compare}) 参照）；
% P10-E0471: preserve the frozen malformed reference placement sidecar-only.
\item $A_2 \to A$ が小拡大であるときはいつでも，(2) と同じ条件が成り立つ。
\end{enumerate}
\end{lemma}

\begin{proof}
(1) を仮定する。Lemma \ref{formal-defos-lemma-RS-fiber-square} により，
$\mathcal{F}(A_1 \times_A A_2)$ の各対象は
$x_1 \times_x x_2$ の形である。さらに
$$
\Mor_{A_1 \times_A A_2}(x_1 \times_x x_2, y_1 \times_y y_2) =
\Mor_{A_1}(x_1, y_1) \times_{\Mor_A(x, y)}
\Mor_{A_2}(x_2, y_2).
$$
である。したがって Categories, Remark
\ref{categories-remark-other-description-2-fibre-product} により，
$\mathcal{F}(A_1 \times_A A_2)$ は $\mathcal{F}(A_1)$ と
$\mathcal{F}(A_2)$ の $\mathcal{F}(A)$ 上の $2$-ファイバー積である。
言い換えれば，(2) が成り立つ。

\medskip\noindent
(2) $\Rightarrow$ (3) は直ちに従う。

\medskip\noindent
(3) を仮定する。$q_1 : A_1 \to A$ と $q_2 : A_2 \to A$ が与えられ，
$q_2$ は小拡大であるとする。Categories, Remark
\ref{categories-remark-other-description-2-fibre-product} にある
$2$-ファイバー積
$\mathcal{F}(A_1) \times_{\mathcal{F}(A)} \mathcal{F}(A_2)$ の記述を用いる。
そこで，$y \in \mathcal{F}(A_1 \times_A A_2)$ が
$(x_1, x_2, x, a_1 : x_1 \to x, a_2 : x_2 \to x)$ に対応するとする。
$z$ を $C$ 上にある $\mathcal{F}$ の対象とする。このとき
% P10-E0472: preserve the frozen ill-typed alpha compatibility equation.
\begin{align*}
\Mor_\mathcal{F}(z, y) & =
\{(f, \alpha) \mid f : C \to A_1 \times_A A_2,
\alpha : f_*z \to y\} \\
& = \{(f_1, f_2, \alpha_1, \alpha_2) \mid
f_i : C \to A_i, \ \alpha_i : f_{i, *}z \to x_i, \\
& \quad\quad\quad\quad
q_1 \circ f_1 = q_2 \circ f_2, \ q_{1, *} \alpha_1 = q_{2, *}\alpha_2\} \\
& =
\Mor_\mathcal{F}(z, x_1) \times_{\Mor_\mathcal{F}(z, x)}
\Mor_\mathcal{F}(z, x_2)
\end{align*}
であるから，$y$ は $x$ 上の $x_1$ と $x_2$ のファイバー積である。
したがって，$A_2 \to A$ が小拡大の場合に
$\mathcal{F}$ は (RS) を満たす。ゆえに
Lemma \ref{formal-defos-lemma-RS-small-extension} により (RS) が成り立つ。
\end{proof}

\begin{remark}
\label{formal-defos-remark-compare-schlessinger-H4}
$\mathcal{F}$ が集合をファイバーとする場合，条件 (RS) は
Schlessinger の論文 \cite[Theorem 2.11]{Sch} の条件 (H4) にほかならない。
実際，関手 $F: \mathcal{C}_\Lambda \to \textit{Sets}$ に対して，
条件 (RS) は次を述べる：$\mathcal{C}_\Lambda$ の写像
$A_1 \to A$ と $A_2 \to A$ があり，$A_2 \to A$ が全射ならば，
誘導写像 $F(A_1 \times_A A_2) \to F(A_1) \times_{F(A)} F(A_2)$ は
全単射である。
\end{remark}

\begin{lemma}
\label{formal-defos-lemma-RS-implies-S1-S2}
$\mathcal{F}$ を $\mathcal{C}_\Lambda$ 上で群亜群をファイバーとする
余ファイバー圏とする。$\mathcal{F}$ に対する条件 (RS) は，
$\mathcal{F}$ に対する (S1) と (S2) の両方を含意する。
\end{lemma}

\begin{proof}
Lemma \ref{formal-defos-lemma-RS-2-categorical} の言い換えと，
Definition \ref{formal-defos-definition-S1-S2} の後の (S1) の説明から，
(RS) が (S1) を含意することは直ちに分かる。
これで (S2) の第一部分が示された。(S2) の第二部分は，
Lemma \ref{formal-defos-lemma-RS-fiber-square} が，
Definition \ref{formal-defos-definition-S1-S2} の (S2) の第二部分に現れる
$y, y'$ に対して
$y = x_1 \times_{d, x_0, e} x_2 = y'$ を与えることから従う
（実際，射 $y \to y'$ は $a, a'$ と $c, c'$ の両方と両立する！）。
\end{proof}

\noindent
次の補題は Lemma \ref{formal-defos-lemma-S1-S2-associated-functor} の類似である。
$\mathcal{F}$ を $\mathcal{C}_\Lambda$ 上で群亜群をファイバーとする
余ファイバー圏とし，$x$ を $A$ 上にある $\mathcal{F}$ の対象とするとき，
$\text{Aut}_A(x) = \Mor_A(x, x) = \Mor_{\mathcal{F}(A)}(x, x)$
と書くことを思い出そう。$x' \to x$ が $A' \to A$ 上にある
$\mathcal{F}$ の射ならば，群の良定義な写像
$\text{Aut}_{A'}(x') \to \text{Aut}_A(x)$ が存在する。

\begin{lemma}
\label{formal-defos-lemma-RS-associated-functor}
$\mathcal{F}$ を (RS) を満たす，$\mathcal{C}_\Lambda$ 上で
群亜群をファイバーとする余ファイバー圏とする。次の条件は同値である：
\begin{enumerate}
\item $\overline{\mathcal{F}}$ は (RS) を満たす。
\item $f_1: A_1 \to A$ と $f_2: A_2 \to A$ を
$\mathcal{C}_\Lambda$ の環写像とし，$f_2$ は全射であるとする。
同型類の集合の誘導写像
$$
\overline{\mathcal{F}(A_1) \times_{\mathcal{F}(A)} \mathcal{F}(A_2)}
\to \overline{\mathcal{F}}(A_1) \times_{\overline{\mathcal{F}}(A)}
\overline{\mathcal{F}}(A_2)
$$
は単射である。
\item $\mathcal{F}$ の任意の射 $x' \to x$ が全射な環写像
$A' \to A$ 上にあるとき，写像
$\text{Aut}_{A'}(x') \to \text{Aut}_A(x)$ は全射である。
\item $\mathcal{F}$ の任意の射 $x' \to x$ が小拡大
$A' \to A$ 上にあるとき，写像
$\text{Aut}_{A'}(x') \to \text{Aut}_A(x)$ は全射である。
\end{enumerate}
\end{lemma}

\begin{proof}
(1) と (2) が同値であり，(2) と (3) が同値であることを示す。
(3) と (4) の同値性は Lemma \ref{formal-defos-lemma-factor-small-extension} から従う。

\medskip\noindent
$f_1: A_1 \to A$ と $f_2: A_2 \to A$ を
$\mathcal{C}_\Lambda$ の環写像とし，$f_2$ は全射であるとする。
Remark \ref{formal-defos-remark-compare-schlessinger-H4} により，
$\overline{\mathcal{F}}$ が (RS) を満たすことと，
任意のそのような $f_1, f_2$ に対して写像
$$
\overline{\mathcal{F}}(A_1 \times_A A_2) \to \overline{\mathcal
F}(A_1) \times_{\overline{\mathcal{F}}(A)} \overline{\mathcal{F}}(A_2)
$$
が全単射であることとは同値である。
% P10-E0473: preserve the frozen awkward demonstrative sidecar-only.
この写像は少なくとも全射である。実際，全射性は (S1) の条件であり，
Lemmas \ref{formal-defos-lemma-RS-implies-S1-S2} and
\ref{formal-defos-lemma-S1-S2-associated-functor} により
$\overline{\mathcal{F}}$ は (S1) を満たす。さらに，この写像は
$$
\overline{\mathcal{F}}(A_1 \times_A A_2)
\longrightarrow
\overline{\mathcal{F}(A_1) \times_{\mathcal{F}(A)} \mathcal{F}(A_2)}
\longrightarrow
\overline{\mathcal{F}}(A_1) \times_{\overline{\mathcal{F}}(A)}
\overline{\mathcal{F}}(A_2),
$$
と分解する。ここで最初の写像は，
$$
\mathcal{F}(A_1 \times_A A_2)
\longrightarrow
\mathcal{F}(A_1) \times_{\mathcal{F}(A)} \mathcal{F}(A_2)
$$
が $\mathcal{F}$ に対する (RS) により同値なので，全単射である。
したがって (1) と (2) は同値である。

\medskip\noindent
(2) が成り立つと仮定する。$x' \to x$ を，全射な環写像
$f: A' \to A$ 上にある $\mathcal{F}$ の射とする。
$a \in \text{Aut}_A(x)$ とする。
$$
(x', x', a : x \to x), \ (x', x', \text{id} : x \to x)
$$
という
$\mathcal{F}(A') \times_{\mathcal{F}(A)} \mathcal{F}(A')$ の二対象は，
$\overline{\mathcal{F}}(A') \times_{\overline{\mathcal{F}}(A)}
\overline{\mathcal{F}}(A')$ で同じ像をもつ。(2) により，写像
% P10-E0474: preserve the frozen plural-agreement defect sidecar-only.
$b_1, b_2 : x' \to x'$ が存在し，図式
$$
\xymatrix{
x \ar[r]_a \ar[d]_{f_*b_1} & x \ar[d]^{f_*b_2} \\
x \ar[r]^{\text{id}} & x
}
$$
は可換である。したがって $b_2^{-1} \circ b_1 \in \text{Aut}_{A'}(x')$
の像は $a \in \text{Aut}_A(x)$ である。ゆえに (3) が成り立つ。

\medskip\noindent
(3) が成り立つと仮定する。
$$
(x_1, x_2, a : (f_1)_*x_1 \to (f_2)_*x_2),
\ (x'_1, x'_2, a' : (f_1)_*x'_1 \to (f_2)_*x'_2)
$$
を
$\mathcal{F}(A_1) \times_{\mathcal{F}(A)} \mathcal{F}(A_2)$ の対象で，
$\overline{\mathcal{F}}(A_1) \times_{\overline{\mathcal{F}}(A)}
\overline{\mathcal{F}}(A_2)$ で同じ像をもつものとする。
すると $\mathcal{F}(A_1)$ の射 $b_1: x_1 \to x'_1$ と
$\mathcal F(A_2)$ の射 $b_2: x_2 \to x'_2$ が存在する。
(3) により，$A_2$ 上の $x_2$ の自己同型で $b_2$ を修正して，図式
$$
\xymatrix{
(f_1)_*x_1 \ar[r]_a \ar[d]_{(f_1)_*b_1} & (f_2)_*x_2 \ar[d]^{(f_2)_*b_2} \\
(f_1)_*x'_1 \ar[r]^{a'} & (f_2)_*x'_2.
}
$$
を可換にできる。これは
% P10-E0475: preserve the frozen object/isomorphism-class notation.
$(x_1, x_2, a) \cong (x'_1, x'_2, a')$ が
$\overline{\mathcal{F}(A_1) \times_{\mathcal{F}(A)} \mathcal{F}(A_2)}$
で成り立つことを示す。ゆえに (2) が成り立つ。
\end{proof}

\noindent
最後に，変形圏の概念を定義する。

\begin{definition}
\label{formal-defos-definition-deformation-category}
{\it 変形圏}とは，(RS) を満たす前変形圏 $\mathcal{F}$ のことである。
変形圏の射とは，$\mathcal{C}_\Lambda$ 上の圏の射のことである。
\end{definition}

\begin{remark}
\label{formal-defos-remark-deformation-functor}
関手 $F: \mathcal{C}_\Lambda \to \textit{Sets}$ の随伴する
集合をファイバーとする余ファイバー圏が変形圏であるとき，
これを {\it 変形関手}という。すなわち\ $F(k)$ は一点集合であり，
$F$ は (RS) を満たす。$\mathcal{F}$ が変形圏ならば，
$\overline{\mathcal{F}}$ は前変形関手だが，必ずしも変形関手ではない。
これは Lemma \ref{formal-defos-lemma-RS-associated-functor} が示すとおりである。
\end{remark}

\begin{example}
\label{formal-defos-example-prorepresentable-deformation-functor}
プロ表現可能関手 $F$ は変形関手である。実際，
$R \in \Ob(\widehat{\mathcal{C}}_\Lambda)$ であり，
$F(A) = \Mor_{\widehat{\mathcal{C}}_\Lambda}(R, A)$ と仮定する。
射 $R \to k$ は一意なので，$F(k)$ は一点集合である。また
$$
\Hom_\Lambda(R, A_1 \times_A A_2) =
\Hom_\Lambda(R, A_1) \times_{\Hom_\Lambda(R, A)}
\Hom_\Lambda(R, A_2)
$$
であるから，$\widehat{\mathcal{C}}_\Lambda$ の写像についても同じことが
成り立ち，$F$ は (RS) を満たす。
\end{example}

\noindent
次は，群亜群をファイバーとする余ファイバー圏から
$k$ 上の一点における前変形圏へ移る際の典型的な注意の一つであり，
この過程が (RS) を保つことを述べる。

\begin{lemma}
\label{formal-defos-lemma-localize-RS}
$\mathcal{F}$ を $\mathcal{C}_\Lambda$ 上で群亜群をファイバーとする
余ファイバー圏とし，$x_0 \in \Ob(\mathcal{F}(k))$ とする。
$\mathcal{F}_{x_0}$ を
Remark \ref{formal-defos-remark-localize-cofibered-groupoid} で構成した
$\mathcal{C}_\Lambda$ 上で群亜群をファイバーとする余ファイバー圏とする。
$\mathcal{F}$ が (RS) を満たすならば，$\mathcal{F}_{x_0}$ も満たす。
とくに，$\mathcal{F}_{x_0}$ は変形圏である。
\end{lemma}

\begin{proof}
$\mathcal{F}_{x_0}$ における Definition \ref{formal-defos-definition-RS} の形の
任意の図式は，$\mathcal{F}$ の図式を与える。この図式に対する
$\mathcal{F}$ での (RS) の出力は，
$\mathcal{F}_{x_0}$ での出力ともみなせる。
\end{proof}

\noindent
次の補題は，代数的スタックの $2$-ファイバー積が
代数的スタックであるという事実の類似である。

\begin{lemma}
\label{formal-defos-lemma-RS-fiber-product-morphisms}
% P10-E0476: preserve the frozen missing-article locus sidecar-only.
次を，$\mathcal{C}_\Lambda$ 上で群亜群をファイバーとする
余ファイバー圏の $2$-ファイバー積とする：
$$
\xymatrix{
\mathcal{H} \times_\mathcal{F} \mathcal{G} \ar[r] \ar[d] &
\mathcal{G} \ar[d]^g \\
\mathcal{H} \ar[r]^f & \mathcal{F}
}
$$
$\mathcal{F}, \mathcal{G}, \mathcal{H}$ がすべて (RS) を満たすならば，
$\mathcal{H} \times_\mathcal{F} \mathcal{G}$ は (RS) を満たす。
\end{lemma}

\begin{proof}
$A$ を $\mathcal{C}_\Lambda$ の対象とする。このとき，
$\mathcal{H} \times_\mathcal{F} \mathcal{G}$ の $A$ 上のファイバー圏の
対象は三つ組 $(u, v, a)$ である。ここで
$u \in \mathcal{H}(A)$，$v \in \mathcal{G}(A)$ であり，
$a : f(u) \to g(v)$ は $\mathcal{F}(A)$ の射である。
$\mathcal{H} \times_\mathcal{F} \mathcal{G}$ の図式
$$
\vcenter{
\xymatrix{
           & (u_2, v_2, a_2) \ar[d] \\
(u_1, v_1, a_1) \ar[r] & (u, v, a)
}
}
\quad\text{次の図式の上にある}\quad
\vcenter{
\xymatrix{
           & A_2 \ar[d] \\
A_1 \ar[r] & A
}
}
$$
を考え，$A_2 \to A$ は $\mathcal{C}_\Lambda$ の全射であるとする。
$\mathcal{H}$ と $\mathcal{G}$ は (RS) を満たすので，
$A_1 \times_A A_2$ 上にあるファイバー積
$u_1 \times_u u_2$ および $v_1 \times_v v_2$ が存在する。
$\mathcal{F}$ も (RS) を満たすので，
Lemma \ref{formal-defos-lemma-RS-fiber-square} により
$$
\vcenter{
\xymatrix{
f(u_1 \times_u u_2) \ar[r] \ar[d] & f(u_2) \ar[d] \\
f(u_1) \ar[r] & f(u)
}
}
\quad\text{および}\quad
\vcenter{
\xymatrix{
g(v_1 \times_v v_2) \ar[r] \ar[d] & g(v_2) \ar[d] \\
g(v_1) \ar[r] & g(v)
}
}
$$
はいずれも $\mathcal{F}$ のファイバー平方である。したがって，
$a_1 \times_a a_2$ を，$A_1 \times_A A_2$ 上の
$f(u_1 \times_u u_2)$ から $g(v_1 \times_v v_2)$ への射とみなせる。
ゆえに
$$
\xymatrix{
 (u_1 \times_u u_2, v_1 \times_v v_2, a_{1} \times_a a_2) \ar[d] \ar[r] &
(u_2, v_2, a_2) \ar[d] \\
(u_1, v_1, a_1) \ar[r] & (u, v, a)
}
$$
は，所望のように
$\mathcal{H} \times_\mathcal{F} \mathcal{G}$ のファイバー平方である。
\end{proof}

\section{対象の持ち上げ}
\label{formal-defos-section-lifts}

\noindent
この節の内容は，変形圏の場合に，接空間が小拡大に沿う
持ち上げの集合へ自由かつ推移的に作用するということである。

\begin{definition}
\label{formal-defos-definition-lifts}
$\mathcal{F}$ を $\mathcal{C}_\Lambda$ 上で群亜群をファイバーとする
余ファイバー圏とする。$f: A' \to A$ を $\mathcal{C}_\Lambda$ の写像とし，
$x \in \mathcal{F}(A)$ とする。$f$ に沿う $x$ の持ち上げの圏
$\textit{Lift}(x, f)$ を，次の対象と射をもつ圏として定める。
\begin{enumerate}
\item 対象：{\it $f$ に沿う $x$ の持ち上げ}とは，
$f$ 上にある射 $x' \to x$ のことである。
\item 射：$a_1 : x'_1 \to x$ から $a_2 : x'_2 \to x$ への
{\it 持ち上げの射}とは，$\mathcal{F}(A')$ の射
$b : x'_1 \to x'_2$ であって，
% P10-E0477: preserve the frozen ill-typed lift-morphism equality.
$a_2 = a_1 \circ b$ を満たすものをいう。
\end{enumerate}
$f$ に沿う $x$ の持ち上げの集合 $\text{Lift}(x, f)$ とは，
$\textit{Lift}(x, f)$ の同型類の集合である。
\end{definition}

\begin{remark}
\label{formal-defos-remark-omit-arrow}
写像 $f: A' \to A$ が文脈から明らかなときは，
$\textit{Lift}(x, f)$ および $\text{Lift}(x, f)$ の代わりに
$\textit{Lift}(x, A')$ および $\text{Lift}(x, A')$ と書くことがある。
\end{remark}

\begin{remark}
\label{formal-defos-remark-tangent-space-lifting}
$\mathcal{F}$ を $\mathcal{C}_\Lambda$ 上で群亜群をファイバーとする
余ファイバー圏とし，$x_0 \in \Ob(\mathcal{F}(k))$ とする。
$V$ を有限次元ベクトル空間とする。このとき
$\text{Lift}(x_0, k[V])$ は $\mathcal{F}_{x_0}(k[V])$ の
同型類の集合である。ここで $\mathcal{F}_{x_0}$ は
$x_0$ 上にある $\mathcal{F}$ の対象の前変形圏である
（Remark \ref{formal-defos-remark-localize-cofibered-groupoid} 参照）。
したがって $\mathcal{F}$ が (S2) を満たすならば，
$\mathcal{F}_{x_0}$ も満たし
（Lemma \ref{formal-defos-lemma-S1-S2-localize} 参照），
Lemma \ref{formal-defos-lemma-tangent-space-vector-space} により
$$
\text{Lift}(x_0, k[V]) = T\mathcal{F}_{x_0} \otimes_k V
$$
が $k$-ベクトル空間として成り立つ。
\end{remark}

\begin{remark}
\label{formal-defos-remark-lift-bijections}
$\mathcal{F}$ を (RS) を満たす，$\mathcal C_\Lambda$ 上で
群亜群をファイバーとする余ファイバー圏とする。
$$
\xymatrix{
A_1 \times_A A_2 \ar[r] \ar[d] & A_2 \ar[d] \\
A_1 \ar[r] & A
}
$$
を $\mathcal{C}_\Lambda$ のファイバー平方とし，
$A_1 \to A$ または $A_2 \to A$ のいずれかが全射であるとする。
$x \in \Ob(\mathcal{F}(A))$ とする。$x$ の $A_1$ および $A_2$ への
持ち上げ $x_1 \to x$ および $x_2 \to x$ が与えられると，
(RS) により，$x$ の $A_1 \times_A A_2$ への持ち上げ
$x_1 \times_x x_2 \to x$ を得る。逆に
Lemma \ref{formal-defos-lemma-RS-fiber-square} により，
$x$ の $A_1 \times_A A_2$ への任意の持ち上げはこの形である。
% P10-E0478: preserve the frozen missing finite verb sidecar-only.
したがって全単射
$$
\text{Lift}(x, A_1) \times \text{Lift}(x, A_2)
\longrightarrow
\text{Lift}(x, A_1 \times_A A_2).
$$
を得る。同様に，$x_1 \to x$ が $x$ の $A_1$ への固定した持ち上げならば，
全単射
$$
\text{Lift}(x_1, A_1 \times_A A_2)
\longrightarrow
\text{Lift}(x, A_2).
$$
が存在する。いま
$$
\xymatrix{
A_1' \times_A A_2 \ar[r] \ar[d] & A_1 \times_A A_2 \ar[r] \ar[d] & A_2
\ar[d] \\
A_1' \ar[r] & A_1 \ar[r] & A
}
$$
を $\mathcal{C}_\Lambda$ のファイバー平方の合成とし，
$A'_1 \to A_1$ と $A_1 \to A$ はともに全射であるとする。
$x_1 \to x$ を $A_1 \to A$ 上にある射とする。このとき上の議論により
\begin{align*}
\text{Lift}(x_1, A_1' \times_A A_2)
& = \text{Lift}(x_1, A_1') \times \text{Lift}(x_1, A_1 \times_A A_2) \\
& = \text{Lift}(x_1, A_1') \times \text{Lift}(x, A_2).
\end{align*}
という全単射を得る。
\end{remark}

\begin{lemma}
\label{formal-defos-lemma-free-transitive-action}
$\mathcal{F}$ を変形圏とする。$A' \to A$ を
$\mathcal{C}_\Lambda$ の全射な環写像とし，その核 $I$ は
$\mathfrak m_{A'}$ で消されるとする。$x \in \Ob(\mathcal{F}(A))$ とする。
$\text{Lift}(x, A')$ が空でないならば，
$T\mathcal{F} \otimes_k I$ は $\text{Lift}(x, A')$ に
自由かつ推移的に作用する。
\end{lemma}

\begin{proof}
規則 $g(a_1, a_2) = \overline{a_1} \oplus a_2 - a_1$ で定まる環写像
$g : A' \times_A A' \to k[I]$ を考える
（Lemma \ref{formal-defos-lemma-lifting-along-small-extension} と比較せよ）。
$(a_1, a_2) \mapsto (a_1, g(a_1, a_2))$ で与えられる同型
$$
A' \times_A A' \xrightarrow{\sim} A' \times_k k[I]
$$
が存在する。この同型は第一成分の $A'$ への射影と可換であり，
したがって $A' \times_A A'$ と $A' \times_k k[I]$ から
$A$ への射影とも可換である。ゆえに全単射
\begin{equation}
\label{formal-defos-equation-one}
\text{Lift}(x, A' \times_A A')
\longrightarrow
\text{Lift}(x, A' \times_k k[I])
\end{equation}
が存在する。Remark \ref{formal-defos-remark-lift-bijections} により全単射
\begin{equation}
\label{formal-defos-equation-two}
\text{Lift}(x, A') \times \text{Lift}(x, A')
\longrightarrow
\text{Lift}(x, A' \times_A A')
\end{equation}
が存在する。また可換図式
$$
\xymatrix{
A' \times_k k[I] \ar[r] \ar[d] & A \times_k k[I] \ar[r] \ar[d] & k[I] \ar[d] \\
A' \ar[r] & A \ar[r] & k.
}
$$
がある。したがって $x$ の $A \to k$ に沿う順像
$x \to x_0$ を一つ選べば，Remark \ref{formal-defos-remark-lift-bijections} の
末尾により全単射
\begin{equation}
\label{formal-defos-equation-three}
\text{Lift}(x, A' \times_k k[I])
\longrightarrow
\text{Lift}(x, A') \times \text{Lift}(x_0, k[I])
\end{equation}
を得る。(\ref{formal-defos-equation-two})，(\ref{formal-defos-equation-one})，および
(\ref{formal-defos-equation-three}) を合成して，全単射
$$
\Phi :
\text{Lift}(x, A') \times \text{Lift}(x, A')
\longrightarrow
\text{Lift}(x, A') \times \text{Lift}(x_0, k[I]).
$$
を得る。この全単射は第一成分への射影と可換である。
Remark \ref{formal-defos-remark-tangent-space-lifting} により
$\text{Lift}(x_0, k[I]) = T\mathcal{F} \otimes_k I$ である。
$\text{pr}_2$ を $\text{Lift}(x, A') \times \text{Lift}(x, A')$ の
第二射影とすれば，写像
$$
a = \text{pr}_2 \circ \Phi^{-1} :
\text{Lift}(x, A') \times (T\mathcal{F} \otimes_k I)
\longrightarrow
\text{Lift}(x, A').
$$
を得る。以上を展開すると，$a(x' \to x, \theta)$ は
$g_*(x', x'') = \theta$ を満たす一意な持ち上げ $x'' \to x$ である。
ここで等式は
$\text{Lift}(x_0, k[I]) = T\mathcal{F} \otimes_k I$ におけるものである。
これが $\text{Lift}(x, A')$ への $T\mathcal{F} \otimes_k I$ の
作用であることを示すには，次を示せばよい：
$x', x'', x'''$ を $x$ の持ち上げとし，
$g_*(x', x'') = \theta$，$g_*(x'', x''') = \theta'$ ならば，
$g_*(x', x''') = \theta + \theta'$ である。これは可換図式
$$
\xymatrix{
A' \times_A A' \times_A A'
\ar[rrrrr]_-{(a_1, a_2, a_3) \mapsto (g(a_1, a_2), g(a_2, a_3))}
\ar[rrrrrd]_{(a_1, a_2, a_3) \mapsto g(a_1, a_3)} & & & & &
k[I] \times_k k[I] = k[I \times I] \ar[d]^{+} \\
& & & & & k[I]
}
$$
から従う。$\Phi$ は全単射なので，作用は自由かつ推移的である。
\end{proof}

\begin{remark}
\label{formal-defos-remark-free-transitive-action-functorial}
Lemma \ref{formal-defos-lemma-free-transitive-action} の作用は関手的である。
$\varphi : \mathcal{F} \to \mathcal{G}$ を変形圏の射とする。
$A' \to A$ を全射な環写像とし，その核 $I$ は
$\mathfrak m_{A'}$ で消されるとする。
$x \in \Ob(\mathcal{F}(A))$ とする。この状況で $\varphi$ は，
次の可換図式の縦の矢印を誘導する：
$$
\xymatrix{
\text{Lift}(x, A') \times (T\mathcal{F} \otimes_k I)
\ar[d]_{(\varphi, d\varphi \otimes \text{id}_I)} \ar[r] &
\text{Lift}(x, A') \ar[d]^\varphi \\
\text{Lift}(\varphi(x), A') \times (T\mathcal{G} \otimes_k I) \ar[r] &
\text{Lift}(\varphi(x), A')
}
$$
可換性は，各写像 (\ref{formal-defos-equation-two})，(\ref{formal-defos-equation-one})，および
(\ref{formal-defos-equation-three})（Lemma \ref{formal-defos-lemma-free-transitive-action} の
証明に現れる）が同様の可換図式を生むことから従う。
\end{remark}

\section{プロ表現可能関手に関する Schlessinger の定理}
\label{formal-defos-section-schlessingers-theorem}

\noindent
$\mathcal{C}_\Lambda$ 上のプロ表現可能関手を特徴付ける
Schlessinger の定理を導く。

\begin{lemma}
\label{formal-defos-lemma-minimal-smooth-morphism-functors}
$F, G: \mathcal{C}_\Lambda \to \textit{Sets}$ を変形関手とする。
$\varphi : F \to G$ を滑らかな射とし，接空間の同型
$d\varphi : TF \to TG$ を誘導するとする。このとき $\varphi$ は同型である。
\end{lemma}

\begin{proof}
すべての $A \in \Ob(\mathcal{C}_\Lambda)$ に対し
$F(A) \to G(A)$ が全単射であることを，
$\text{length}_A(A)$ に関する帰納法で示す。$A = k$ のときは，
$F$ と $G$ が変形関手であるという仮定から従う。
長さが $n$ 未満の環について主張が成り立つとし，
$A'$ を長さ $n$ の環とする。小拡大 $f : A' \to A$ を選ぶ。
可換図式
$$
\xymatrix{
F(A') \ar[r] \ar[d]_{F(f)} & G(A') \ar[d]^{G(f)} \\
F(A) \ar[r]^{\sim} & G(A)
}
$$
があり，写像 $F(A) \to G(A)$ は全単射である。
$F \to G$ の滑らかさにより，$F(A') \to G(A')$ は全射である
（Lemma \ref{formal-defos-lemma-smooth-morphism-essentially-surjective}）。
したがって，$F(f)^{-1}(x)$ が空でない $x \in F(A)$ に対するファイバー
$F(f)^{-1}(x) \to G(f)^{-1}(\varphi(x))$ 上で全単射性を
確認すればよい。これらのファイバーはそれぞれ
$\text{Lift}(x, A')$ と $\text{Lift}(\varphi(x), A')$ であり，
仮定により同型
$d\varphi \otimes \text{id} :
TF \otimes_k \Ker(f)  \to TG \otimes_k \Ker(f)$
をもつ。したがって Lemma \ref{formal-defos-lemma-free-transitive-action} と
Remark \ref{formal-defos-remark-free-transitive-action-functorial} により，
$F(f)^{-1}(x)$ が空でない $x \in F(A)$ に対する写像
$F(f)^{-1}(x) \to G(f)^{-1}(\varphi(x))$ は，
$TF \otimes_k \Ker(f)$ の自由かつ推移的な作用と可換する集合の写像である。
ゆえに全単射である。
\end{proof}

\noindent
$k' \subset k$ が分離的な場合，次の定理の条件 (c) は空であることに注意する。

\begin{theorem}
\label{formal-defos-theorem-Schlessinger-prorepresentability}
$F: \mathcal{C}_\Lambda \to \textit{Sets}$ を関手とする。このとき，
$F$ がプロ表現可能であるための必要十分条件は，
(a) $F$ が変形関手であり，(b) $\dim_k TF$ が有限であり，かつ
(c) $\gamma : \text{Der}_\Lambda(k, k) \to TF$ が単射であることである。
\end{theorem}

\begin{proof}
$F$ が $R \in \widehat{\mathcal{C}}_\Lambda$ によりプロ表現されると仮定する。
Example \ref{formal-defos-example-prorepresentable-deformation-functor} により，
$F$ は変形関手である。
Example \ref{formal-defos-example-tangent-space-prorepresentable-functor} により，
$\dim_k TF$ は有限である。最後に，
$\text{Der}_\Lambda(k, k) \to TF$ は
Example \ref{formal-defos-example-tangent-space-map-prorepresentable-functor} により
$\text{Der}_\Lambda(k, k) \to \text{Der}_\Lambda(R, k)$ と同一視され，
$R \to k$ が全射なので単射である。

\medskip\noindent
逆に，(a)，(b)，(c) が成り立つと仮定する。
Lemma \ref{formal-defos-lemma-RS-implies-S1-S2} により (S1) と (S2) が成り立つ。
したがって Theorem \ref{formal-defos-theorem-miniversal-object-existence} により，
(\ref{formal-defos-equation-bijective-orbits}) を満たす $F$ の極小 versal 形式対象
$\xi$ が存在する。$\xi$ は $R$ 上にあるとする。写像
% P10-E0479: preserve the frozen T-mathcal-F codomain.
$$
d\underline{\xi} : \text{Der}_\Lambda(R, k) \to T\mathcal{F}
$$
は $\text{Der}_\Lambda(k, k)$-軌道上で全単射である。
左辺への $\text{Der}_\Lambda(k, k)$ の作用は，(c) と
Lemma \ref{formal-defos-lemma-action-linear} により自由なので，この写像は全単射である。
したがって Lemma \ref{formal-defos-lemma-minimal-smooth-morphism-functors} により，
$\underline{\xi}$ は同型である。
\end{proof}

\section{無限小自己同型}
\label{formal-defos-section-infinitesimal-automorphisms}

\noindent
$\mathcal{F}$ を $\mathcal{C}_\Lambda$ 上で群亜群をファイバーとする
余ファイバー圏とする。$A' \to A$ 上にある $\mathcal{F}$ の射
$x' \to x$ が与えられると，準同型
$$
\text{Aut}_{A'}(x') \to \text{Aut}_A(x).
$$
が誘導される。Lemma \ref{formal-defos-lemma-RS-associated-functor} は，
この準同型の余核が，$\mathcal{F}$ の条件 (RS) が
$\overline{\mathcal{F}}$ へ移るかどうかを決定することを述べる。
この節では，この準同型の核を調べる。それもまた，
$\mathcal{F}$ と $\overline{\mathcal{F}}$ の隔たりを測ることが分かる。

\begin{definition}
\label{formal-defos-definition-relative-infinitesimal-auts}
$\mathcal{F}$ を $\mathcal C_\Lambda$ 上で群亜群をファイバーとする
余ファイバー圏とする。$x' \to x$ を $A' \to A$ 上にある
$\mathcal{F}$ の射とする。核
$$
\text{Inf}(x'/x) = \Ker(\text{Aut}_{A'}(x') \to \text{Aut}_A(x))
$$
を {\it $x$ 上の $x'$ の無限小自己同型群}という。
\end{definition}

\begin{definition}
\label{formal-defos-definition-infinitesimal-auts}
$\mathcal{F}$ を $\mathcal C_\Lambda$ 上で群亜群をファイバーとする
余ファイバー圏とする。$x_0 \in \Ob(\mathcal{F}(k))$ とする。
写像 $k \to k[\epsilon], a \mapsto a$ に沿う $x_0$ の順像
$x_0 \to x_0'$ を一つ選んだと仮定する。このとき，
$x_0 \to x_0' \to x_0$ が $x_0$ 上の恒等射となるような
一意な写像 $x'_0 \to x_0$ が存在する。このとき
$$
\text{Inf}_{x_0}(\mathcal F) = \text{Inf}(x'_0/x_0)
$$
は {\it $x_0$ の無限小自己同型群}である。
% P10-E0480: preserve the frozen missing terminal punctuation sidecar-only.
\end{definition}

\begin{remark}
\label{formal-defos-remark-choice-pushforward-immaterial-infinitesimal-aut}
標準同型を除けば，$\text{Inf}_{x_0}(\mathcal{F})$ は
順像 $x_0 \to x_0'$ の選択に依存しない。実際，任意の二つの順像は
標準的に同型である。さらに，$y_0 \in \mathcal{F}(k)$ かつ
$\mathcal{F}(k)$ において $x_0 \cong y_0$ ならば，
$\text{Inf}_{x_0}(\mathcal{F}) \cong \text{Inf}_{y_0}(\mathcal{F})$
であり，この同型は同型 $x_0 \to y_0$ の選択（だけ）に依存する。
とくに，$\text{Aut}_k(x_0)$ は $\text{Inf}_{x_0}(\mathcal{F})$ に作用する。
\end{remark}

\begin{remark}
\label{formal-defos-remark-trivial-aut-point}
$\mathcal{F}$ を前変形圏と仮定する。このとき
\begin{enumerate}
\item $x_0 \in \Ob(\mathcal{F}(k))$ に対して自己同型群
$\text{Aut}_k(x_0)$ は自明であり，したがって
$\text{Inf}_{x_0}(\mathcal{F}) = \text{Aut}_{k[\epsilon]}(x'_0)$ である；
\item $x_0, y_0 \in \Ob(\mathcal{F}(k))$ に対して
一意な同型 $x_0 \to y_0$ が存在し，したがって標準的な同一視
$\text{Inf}_{x_0}(\mathcal{F}) = \text{Inf}_{y_0}(\mathcal{F})$ が存在する。
\end{enumerate}
$\mathcal{F}(k)$ は空でないので，$x_0 \in \Ob(\mathcal{F}(k))$ を選び
$$
\text{Inf}(\mathcal{F}) = \text{Inf}_{x_0}(\mathcal{F})
$$
と置くことで，良定義な {\it $\mathcal{F}$ の無限小自己同型群}を得る。
この記法では
$\text{Inf}(\mathcal{F}_{x_0}) = \text{Inf}_{x_0}(\mathcal{F})$ である。
Remark \ref{formal-defos-remark-tangent-space-cofibered-groupoid} の等式
$T\mathcal{F}_{x_0} = T_{x_0}\mathcal{F}$ と比較されたい。
\end{remark}

\noindent
$\mathcal{F}$ が (RS) を満たすとき，
$\text{Inf}_{x_0}(\mathcal{F})$ は自然な $k$-ベクトル空間構造を
もつことが分かる。同時に，$\mathcal{F}$ が (RS) を満たすならば，
小拡大上にある射 $x' \to x$ の無限小自己同型
$\text{Inf}(x'/x)$ は $\text{Inf}_{x_0}(\mathcal{F})$ によって
支配されることも分かる。ここで $x_0$ は $x$ の
$\mathcal{F}(k)$ への順像である。このため，任意の対象
$x \in \Ob(\mathcal{F})$ の自己同型関手を次のように導入する。

\begin{definition}
\label{formal-defos-definition-automorphism-functor}
$p : \mathcal{F} \to \mathcal{C}$ を，任意の基底圏 $\mathcal{C}$ 上で
群亜群をファイバーとする余ファイバー圏とする。順像を選んだと仮定する。
$x \in \Ob(\mathcal{F})$ とし，$U = p(x)$ と置く。
$U/\mathcal{C}$ を $U$ の下の対象の圏とする。
{\it $x$ の自己同型関手}とは，対象 $f : U \to V$ を
$\text{Aut}_V(f_*x)$ へ写し，射
$$
\xymatrix{
V' \ar[rr] &                    & V\\
          & U \ar[ul]^{f'}  \ar[ur]_f &
}
$$
を準同型
$\text{Aut}_{V'}(f'_*x) \to \text{Aut}_V(f_*x)$
へ写す関手
$\mathit{Aut}(x) : U/\mathcal{C} \to \textit{Sets}$ のことである。
この準同型は，$V' \to V$ 上にあり，$x \to f'_*x$ および
$x \to f_*x$ と両立する一意な射 $f'_*x \to f_*x$ から生じる。
\end{definition}

\noindent
以下では，$\mathcal{C}_\Lambda$ 上で群亜群をファイバーとする
余ファイバー圏 $\mathcal{F}$ の対象の自己同型関手を扱う。
$A \in \Ob(\mathcal{C}_\Lambda)$ ならば，圏
$A/\mathcal{C}_\Lambda$ は圏 $\mathcal{C}_A$ にほかならない。
すなわち\ Section \ref{formal-defos-section-CLambda} で
$\Lambda = A$ および $k = A/\mathfrak m_A$ として定義した圏である。
したがって，対象 $x \in \Ob(\mathcal{F}(A))$ の自己同型関手は関手
$\mathit{Aut}(x) : \mathcal{C}_A \to \textit{Sets}$ である。

\medskip\noindent
次の補題は，群亜群をファイバーとする余ファイバー圏の「惰性」を
考えることで Lemma \ref{formal-defos-lemma-RS-fiber-product-morphisms} から
導くこともできる。例えば
Stacks, Section \ref{stacks-section-the-inertia-stack} および
Categories, Section \ref{categories-section-inertia} を参照せよ。
しかし，直接見る方が容易である。

\begin{lemma}
\label{formal-defos-lemma-Aut-functor-RS}
$\mathcal{F}$ を (RS) を満たす，$\mathcal{C}_\Lambda$ 上で
群亜群をファイバーとする余ファイバー圏とする。
$x \in \Ob(\mathcal{F}(A))$ とする。このとき
$\mathit{Aut}(x): \mathcal{C}_A \to \textit{Sets}$ は (RS) を満たす。
\end{lemma}

\begin{proof}
% P10-E0481: preserve the frozen malformed full-faithfulness phrase sidecar-only.
$\mathit{Aut}(x)$ が (RS) を満たすことは，
Lemma \ref{formal-defos-lemma-RS-2-categorical} の関手
$\mathcal{F}(A_1 \times_A A_2) \to
\mathcal{F}(A_1) \times_{\mathcal{F}(A)} \mathcal{F}(A_2)$ の
充満忠実性から従う。
\end{proof}

\begin{lemma}
\label{formal-defos-lemma-Aut-functor-tangent-space}
$\mathcal{F}$ を (RS) を満たす，$\mathcal{C}_\Lambda$ 上で
群亜群をファイバーとする余ファイバー圏とする。
$x \in \Ob(\mathcal{F}(A))$ とする。$x_0$ を $x$ の
$\mathcal{F}(k)$ への順像とする。
\begin{enumerate}
\item $T_{\text{id}_{x_0}} \mathit{Aut}(x)$ は，和が
$T_{\text{id}_{x_0}} \mathit{Aut}(x)$ における合成と一致するような
自然な $k$-ベクトル空間構造をもつ。とくに，
$T_{\text{id}_{x_0}} \mathit{Aut}(x)$ における合成は可換である。
\item $k$-ベクトル空間の標準同型
$T_{\text{id}_{x_0}} \mathit{Aut}(x) \to
T_{\text{id}_{x_0}} \mathit{Aut}(x_0)$
が存在する。
\end{enumerate}
\end{lemma}

\begin{proof}
Remark \ref{formal-defos-remark-localize-cofibered-groupoid} を，関手
$\mathit{Aut}(x) : \mathcal{C}_A \to \textit{Sets}$ と元
$\text{id}_{x_0} \in \mathit{Aut}(x)(k)$ に適用すると，前変形関手
$F = \mathit{Aut}(x)_{\text{id}_{x_0}}$ を得る。
Lemmas \ref{formal-defos-lemma-Aut-functor-RS} and \ref{formal-defos-lemma-localize-RS} により
$F$ は変形関手である。定義により
$T_{\text{id}_{x_0}} \mathit{Aut}(x) = TF = F(k[\epsilon])$
であり，これは Lemma \ref{formal-defos-lemma-tangent-space-functor} で指定される
自然な $k$-ベクトル空間構造をもつ。

\medskip\noindent
和は合成
$$
F(k[\epsilon]) \times F(k[\epsilon]) \longrightarrow
F(k[\epsilon] \times_k k[\epsilon]) \longrightarrow
F(k[\epsilon])
$$
として定義される。ここで第一の写像は (RS) が保証する全単射の逆であり，
第二の写像は，$(\epsilon, 0)$ と $(0, \epsilon)$ を $\epsilon$ に写す
$k$-代数写像
$k[\epsilon] \times_k k[\epsilon] \to k[\epsilon]$
によって誘導される。$A \to B$ が $\mathcal{C}_\Lambda$ の環写像ならば，
$F(A) \to F(B)$ は準同型である。ここで
$F(A) = \mathit{Aut}(x)_{\text{id}_{x_0}}(A)$ と
$F(B) = \mathit{Aut}(x)_{\text{id}_{x_0}}(B)$ は，合成を演算とする群である。
したがって，$F(k[\epsilon])$ を合成に関する群とみなすと，
$+ : F(k[\epsilon]) \times F(k[\epsilon])\to F(k[\epsilon])$
は準同型である。$\text{id} \in F(k[\epsilon])$ を単位元とすると，
$+(v, \text{id}) =
+(\text{id}, v) = v$ が任意の $v \in F(k[\epsilon])$ に対して成り立つ。
実際，$(\text{id}, v)$ は，
$\epsilon \mapsto (\epsilon, 0)$ で与えられる環写像
$k[\epsilon] \to k[\epsilon] \times_k k[\epsilon]$ に沿う $v$ の順像である。
% P10-E0482: preserve the frozen malformed Eckmann--Hilton sentence sidecar-only.
一般に，乗法 $\circ$ をもつ群 $G$ と，
$+(g, 1) = +(1, g) = g$ を満たす準同型 $+ : G \times G \to G$ が
与えられたとする。ここで $1$ は $G$ の単位元である。このとき
$+ = \circ$ である。これにより，$F(k[\epsilon])$ の
$k$-ベクトル空間構造における和は合成と一致する。

\medskip\noindent
最後に，(2) は定義を展開すれば分かる。すなわち，
$T_{\text{id}_{x_0}} \mathit{Aut}(x)$ は，$A \to k \to k[\epsilon]$ に
沿う $x$ の順像の自己同型 $\alpha$ で，$\epsilon$ を法として自明なものの
集合である。一方，$T_{\text{id}_{x_0}} \mathit{Aut}(x_0)$ は，
$k \to k[\epsilon]$ に沿う $x_0$ の順像の自己同型で，
$\epsilon$ を法として自明なものの集合である。$x_0$ は
$A \to k$ に沿う $x$ の順像なので，結果は明らかである。
\end{proof}

\begin{remark}
\label{formal-defos-remark-infaut-lifting-equalities}
無限小自己同型群，持ち上げ，および自己同型関手の接空間の間の
基本的な関係をいくつか指摘する。$\mathcal{F}$ を
$\mathcal{C}_\Lambda$ 上で群亜群をファイバーとする余ファイバー圏とする。
$x' \to x$ を環写像 $A' \to A$ 上にある射とする。
このとき定義から等式
$$
\text{Inf}(x'/x) = \text{Lift}(\text{id}_x, A')
$$
を得る。ここで持ち上げとは，$\mathit{Aut}(x')$ の対象としての
$\text{id}_x$ の持ち上げである。$x_0 \in \Ob(\mathcal{F}(k))$ とし，
$x'_0$ を $\mathcal{F}(k[\epsilon])$ への順像とすると，
これを $x'_0 \to x_0$ に適用して
$$
\text{Inf}_{x_0}(\mathcal{F}) =
\text{Lift}(\text{id}_{x_0}, k[\epsilon]) =
T_{\text{id}_{x_0}} \mathit{Aut}(x_0),
$$
を得る。最後の等式は定義から直接従う。
\end{remark}

\begin{lemma}
\label{formal-defos-lemma-infaut-vector-space}
$\mathcal{F}$ を (RS) を満たす，$\mathcal{C}_\Lambda$ 上で
群亜群をファイバーとする余ファイバー圏とする。
$x_0 \in \Ob(\mathcal{F}(k))$ とする。このとき
$\text{Inf}_{x_0}(\mathcal{F})$ は集合として
$T_{\text{id}_{x_0}} \mathit{Aut}(x_0)$ に等しく，したがって，
和が自己同型の合成と一致するような自然な $k$-ベクトル空間構造をもつ。
\end{lemma}

\begin{proof}
集合としての等しさは Remark \ref{formal-defos-remark-infaut-lifting-equalities} の
末尾で述べたものであり，ベクトル空間構造に関する主張は
Lemma \ref{formal-defos-lemma-Aut-functor-tangent-space} から従う。
\end{proof}

\begin{lemma}
\label{formal-defos-lemma-k-linear-infaut}
$\varphi : \mathcal{F} \to \mathcal{G}$ を，(RS) を満たす
$\mathcal{C}_\Lambda$ 上の，群亜群をファイバーとする余ファイバー圏の射とする。
$x_0 \in \Ob(\mathcal{F}(k))$ とする。このとき $\varphi$ は
$k$-線形写像
$\text{Inf}_{x_0}(\mathcal{F}) \to \text{Inf}_{\varphi(x_0)}(\mathcal{G})$
を誘導する。
\end{lemma}

\begin{proof}
$\varphi$ が恒等射を恒等射に写す射
$\mathit{Aut}(x_0) \to \mathit{Aut}(\varphi(x_0))$
を誘導することは明らかである。したがって，接空間に対する結果
Lemma \ref{formal-defos-lemma-k-linear-differential} から従う。
\end{proof}

\begin{lemma}
\label{formal-defos-lemma-lifted-automorphisms-torsor}
$\mathcal{F}$ を (RS) を満たす，$\mathcal{C}_\Lambda$ 上で
群亜群をファイバーとする余ファイバー圏とする。$x' \to x$ を，
核 $I$ が $\mathfrak m_{A'}$ で消される全射な環写像 $A' \to A$ 上にある
射とする。$x_0$ を $x$ の $\mathcal{F}(k)$ への順像とする。
このとき $\text{Inf}(x'/x)$ は
$T_{\text{id}_{x_0}} \mathit{Aut}(x') \otimes_k I
= \text{Inf}_{x_0}(\mathcal{F}) \otimes_k I$
による自由かつ推移的な作用をもつ。
\end{lemma}

\begin{proof}
これは自己同型層の状況における
Lemma \ref{formal-defos-lemma-free-transitive-action} の類似物にすぎない。
正確には，Remark \ref{formal-defos-remark-localize-cofibered-groupoid} を関手
$\mathit{Aut}(x') : \mathcal{C}_{A'} \to \textit{Sets}$ と元
% P10-E0483: preserve the frozen wrong Aut(x)(k) functor value sidecar-only.
$\text{id}_{x_0} \in \mathit{Aut}(x)(k)$ に適用すると，前変形関手
$F = \mathit{Aut}(x')_{\text{id}_{x_0}}$ を得る。
Lemmas \ref{formal-defos-lemma-Aut-functor-RS} and \ref{formal-defos-lemma-localize-RS} により
$F$ は変形関手である。したがって
Lemma \ref{formal-defos-lemma-free-transitive-action} は
$\text{Lift}(\text{id}_x, A')$ 上の $TF \otimes_k I$ による
自由かつ推移的な作用を与える。実際，
$\text{Lift}(\text{id}_x, A')$ は群なので常に空でない。
ベクトル空間の等式
% P10-E0484: preserve the frozen missing tensor factor on the first TF sidecar-only.
$$
TF = T_{\text{id}_{x_0}} \mathit{Aut}(x') \otimes_k I =
\text{Inf}_{x_0}(\mathcal{F}) \otimes_k I
$$
が Lemma \ref{formal-defos-lemma-Aut-functor-tangent-space} により成り立つことに注意する。
Remark \ref{formal-defos-remark-infaut-lifting-equalities} の等式
$\text{Inf}(x'/x) = \text{Lift}(\text{id}_x, A')$ が証明を完結させる。
\end{proof}

\begin{lemma}
\label{formal-defos-lemma-infaut-trivial}
$\mathcal{F}$ を (RS) を満たす，$\mathcal{C}_\Lambda$ 上で
群亜群をファイバーとする余ファイバー圏とする。$x' \to x$ を，
全射な環写像上にある $\mathcal{F}$ の射とする。$x_0$ を $x$ の
$\mathcal{F}(k)$ への順像とする。$\text{Inf}_{x_0}(\mathcal{F}) = 0$ ならば
$\text{Inf}(x'/x) = 0$ である。
\end{lemma}

\begin{proof}
Lemmas \ref{formal-defos-lemma-factor-small-extension} and
\ref{formal-defos-lemma-lifted-automorphisms-torsor} から従う。
\end{proof}

\begin{lemma}
\label{formal-defos-lemma-infdef-trivial}
$\mathcal{F}$ を (RS) を満たす，$\mathcal{C}_\Lambda$ 上で
群亜群をファイバーとする余ファイバー圏とする。
$x_0 \in \Ob(\mathcal{F}(k))$ とする。このとき
$\text{Inf}_{x_0}(\mathcal{F}) = 0$ であるための必要十分条件は，
群亜群をファイバーとする余ファイバー圏の自然な射
$\mathcal{F}_{x_0} \to \overline{\mathcal{F}_{x_0}}$ が同値であることである。
\end{lemma}

\begin{proof}
射 $\mathcal{F}_{x_0} \to \overline{\mathcal{F}_{x_0}}$ が同値であるための
必要十分条件は，$\mathcal{F}_{x_0}$ がセトイドをファイバーとすることである。
cf.\ Categories, Section \ref{categories-section-fibred-in-setoids}
（セトイドとは，定義により，任意の対象の自己同型が恒等射だけである
群亜群である）。$\text{Inf}_{x_0}(\mathcal{F}) = 0$ であることと，
この条件が $\mathcal{F}_{x_0}$ に対して成り立つことが同値であると示す。
$\mathcal{F}_{x_0}$ がセトイドをファイバーとするならば
$\text{Inf}_{x_0}(\mathcal{F}) = 0$ であることは明らかである。
逆に，$\text{Inf}_{x_0}(\mathcal{F}) = 0$ と仮定する。
$A$ を $\mathcal{C}_\Lambda$ の対象とする。このとき
Lemma \ref{formal-defos-lemma-infaut-trivial} により，
$\mathcal{F}_{x_0}(A)$ の任意の対象 $x \to x_0$ に対して
$\text{Inf}(x/x_0) = 0$ である。定義により $\text{Inf}(x/x_0)$ は
$\mathcal{F}_{x_0}(A)$ における $x \to x_0$ の自己同型群に等しいので，
$\mathcal{F}_{x_0}(A)$ がセトイドであることが示された。
\end{proof}

\section{応用}
\label{formal-defos-section-applications}

\noindent
後で用いる変形圏に関する結果をいくつかまとめる。

\begin{lemma}
\label{formal-defos-lemma-deformation-categories-fiber-product-morphisms}
$f : \mathcal{H} \to \mathcal{F}$ および
$g : \mathcal{G} \to \mathcal{F}$ を変形圏の $1$-射とする。このとき
\begin{enumerate}
\item $\mathcal{W} = \mathcal{H} \times_\mathcal{F} \mathcal{G}$ は
変形圏であり，
\item ベクトル空間の $6$-項完全列
$$
0 \to \text{Inf}(\mathcal{W})
\to \text{Inf}(\mathcal{H}) \oplus \text{Inf}(\mathcal{G})
\to \text{Inf}(\mathcal{F}) \to
T\mathcal{W} \to T\mathcal{H} \oplus T\mathcal{G} \to T\mathcal{F}
$$
が存在する。
\end{enumerate}
\end{lemma}

\begin{proof}
(1) は Lemma \ref{formal-defos-lemma-RS-fiber-product-morphisms} と，
$\mathcal{W}(k)$ が，一意な同型類をもつセトイド上の，
それぞれ一意な同型類をもつ二つのセトイドのファイバー積であることから従う。

\medskip\noindent
(2)．$w_0 \in \Ob(\mathcal{W}(k))$ とし，$x_0, y_0, z_0$ をそれぞれ
$\mathcal{F}, \mathcal{H}, \mathcal{G}$ における $w_0$ の像とする。
このとき $\text{Inf}(\mathcal{W}) = \text{Inf}_{w_0}(\mathcal{W})$ であり，
$\mathcal{H}$，$\mathcal{G}$，$\mathcal{F}$ についても同様である。
Remark \ref{formal-defos-remark-trivial-aut-point} を参照せよ。
Lemmas \ref{formal-defos-lemma-k-linear-differential} and
\ref{formal-defos-lemma-k-linear-infaut} を適用すると，「境界写像」
$\delta : \text{Inf}_{x_0}(\mathcal{F}) \to T\mathcal{W}$
を除くすべての線形写像を得る。適切な符号は後で挿入する。

\medskip\noindent
$\delta$ の構成。$k \to k[\epsilon]$ に沿う順像
$w_0 \to w'_0$ を選ぶ。$\mathcal{F}, \mathcal{H}, \mathcal{G}$ における
$w'_0$ の像を $x'_0, y'_0, z'_0$ と書く。とくに，同型
$b' : f(y'_0) \to x'_0$ および $c' : x'_0 \to g(z'_0)$ を得る。
$k[\epsilon] \to k$ に沿う順像を
$b : f(y_0) \to x_0$ および $c : x_0 \to g(z_0)$ と書く。
圏のファイバー積の明示的な形で述べれば，これは
$w'_0 = (k[\epsilon], y'_0, z'_0, c' \circ b')$ および
$w_0 = (k, y_0, z_0, c \circ b)$ であることを意味する。
Remarks \ref{formal-defos-remarks-cofibered-groupoids} (\ref{formal-defos-item-fibre-product}) を参照せよ。
$\alpha : x'_0 \to x'_0$ が与えられたとき，
$\delta(\alpha) = (k[\epsilon], y'_0, z'_0, c' \circ \alpha \circ b')$
と置く。これは確かに $k[\epsilon]$ 上の $\mathcal{W}$ の対象であり，
$\alpha$ の $k$ 上への順像は恒等射なので，$k[\epsilon] \to k$ 上の射
$(k[\epsilon], y'_0, z'_0, c' \circ \alpha \circ b') \to w_0$
を伴う。より一般に，任意の $k$-ベクトル空間 $V$ に対し，
まったく同じ公式を用いて写像
$$
\text{Lift}(\text{id}_{x_0}, k[V])
\longrightarrow
\text{Lift}(w_0, k[V])
$$
を定義できる。この構成はベクトル空間 $V$ に関して関手的である
（詳細は省略する）。したがって Lemma \ref{formal-defos-lemma-morphism-linear-functors}
を適用すれば，$\delta$ は $k$-線形である。

\medskip\noindent
これらの写像を構成すると，列が完全であることは容易に示せる。
第一の写像の単射性は，
$f \times g : \mathcal{W} \to \mathcal{H} \times \mathcal{G}$
が忠実であることから従う。
$(\beta, \gamma) \in
\text{Inf}_{y_0}(\mathcal{H}) \oplus \text{Inf}_{z_0}(\mathcal{G})$
が $\text{Inf}_{x_0}(\mathcal{F})$ の同じ元に写るならば，
$(\beta, \gamma)$ は
$w'_0 = (k[\epsilon], y'_0, z'_0, c' \circ b')$ の自己同型を定めるので，
第二項での完全性を得る。上の $\alpha$ が $w_0$ の自明な変形
$(k[\epsilon], y'_0, z'_0, c' \circ \alpha \circ b')$ を与えるならば，
同型
$w'_0 = (k[\epsilon], y'_0, z'_0, c' \circ b') \to
(k[\epsilon], y'_0, z'_0, c' \circ \alpha \circ b')$
は $\alpha$ の逆像となる組 $(\beta, \gamma)$ を生じる。
$w = (k[\epsilon], y, z, \phi)$ が $w_0$ の変形で，
$y'_0 \cong y$ かつ $z \cong z'_0$ ならば，写像
$$
f(y'_0) \to f(y) \xrightarrow{\phi} g(z) \to g(z'_0)
$$
は $\delta$ によって $w$ に写る $\alpha$ である。
最後に，$y$ と $z$ が $y_0$ と $z_0$ の変形であり，
$f(y_0) = x_0 = g(z_0)$ の変形の同型
$\phi : f(y) \to g(z)$ が存在するならば，
% P10-E0485: preserve the frozen wrong target pair (x, y) sidecar-only.
$T\mathcal{W}$ における $(x, y)$ の逆像
$w = (k[\epsilon], y, z, \phi)$ を得る。これで証明は完了する。
\end{proof}

\begin{lemma}
\label{formal-defos-lemma-map-fibre-products-smooth}
$\mathcal{H}_1 \to \mathcal{G}$，$\mathcal{H}_2 \to \mathcal{G}$，および
$\mathcal{G} \to \mathcal{F}$ を，$\mathcal{C}_\Lambda$ 上で群亜群を
ファイバーとする余ファイバー圏の写像とする。次を仮定する：
\begin{enumerate}
\item $\mathcal{F}$ と $\mathcal{G}$ は変形圏である；
\item $T\mathcal{G} \to T\mathcal{F}$ は単射である；
\item $\text{Inf}(\mathcal{G}) \to \text{Inf}(\mathcal{F})$ は全射である。
\end{enumerate}
このとき
$\mathcal{H}_1 \times_\mathcal{G} \mathcal{H}_2 \to
\mathcal{H}_1 \times_\mathcal{F} \mathcal{H}_2$ は滑らかである。
\end{lemma}

\begin{proof}
% P10-E0486: render the intended notation while preserving the frozen grammar sidecar-only.
与えられた写像を $p_i : \mathcal{H}_i \to \mathcal{G}$ および
$q : \mathcal{G} \to \mathcal{F}$ と書く。
$A' \to A$ を $\mathcal{C}_\Lambda$ の小拡大とする。
$A'$ 上の $\mathcal{H}_1 \times_\mathcal{F} \mathcal{H}_2$ の対象は
三つ組 $(x'_1, x'_2, a')$ である。ここで $x'_i$ は
$A'$ 上の $\mathcal{H}_i$ の対象であり，
$a' : q(p_1(x'_1)) \to q(p_2(x'_2))$ は $A'$ 上の
$\mathcal{F}$ のファイバー圏の射である。$A' \to A$ に沿う順像をとると
$(x_1, x_2, a)$ を得る。これを $A$ 上の
$\mathcal{H}_1 \times_\mathcal{G} \mathcal{H}_2$ の対象へ持ち上げるとは，
$q(b) = a$ を満たす $A$ 上の射
$b : p_1(x_1) \to p_2(x_2)$ を見つけることである。
したがって，$q$ による像が $a'$ であるような射
$b' : p_1(x'_1) \to p_2(x'_2)$ へ $b$ を持ち上げられることを示せばよい。

\medskip\noindent
$$
p_1(x'_1) \to p_1(x_1) \xrightarrow{b} p_2(x_2)
\quad\text{および}\quad
p_2(x'_2) \to p_2(x_2)
$$
を $\textit{Lift}(p_2(x_2), A' \to A)$ の二つの対象とみなせることに注意する。
関手 $q$ は，これらの対象を
% P10-E0487: preserve the frozen b label after applying q sidecar-only.
$$
q(p_1(x'_1)) \to q(p_1(x_1)) \xrightarrow{b} q(p_2(x_2))
\quad\text{および}\quad
q(p_2(x'_2)) \to q(p_2(x_2))
$$
という $\textit{Lift}(q(p_2(x_2)), A' \to A)$ の二つの対象へ写す。
これらは写像 $a' : q(p_1(x'_1)) \to q(p_2(x'_2))$ を用いて同型になる。
一方，関手
$$
q : \textit{Lift}(p_2(x_2), A' \to A) \to
\textit{Lift}(q(p_2(x_2)), A' \to A)
$$
は，Lemma \ref{formal-defos-lemma-free-transitive-action} と接空間に関する仮定により，
同型類上で単射を定める。したがって，$A$ への順像が $b$ である射
$b' : p_1(x_1') \to p_2(x'_2)$ が存在することが分かる。
しかし，$q(b') = a'$ を達成するには $b'$ の選択を調整する必要が
あるかもしれない。そのためには
$q : \text{Inf}(p_2(x'_2)/p_2(x_2)) \to
\text{Inf}(q(p_2(x'_2))/q(p_2(x_2)))$
が全射であることを示せば十分である。これは無限小自己同型に関する仮定と
Lemma \ref{formal-defos-lemma-lifted-automorphisms-torsor} から従う。
\end{proof}

\begin{lemma}
\label{formal-defos-lemma-easy-check-smooth}
$f : \mathcal{F} \to \mathcal{G}$ を変形圏の写像とする。
$x_0 \in \Ob(\mathcal{F}(k))$ とし，その像を
$y_0 \in \Ob(\mathcal{G}(k))$ とする。次を仮定する：
\begin{enumerate}
\item 写像 $T\mathcal{F} \to T\mathcal{G}$ は全射である；
\item $\mathcal{C}_\Lambda$ の任意の小拡大 $A' \to A$ と，
その像が $y \in \mathcal{G}(A)$ である
$x \in \mathcal{F}(A)$ に対し，$y$ の $A'$ への持ち上げが存在するならば，
$x$ の $A'$ への持ち上げが存在する。
\end{enumerate}
このとき $\mathcal{F} \to \mathcal{G}$ は滑らかである（逆も成り立つ）。
\end{lemma}

\begin{proof}
$A' \to A$ を小拡大とする。$x \in \mathcal{F}(A)$ とする。
$y' \to f(x)$ を $A' \to A$ 上にある $\mathcal{G}$ の射とする。
% P10-E0488: preserve frozen reversed Lift arguments and set/category notation sidecar-only.
$f$ により誘導される関手
$\text{Lift}(A', x) \to \text{Lift}(A', f(x))$ を考える。
Lemma \ref{formal-defos-lemma-smoothness-small-extensions} により，
$y' \to f(x)$ に写る $\text{Lift}(A', x)$ の対象
$x' \to x$ が存在することを示さなければならない。
条件 (2) により，$\text{Lift}(A', x)$ が空圏でないことが分かる。
% P10-E0489: preserve the frozen second citation to condition (2) sidecar-only.
条件 (2) と Lemma \ref{formal-defos-lemma-free-transitive-action} により，
同型類上の写像は所望のとおり全射であると結論できる。
\end{proof}

\begin{lemma}
\label{formal-defos-lemma-map-between-smooth}
$\mathcal{F} \to \mathcal{G} \to \mathcal{H}$ を，
$\mathcal{C}_\Lambda$ 上で群亜群をファイバーとする余ファイバー圏の写像とする。
次を仮定する：
\begin{enumerate}
\item $\mathcal{F}$，$\mathcal{G}$ は変形圏である；
\item 写像 $T\mathcal{F} \to T\mathcal{G}$ は全射である；
\item $\mathcal{F} \to \mathcal{H}$ は滑らかである。
\end{enumerate}
このとき $\mathcal{F} \to \mathcal{G}$ は滑らかである。
\end{lemma}

\begin{proof}
$A' \to A$ を $\mathcal{C}_\Lambda$ の小拡大とし，
$x \in \mathcal{F}(A)$ の像を $y \in \mathcal{G}(A)$ とする。
持ち上げ $y' \in \mathcal{G}(A')$ が存在すると仮定する。
Lemma \ref{formal-defos-lemma-easy-check-smooth} によれば，$x$ にも持ち上げがあることだけを
確認すればよい。$y'$ の像を
$z' \in \mathcal{H}(A')$ とする。$\mathcal{F} \to \mathcal{H}$ は
滑らかなので，$x \in \mathcal{F}(A)$ と $z' \in \mathcal{H}(A')$ の
双方へ写る $x' \in \mathcal{F}(A')$ が存在する。
Definition \ref{formal-defos-definition-smooth-morphism} を参照せよ。これで証明は完了する。
\end{proof}

\section{任意の圏上の関手における群亜群}
\label{formal-defos-section-groupoids-arbitrary}

\noindent
まず，任意の圏上の関手における群亜群について一般的な事項を述べる。
次節では圏 $\mathcal{C}_\Lambda$ へ移る。明確に区別するため，
通常の群亜群，すなわちすべての射が同型である圏を，
群亜群圏と呼ぶことがある。

\begin{definition}
\label{formal-defos-definition-groupoid-in-functors}
$\mathcal{C}$ を圏とする。$\mathcal{C}$ 上の
{\it 関手における群亜群の圏}とは，次の対象と射をもつ圏である。
\begin{enumerate}
\item 対象：$\mathcal{C}$ 上の{\it 関手における群亜群}とは五つ組
$(U, R, s, t, c)$ であって，$U, R : \mathcal{C} \to \textit{Sets}$ は
関手，$s, t : R \to U$ および $c : R \times_{s, U, t} R \to R$ は
次の性質をもつ射である：$\mathcal{C}$ の任意の対象 $T$ に対して，五つ組
$$
(U(T), R(T), s, t, c)
$$
は群亜群圏である。
\item 射：$\mathcal{C}$ 上の関手における群亜群の
{\it 射 $(U, R, s, t, c) \to (U', R', s', t', c')$}とは，
次の性質をもつ射 $U \to U'$ と $R \to R'$ からなる：
$\mathcal{C}$ の任意の対象 $T$ に対し，誘導写像
$U(T) \to U'(T)$ と $R(T) \to R'(T)$ は群亜群圏の間の関手
$$
(U(T), R(T), s, t, c) \to (U'(T), R'(T), s', t', c').
$$
を定める。
\end{enumerate}
\end{definition}

\begin{remark}
\label{formal-defos-remark-confusion-groupoids-in-functors}
$\mathcal{C}$ 上の関手における群亜群は，関手
$\mathcal{C} \to \textit{Groupoids}$ のデータと同じであり，
$\mathcal{C}$ 上の関手における群亜群の射は，対応する関手
$\mathcal{C} \to \textit{Groupoids}$ の射と同じである
（ここで $\textit{Groupoids}$ は 1-圏とみなす）。しかし本章の目的には，
関手における群亜群という用語を使う方が便利である。実際，
関手における群亜群を対応する関手
$\mathcal{C} \to \textit{Groupoids}$ と考えたり，同じことであるが，
その関手に付随する，群亜群をファイバーとする余ファイバー圏と考えたりすると，
混同を招くことがある
（Remark \ref{formal-defos-remark-smooth-groupoid-in-functors-warning}）。
\end{remark}

\begin{remark}
\label{formal-defos-remark-identity-inverse}
$(U, R, s, t, c)$ を圏 $\mathcal{C}$ 上の関手における群亜群とする。
一意な射 $e : U \to R$ と $i : R \to R$ が存在して，
$\mathcal{C}$ の任意の対象 $T$ に対し，$e: U(T) \to R(T)$ は
$x \in U(T)$ を $x$ の恒等射へ写し，$i: R(T) \to R(T)$ は
% P10-E0491: preserve the frozen wrong membership a in U(T) sidecar-only.
$a \in U(T)$ を群亜群圏 $(U(T), R(T), s, t, c)$ における
$a$ の逆射へ写す。$s$，$t$，$c$，$e$，$i$ をそれぞれ
「始域」，「終域」，「合成」，「恒等」，「逆」と呼ぶことがある。
\end{remark}

\begin{definition}
\label{formal-defos-definition-representable}
$\mathcal{C}$ を圏とする。$\mathcal{C}$ 上の関手における群亜群が
{\it 表現可能}であるとは，$U$ と $R$ が $\mathcal{C}$ の対象であり，
押し出し $R \amalg_{s, U, t} R$ が存在するような形
$(\underline{U}, \underline{R}, s, t, c)$ のものに同型であることをいう。
\end{definition}

\begin{remark}
\label{formal-defos-remark-reason-existence-coproduct}
したがって，$\mathcal{C}$ 上の関手における表現可能な群亜群は，
$\mathcal{C}$ の対象 $U$ と $R$，および射 $s, t : U \to R$ と
$c : R \to R \amalg_{s, U, t} R$ であって，
$(\underline{U}, \underline{R}, s, t, c)$ が
Definition \ref{formal-defos-definition-groupoid-in-functors} の条件を満たすものによって
与えられる。押し出し $R \amalg_{s, U, t} R$ の存在を要求する理由は，
合成射 $c$ が $\mathcal{C}$ の射の水準で定義されるようにするためである。
以下で $\widehat{\mathcal{C}}_\Lambda$ 上の関手における表現可能な群亜群を
考えるとき，この要件は常に満たされる。実際，
Lemma \ref{formal-defos-lemma-CLambdahat-pushouts} により，圏
$\widehat{\mathcal{C}}_\Lambda$ は押し出しをもつ。
\end{remark}

\begin{remark}
\label{formal-defos-remark-simplify-terminology}
表現可能な群亜群が与えられていることを表すため，
「{\it $(\underline{U}, \underline{R}, s, t, c)$ を $\mathcal{C}$ 上の
関手における群亜群とする}」という。このことは，$U$ と $R$ が
$\mathcal{C}$ の対象であり，射 $s, t : U \to R$ が存在し，
押し出し $R \amalg_{s, U, t} R$ が存在し，射
$c : R \to R \amalg_{s, U, t} R$ が存在し，かつ
$(\underline{U}, \underline{R}, s, t, c)$ が $\mathcal{C}$ 上の
関手における群亜群であることを意味する。
\end{remark}

\noindent
関手における群亜群の制限に対する記法を導入する。これは以下で
$\widehat{\mathcal C}_\Lambda$ から $\mathcal{C}_\Lambda$ へ
制限する状況に関係する。

\begin{definition}
\label{formal-defos-definition-restricting-groupoids-in-functors}
$(U, R, s, t, c)$ を圏 $\mathcal{C}$ 上の関手における群亜群とする。
$\mathcal{C}'$ を $\mathcal{C}$ の部分圏とする。
$(U, R, s, t, c)$ の $\mathcal{C}'$ への
{\it 制限 $(U, R, s, t, c)|_{\mathcal{C}'}$}とは，
$(U|_{\mathcal{C}'}, R|_{\mathcal C'}, s|_{\mathcal{C}'},
t|_{\mathcal{C}'}, c|_{\mathcal{C}'})$ により与えられる，
$\mathcal{C}'$ 上の関手における群亜群である。
\end{definition}

\begin{remark}
\label{formal-defos-remark-notation-restriction}
Definition \ref{formal-defos-definition-restricting-groupoids-in-functors} の状況では，
$s|_{\mathcal{C}'}, t|_{\mathcal{C}'}, c|_{\mathcal{C}'}$ を単に
$s, t, c$ と書くことが多い。
\end{remark}

\begin{definition}
\label{formal-defos-definition-quotient}
$(U, R, s, t, c)$ を圏 $\mathcal{C}$ 上の関手における群亜群とする。
\begin{enumerate}
\item 対応 $T \mapsto  (U(T), R(T), s, t, c)$ は関手
$\mathcal{C} \to \textit{Groupoids}$ を定める。
{\it 商となる，群亜群をファイバーとする余ファイバー圏
$[U/R] \to \mathcal{C}$}とは，この関手に付随する
$\mathcal{C}$ 上で群亜群をファイバーとする余ファイバー圏である
（Remarks \ref{formal-defos-remarks-cofibered-groupoids}
(\ref{formal-defos-item-construction-associated-cofibered-groupoid}) と同様）。
\item {\it 商射 $U \to [U/R]$}とは，次の規則で定められる，
$\mathcal{C}$ 上で群亜群をファイバーとする余ファイバー圏の射である：
\begin{enumerate}
\item $x \in U(T)$ を対象 $(T, x) \in \Ob([U/R](T))$ へ写す；
\item $x \in U(T)$ と $f : T \to T'$ から，$f : T \to T'$ 上にある射
$(f, \text{id}_{U(f)(x)}): (T, x) \to (T, U(f)(x))$ が生じる。
\end{enumerate}
\end{enumerate}
\end{definition}

\section{基底圏上の関手における群亜群}
\label{formal-defos-section-prorepresentable-groupoids-in-functors}

\noindent
この節では，$\mathcal{C}_\Lambda$ 上の関手における群亜群を論じる。
最終的な目標は，$\mathcal{C}_\Lambda$ 上の関手における
プロ表現可能な群亜群が，代数空間における滑らかな群亜群が
代数スタックの「表示」として働くのと同じ仕方で，
振る舞いのよい変形圏の「表示」として働くことを示すことである。
cf.\ 『代数スタック』Section \ref{algebraic-section-stack-to-presentation}。

\begin{definition}
\label{formal-defos-definition-prorepresentable-groupoid-in-functors}
$\mathcal{C}_\Lambda$ 上の関手における群亜群が
{\it プロ表現可能}であるとは，$\widehat{\mathcal{C}}_\Lambda$ 上の
関手におけるある表現可能な群亜群
$(\underline{R_0}, \underline{R_1}, s, t, c)$ に対する
$(\underline{R_0}, \underline{R_1}, s, t, c)|_{\mathcal{C}_\Lambda}$
に同型であることをいう。
\end{definition}

\noindent
$(U, R, s, t, c)$ を $\mathcal{C}_\Lambda$ 上の関手における
群亜群とする。完備化をとると，五つ組
$(\widehat{U}, \widehat{R}, \widehat{s}, \widehat{t}, \widehat{c})$ を得る。
Remark \ref{formal-defos-remark-completion-restriction-cofset-adjoint} により，
$\text{CofSet}(\mathcal{C}_\Lambda)$ 上の関手としての完備化は右随伴なので，
極限と可換する。とくに標準同型
$$
\widehat{R \times_{s, U, t} R}
\longrightarrow
\widehat{R} \times_{\widehat{s}, \widehat{U}, \widehat{t}} \widehat{R},
$$
があるので，$\widehat{c}$ を関手
$\widehat{R} \times_{\widehat{s}, \widehat{U}, \widehat{t}} \widehat{R} \to
\widehat{R}$ とみなせる。このとき
$(\widehat{U}, \widehat{R}, \widehat{s}, \widehat{t}, \widehat{c})$ は
$\widehat{\mathcal{C}}_\Lambda$ 上の関手における群亜群であり，
その恒等射と逆射は $(U, R, s, t, c)$ のそれらの完備化である。

\begin{definition}
\label{formal-defos-definition-completion-groupoid-in-functors}
$(U, R, s, t, c)$ を $\mathcal{C}_\Lambda$ 上の関手における
群亜群とする。$(U, R, s, t, c)$ の
{\it 完備化 $(U, R, s, t, c)^{\wedge}$}とは，上で記述した
$\widehat{\mathcal{C}}_\Lambda$ 上の関手における群亜群
$(\widehat{U}, \widehat{R}, \widehat{s}, \widehat{t}, \widehat{c})$ である。
\end{definition}

\begin{remark}
\label{formal-defos-remark-groupoid-in-functors-complete-restrict}
$(U, R, s, t, c)$ を $\mathcal{C}_\Lambda$ 上の関手における
群亜群とする。このとき標準同型
$(U, R, s, t, c)^{\wedge}|_{\mathcal{C}_\Lambda} \cong (U, R, s, t, c)$
がある。Remark \ref{formal-defos-remark-restrict-completion} を参照せよ。
一方，$(U, R, s, t, c)$ を $\widehat{\mathcal{C}}_\Lambda$ 上の
関手における群亜群とし，
$U, R : \widehat{\mathcal{C}}_\Lambda \to \textit{Sets}$ がともに
極限と可換するとする。例えば $U, R$ が表現可能ならばそうである。
このとき標準同型
$((U, R, s, t, c)|_{\mathcal{C}_\Lambda})^{\wedge} \cong (U, R, s, t, c)$
がある。これは Remark \ref{formal-defos-remark-restrict-complete-continuous-functor}
から従う。
\end{remark}

\begin{lemma}
\label{formal-defos-lemma-groupoid-in-functors-prorep-equivalences}
$(U, R, s, t, c)$ を $\mathcal{C}_\Lambda$ 上の関手における群亜群とする。
\begin{enumerate}
\item $(U, R, s, t, c)$ がプロ表現可能であるための必要十分条件は，
その完備化が $\widehat{\mathcal{C}}_\Lambda$ 上の関手における
群亜群として表現可能であることである。
\item $(U, R, s, t, c)$ がプロ表現可能であるための必要十分条件は，
$U$ と $R$ がプロ表現可能であることである。
\end{enumerate}
\end{lemma}

\begin{proof}
(1) は Remark \ref{formal-defos-remark-groupoid-in-functors-complete-restrict} から従う。
(2) の「必要」方向は，関手におけるプロ表現可能な群亜群の定義から明らかである。
逆に，$U$ と $R$ がプロ表現可能であると仮定する。すなわち，
$\widehat{\mathcal{C}}_\Lambda$ の対象 $R_0$ と $R_1$ に対して
$U \cong \underline{R_0}|_{\mathcal{C}_\Lambda}$ および
$R \cong \underline{R_1}|_{\mathcal{C}_\Lambda}$ であるとする。
Remark \ref{formal-defos-remark-restrict-complete-continuous-functor} により
$\underline{R_0} \cong \widehat{\underline{R_0}|_{\mathcal{C}_\Lambda}}$
および
$\underline{R_1} \cong \widehat{\underline{R_1}|_{\mathcal{C}_\Lambda}}$
なので，完備化 $(U, R, s, t, c)^\wedge$ は
$(\underline{R_0}, \underline{R_1}, \widehat{s}, \widehat{t}, \widehat{c})$
という形の関手における群亜群である。
Lemma \ref{formal-defos-lemma-CLambdahat-pushouts} により，押し出し
% P10-E0492: preserve the frozen fibre-product symbol and wrong middle object sidecar-only.
$\underline{R_1} \times_{\widehat{s}, \underline{R_1}, \widehat{t}}
\underline{R_1}$
が存在する。したがって
$(\underline{R_0}, \underline{R_1}, \widehat{s}, \widehat{t}, \widehat{c})$
は $\widehat{\mathcal{C}}_\Lambda$ 上の関手における表現可能な群亜群である。
最後に，Remark \ref{formal-defos-remark-groupoid-in-functors-complete-restrict} により，制限
$(\underline{R_0}, \underline{R_1}, s, t, c)|_{\mathcal{C}_\Lambda}$
は $(U, R, s, t, c)$ を再び与える。ゆえに定義により
$(U, R, s, t, c)$ はプロ表現可能である。
\end{proof}

\section{基底圏上の関手における滑らかな群亜群}
\label{formal-defos-section-smooth-minimal-groupoids-in-functors}

\noindent
$\mathcal{C}_\Lambda$ 上の関手における群亜群の滑らかさを
次のように定義する。

\begin{definition}
\label{formal-defos-definition-smooth-groupoid-in-functors}
$(U, R, s, t, c)$ を $\mathcal{C}_\Lambda$ 上の関手における群亜群とする。
$s, t: R \to U$ が滑らかなとき，$(U, R, s, t, c)$ は
{\it 滑らか}であるという。
\end{definition}

\begin{remark}
\label{formal-defos-remark-smooth-groupoid-in-functors-warning}
この用語は混同を招きうることに注意する：
$(U, R, s, t, c)$ が関手における滑らかな群亜群であっても，商
$[U/R]$ は Definition \ref{formal-defos-definition-cofibered-groupoid-projection-smooth}
で定義した滑らかな，群亜群をファイバーとする余ファイバー圏とは限らない。
しかし，$(U, R, s, t, c)$ の滑らかさは商射 $U \to [U/R]$ の
滑らかさを含意する（実際には同値である）ことを
Lemma \ref{formal-defos-lemma-smooth-quotient-morphism} で見る。
\end{remark}

\begin{remark}
\label{formal-defos-remark-smooth-power-series-prorepresentable-smooth-groupoid-in-functors}
$(\underline{R_0}, \underline{R_1}, s, t, c)|_{\mathcal{C}_\Lambda}$ を
$\mathcal{C}_\Lambda$ 上の関手におけるプロ表現可能な群亜群とする。
このとき $(\underline{R_0}, \underline{R_1}, s, t, c)|_{\mathcal{C}_\Lambda}$
が滑らかであるための必要十分条件は，
% P10-E0493: preserve the frozen missing ring/algebra noun sidecar-only.
$R_1$ が $s$ と $t$ の双方を介して $R_0$ 上の冪級数であることである。
これは Lemma \ref{formal-defos-lemma-smooth-morphism-power-series} から従う。
\end{remark}

\begin{lemma}
\label{formal-defos-lemma-smooth-quotient-morphism}
$(U, R, s, t, c)$ を $\mathcal{C}_\Lambda$ 上の関手における群亜群とする。
次は同値である：
\begin{enumerate}
\item 関手における群亜群 $(U, R, s, t, c)$ は滑らかである。
\item 射 $s : R \to U$ は滑らかである。
\item 射 $t : R \to U$ は滑らかである。
\item 商射 $U \to [U/R]$ は滑らかである。
\end{enumerate}
\end{lemma}

\begin{proof}
$(U, R, s, t, c)$ の逆射 $i: R \to R$ は同型であり，
$t = s \circ i$ なので，主張 (2) と (3) は同値である。
定義により (1) は (2) と (3) がともに成り立つことと同値であり，
したがってそれぞれ一方だけとも同値である。

\medskip\noindent
最後に，(2) と (4) が同値であることを示す。定義を展開すると：
\begin{enumerate}
\item[(2)] $s: R \to U$ の滑らかさは次の条件を意味する：
$f: B \to A$ が $\mathcal{C}_\Lambda$ の全射な環写像であり，
$a \in R(A)$ および $y \in U(B)$ が $s(a) = U(f)(y)$ を満たすならば，
$R(f)(a') = a$ および $s(a') = y$ を満たす $a' \in R(B)$ が存在する。
\item[(4)] $U \to [U/R]$ の滑らかさは次の条件を意味する：
$f: B \to A$ が $\mathcal{C}_\Lambda$ の全射な環写像であり，
$(f, a) : (B, y) \to (A, x)$ が $[U/R]$ の射であるならば，
$s(b) = x'$，$t(b) = y$，$c(a, R(f)(b)) = e(x)$ を満たす
$x' \in U(B)$ と $b \in R(B)$ が存在する。ここで $e : U \to R$ は
恒等射を表し，記法 $(f, a)$ は Remarks \ref{formal-defos-remarks-cofibered-groupoids}
(\ref{formal-defos-item-construction-associated-cofibered-groupoid}) と同じである。
とくに，$a \in R(A)$ であり，$s(a) = U(f)(y)$ および $t(a) = x$ を満たす。
\end{enumerate}
(4) が成り立ち，(2) のとおり $f, a, y$ が与えられたとする。
$x = t(a)$ と置けば，射 $(f, a): (B, y) \to (A, x)$ がある。
このとき (4) は $x'$ と $b$ を与え，$a' = i(b)$ は (2) の要件を満たす。
逆に，(2) が成り立ち，(4) のとおり
$(f, a): (B, y) \to (A, x)$ が与えられたとする。このとき (2) は
$a' \in R(B)$ を与え，$x' = t(a')$ と $b = i(a')$ は (4) の要件を満たす。
\end{proof}

\section{関手における群亜群の商としての変形圏}
\label{formal-defos-section-deformation-categories-as-quotients}

\noindent
$\mathcal{C}_\Lambda$ 上の関手における群亜群について，その商が
変形圏となることを保証する条件を論じ，そのような商の接空間と
無限小自己同型空間を計算する。

\begin{lemma}
\label{formal-defos-lemma-smooth-RS-groupoid-in-functors-quotient}
$(U, R, s, t, c)$ を $\mathcal{C}_\Lambda$ 上の関手における
滑らかな群亜群とする。$U$ と $R$ が (RS) を満たすと仮定する。
このとき $[U/R]$ は (RS) を満たす。
\end{lemma}

\begin{proof}
$[U/R]$ における図式
$$
\xymatrix{
                           &     (A_2, x_2) \ar[d]^{(f_2, a_2)} \\
(A_1, x_1) \ar[r]^{(f_1, a_1)} &     (A, x)
}
$$
であって，$f_2: A_2 \to A$ が全射であるものを考える。記法は
Remarks \ref{formal-defos-remarks-cofibered-groupoids}
(\ref{formal-defos-item-construction-associated-cofibered-groupoid}) と同じである。
したがって $f_1: A_1 \to A, f_2: A_2 \to A$ は
$\mathcal{C}_\Lambda$ の写像，$x \in U(A)$，$x_1 \in U(A_1)$，
$x_2 \in U(A_2)$ であり，$a_1, a_2 \in R(A)$ は
$s(a_1) = U(f_1)(x_1)$，$t(a_1) = x$ および
$s(a_2) = U(f_2)(x_2)$，$t(a_2) = x$ を満たす。
この $[U/R]$ の図式に対し，$A_1 \times_A A_2$ 上にある
ファイバー積を次のように構成する。

\medskip\noindent
$i: R \to R$ を逆射として，$a = c(i(a_1), a_2)$ と置く。
このとき $a \in R(A)$，$x_2 \in U(A_2)$，かつ
$s(a) = U(f_2)(x_2)$ である。したがって元
$(a, x_2) \in R(A) \times_{s, U(A), U(f_2)} U(A_2)$ がある。
$s : R \to U$ の滑らかさにより，
$R(f_2)(\widetilde{a}) = a$ および $s(\widetilde{a}) = x_2$ を満たす元
$\widetilde{a} \in R(A_2)$ が存在する。とくに
$U(f_2)(t(\widetilde{a})) = t(a) = U(f_1)(x_1)$ である。
したがって $x_1$ と $t(\widetilde{a})$ は元
$$
(x_1, t(\widetilde{a})) \in U(A_1) \times_{U(A)} U(A_2).
$$
を定める。$U$ が (RS) を満たすという仮定により，同一視
$U(A_1) \times_{U(A)} U(A_2) = U(A_1 \times_A A_2)$ がある。
% P10-E0494: render the notation introduction naturally; frozen missing preposition is sidecar-only.
$(x_1, t(\widetilde{a})) \in U(A_1) \times_{U(A)} U(A_2)$ に対応する元を
$x_1 \times t(\widetilde{a}) \in U(A_1 \times_A A_2)$ と書く。
$p_1, p_2$ を $A_1 \times_A A_2$ の射影とする。次を主張する：
$$
\xymatrix{
(A_1 \times_A A_2, x_1 \times t(\widetilde{a})) \ar[d]_{(p_1, e(x_1))}
\ar[rr]_-{(p_2, i(\widetilde{a}))} & & (A_2, x_2) \ar[d]^{(f_2, a_2)} \\
(A_1, x_1) \ar[rr]^{(f_1, a_1)} & & (A, x)
}
$$
これは $[U/R]$ におけるファイバー平方である
（$e: U \to R$ は恒等射を表すことに注意せよ）。

\medskip\noindent
等式 $c(a_2, R(f_2)(i(\widetilde{a}))) = c(a_2, i(a)) = a_1$
により，この図式は可換である。これがファイバー平方であることを確認するため，
$[U/R]$ における可換図式
$$
\xymatrix{
(B, z) \ar[d]_{(g_1, b_1)} \ar[rr]_{(g_2, b_2)} & & (A_2, x_2)
\ar[d]^{(f_2, a_2)} \\
(A_1, x_1) \ar[rr]^{(f_1, a_1)} & & (A, x)
}
$$
を考える。$(A_1, x_1)$ および $(A_2, x_2)$ への射と両立する一意な射
$(g, b) : (B, z) \to (A_1 \times_A A_2, x_1 \times t(\widetilde{a}))$
が存在することを示す。$g = (g_1, g_2) : B \to A_1 \times_A A_2$
と取らなければならない。仮定により $R$ は (RS) を満たすので，同一視
$R(A_1 \times_A A_2) = R(A_1) \times_{R(A)} R(A_2)$ がある。
したがって，$R(A)$ で一致するある
$b'_1 \in R(A_1)$，$b'_2 \in R(A_2)$ により
$b = (b'_1, b'_2)$ と書ける。このとき
$((g_1, g_2), (b'_1, b'_2)) : (B, z) \to
(A_1 \times_A A_2, x_1 \times t(\widetilde{a}))$
が射影と可換するための必要十分条件は，
$b'_1 = b_1$ および $b'_2 = c(\widetilde{a}, b_2)$ である。
これにより一意性と存在が示される。
\end{proof}

\begin{lemma}
\label{formal-defos-lemma-deformation-groupoid-quotient}
$(U, R, s, t, c)$ を $\mathcal{C}_\Lambda$ 上の関手における
滑らかな群亜群とする。$U$ と $R$ が変形関手であると仮定する。このとき：
\begin{enumerate}
\item 商 $[U/R]$ は変形圏である。
\item $[U/R]$ の接空間は
$$
T[U/R] = \Coker(ds-dt: TR \to TU).
$$
である。
\item $[U/R]$ の無限小自己同型空間は
$$
\text{Inf}([U/R]) =
\Ker(ds \oplus dt : TR \to TU \oplus TU).
$$
である。
\end{enumerate}
\end{lemma}

\begin{proof}
$U$ と $R$ は変形関手なので，$[U/R]$ は前変形圏である。
定義により変形関手については (RS) が成り立つので，
% P10-E0495: preserve the frozen bare quotient notation outside math mode sidecar-only.
Lemma \ref{formal-defos-lemma-smooth-RS-groupoid-in-functors-quotient} により
(RS) は [U/R] について成り立つ。したがって $[U/R]$ は変形圏である。
主張 (2) と (3) は定義から直接従う。
\end{proof}

\section{群亜群をファイバーとする余ファイバー圏の表示}
\label{formal-defos-section-presentation-categories-cofibred-in-groupoids}

\noindent
表示を次のように定義する。

\begin{definition}
\label{formal-defos-definition-presentation}
$\mathcal{F}$ を圏 $\mathcal{C}$ 上で群亜群をファイバーとする
余ファイバー圏とする。$(U, R, s, t, c)$ を $\mathcal{C}$ 上の
関手における群亜群とする。
{\it $(U, R, s, t, c)$ による $\mathcal{F}$ の表示}とは，
$\mathcal{C}$ 上で群亜群をファイバーとする余ファイバー圏の同値
$\varphi : [U/R] \to \mathcal{F}$ のことである。
\end{definition}

\noindent
次の二つの一般的な補題を用いて表示を得る。

\begin{lemma}
\label{formal-defos-lemma-presentation-construction}
% P10-E0496: the repeated missing English article is recorded sidecar-only.
$\mathcal{F}$ を圏 $\mathcal{C}$ 上で群亜群をファイバーとする
余ファイバー圏とする。$U : \mathcal{C} \to \textit{Sets}$ を関手とする。
$f : U \to \mathcal{F}$ を，$\mathcal{C}$ 上で群亜群をファイバーとする
余ファイバー圏の射とする。$R, s, t, c$ を次のように定義する：
\begin{enumerate}
\item $R : \mathcal{C} \to \textit{Sets}$ は関手
$U \times_{f, \mathcal{F}, f} U$ である。
\item $t, s : R \to U$ はそれぞれ第一射影と第二射影である。
\item $c : R \times_{s, U, t} R \to R$ は，
$U \times_{f, \mathcal{F}, f} U \times_{f, \mathcal{F}, f} U$
の第一因子と最後の因子への射影によって，標準同型
$R \times_{s, U, t} R \to
U \times_{f, \mathcal{F}, f} U \times_{f, \mathcal{F}, f} U$
の下で与えられる射である。
\end{enumerate}
このとき $(U, R, s, t, c)$ は $\mathcal{C}$ 上の関手における群亜群である。
\end{lemma}

\begin{proof}
省略する。
\end{proof}

\begin{lemma}
\label{formal-defos-lemma-presentation-morphism}
% P10-E0496: the second repeated missing English article is recorded sidecar-only.
$\mathcal{F}$ を圏 $\mathcal{C}$ 上で群亜群をファイバーとする
余ファイバー圏とする。$U : \mathcal{C} \to \textit{Sets}$ を関手とする。
$f : U \to \mathcal{F}$ を，$\mathcal{C}$ 上で群亜群をファイバーとする
余ファイバー圏の射とする。$(U, R, s, t, c)$ を，
Lemma \ref{formal-defos-lemma-presentation-construction} において
$f : U \to \mathcal{F}$ から構成された $\mathcal{C}$ 上の関手における
群亜群とする。このとき自然な射 $[f] : [U/R] \to \mathcal{F}$ が存在し，
次を満たす：
\begin{enumerate}
\item $[f]: [U/R] \to \mathcal{F}$ は充満忠実である。
\item $[f]: [U/R] \to \mathcal{F}$ が同値であるための必要十分条件は，
$f : U \to \mathcal{F}$ が本質的全射であることである。
\end{enumerate}
\end{lemma}

\begin{proof}
省略する。
\end{proof}

\section{変形圏の表示}
\label{formal-defos-section-presentation-deformation-categories}

\noindent
次の補題によれば，前変形関手から前変形圏 $\mathcal{F}$ への
滑らかな射は，$\mathcal{F}$ の，関手における滑らかな群亜群による
表示を与える。

\begin{lemma}
\label{formal-defos-lemma-smooth-groupoid-in-functors-construction}
$\mathcal{F}$ を $\mathcal{C}_\Lambda$ 上で群亜群をファイバーとする
余ファイバー圏とする。$U : \mathcal{C}_\Lambda \to \textit{Sets}$ を
関手とする。$f : U \to \mathcal{F}$ を，群亜群をファイバーとする
余ファイバー圏の滑らかな射とする。このとき：
\begin{enumerate}
\item $(U, R, s, t, c)$ が，Lemma \ref{formal-defos-lemma-presentation-construction}
において $f : U \to \mathcal{F}$ から構成された
$\mathcal{C}_\Lambda$ 上の関手における群亜群ならば，
$(U, R, s, t, c)$ は滑らかである。
\item $f : U(k) \to \mathcal{F}(k)$ が本質的全射ならば，
Lemma \ref{formal-defos-lemma-presentation-morphism} の射
$[f] : [U/R] \to \mathcal{F}$ は同値である。
\end{enumerate}
\end{lemma}

\begin{proof}
Lemma \ref{formal-defos-lemma-presentation-construction} の構成から，可換図式
$$
\xymatrix{
R = U \times_{f, \mathcal{F}, f} U \ar[r]_-s \ar[d]_t & U
\ar[d]^f \\
U \ar[r]^f & \mathcal{F}
}
$$
を得る。ここで $t, s$ は第一射影と第二射影である。したがって
Lemma \ref{formal-defos-lemma-smooth-properties} により $t, s$ は滑らかである。
ゆえに (1) が成り立つ。

\medskip\noindent
(2) の仮定が成り立つならば，
Lemma \ref{formal-defos-lemma-smooth-morphism-essentially-surjective} により，
射 $f : U \to \mathcal{F}$ は本質的全射である。したがって
Lemma \ref{formal-defos-lemma-presentation-morphism} により，射
$[f] : [U/R] \to \mathcal{F}$ は同値である。
\end{proof}

\begin{lemma}
\label{formal-defos-lemma-deformation-functor-diagonal}
$\mathcal{F}$ を変形圏とする。
$U : \mathcal{C}_\Lambda \to \textit{Sets}$ を変形関手とする。
$f: U \to \mathcal{F}$ を，群亜群をファイバーとする余ファイバー圏の
射とする。このとき $U \times_{f, \mathcal{F}, f} U$ は変形関手であり，
その接空間は $k$-ベクトル空間の完全列
$$
0 \to \text{Inf}(\mathcal{F}) \to
T(U \times_{f, \mathcal{F}, f} U) \to TU \oplus TU
$$
に入る。
\end{lemma}

\begin{proof}
Lemma \ref{formal-defos-lemma-deformation-categories-fiber-product-morphisms} と
$\text{Inf}(U) = (0)$ であることから従う。
\end{proof}

\begin{lemma}
\label{formal-defos-lemma-prorepresentable-groupoid-in-functors-construction}
$\mathcal{F}$ を変形圏とする。
$U : \mathcal{C}_\Lambda \to \textit{Sets}$ をプロ表現可能関手とする。
$f : U \to \mathcal{F}$ を，群亜群をファイバーとする余ファイバー圏の
射とする。$(U, R, s, t, c)$ を，Lemma \ref{formal-defos-lemma-presentation-construction}
において $f : U \to \mathcal{F}$ から構成された
$\mathcal{C}_\Lambda$ 上の関手における群亜群とする。
$\dim_k \text{Inf}(\mathcal{F}) < \infty$ ならば，
$(U, R, s, t, c)$ はプロ表現可能である。
\end{lemma}

\begin{proof}
Example \ref{formal-defos-example-prorepresentable-deformation-functor} により，
$U$ は変形関手であることに注意する。
Lemma \ref{formal-defos-lemma-deformation-functor-diagonal} により，
$R = U \times_{f, \mathcal{F}, f} U$ は変形関手であり，その接空間
$TR = T(U \times_{f, \mathcal{F}, f} U)$ は完全列
$0 \to \text{Inf}(\mathcal{F}) \to TR \to TU \oplus TU$ に入る。
第一の空間が有限次元であると仮定しており，また
Example \ref{formal-defos-example-tangent-space-prorepresentable-functor} により
$TU$ は有限次元なので，$\dim TR < \infty$ である。
% P10-E0497: set off the frozen undelimited see-reference parenthetically; sidecar-only.
写像 $\gamma : \text{Der}_\Lambda(k, k) \to TR$
（(\ref{formal-defos-equation-map}) を参照）は単射である。実際，その
$TR \to TU$ との合成は，プロ表現可能関手 $U$ に対する
Theorem \ref{formal-defos-theorem-Schlessinger-prorepresentability} により単射である。
したがって Theorem \ref{formal-defos-theorem-Schlessinger-prorepresentability} により
$R$ はプロ表現可能である。ゆえに
Lemma \ref{formal-defos-lemma-groupoid-in-functors-prorep-equivalences} から，
$(U, R, s, t, c)$ はプロ表現可能である。
\end{proof}

\begin{theorem}
\label{formal-defos-theorem-presentation-deformation-groupoid}
$\mathcal{F}$ を $\mathcal{C}_\Lambda$ 上で群亜群をファイバーとする
余ファイバー圏とする。このとき，$\mathcal{F}$ が
$\mathcal{C}_\Lambda$ 上の関手における滑らかなプロ表現可能群亜群による
表示をもつための必要十分条件は，次の条件が成り立つことである：
\begin{enumerate}
\item $\mathcal{F}$ は変形圏である。
\item $\dim_k T\mathcal{F}$ は有限である。
\item $\dim_k \text{Inf}(\mathcal{F})$ は有限である。
\end{enumerate}
\end{theorem}

\begin{proof}
プロ表現可能関手は変形関手であることを思い出そう；
Example \ref{formal-defos-example-prorepresentable-deformation-functor} を参照せよ。
% P10-E0498: preserve the frozen direct equivalence claim; type mismatch remains sidecar-only.
したがって，$\mathcal{F}$ が関手における滑らかなプロ表現可能群亜群と
同値ならば，条件 (1)，(2)，(3) は
Lemma \ref{formal-defos-lemma-deformation-groupoid-quotient} の (1)，(2)，(3) から従う。

\medskip\noindent
逆に，条件 (1)，(2)，(3) が成り立つと仮定する。条件 (1) は
(S1) と (S2) が満たされることを意味する；
Lemma \ref{formal-defos-lemma-RS-implies-S1-S2} を参照せよ。
Lemma \ref{formal-defos-lemma-versal-object-existence} により，
versal 形式対象 $\xi$ が存在する。
% P10-E0499: Japanese supplies the natural clause boundary; frozen missing comma is sidecar-only.
$U = \underline{R}|_{\mathcal{C}_\Lambda}$ と置くと，付随する写像
$\underline{\xi} : U \to \mathcal{F}$ は滑らかである
（これは versal 形式対象の定義である）。
$(U, R, s, t, c)$ を，Lemma \ref{formal-defos-lemma-presentation-construction} において
写像 $\underline{\xi}$ から構成された関手における群亜群とする。
Lemma \ref{formal-defos-lemma-smooth-groupoid-in-functors-construction} により，
$(U, R, s, t, c)$ は関手における滑らかな群亜群であり，
$[U/R] \to \mathcal{F}$ は同値である。さらに
Lemma \ref{formal-defos-lemma-prorepresentable-groupoid-in-functors-construction} により，
$(U, R, s, t, c)$ はプロ表現可能である。したがって
$[U/R] \to \mathcal{F}$ は求める $\mathcal{F}$ の表示である。
\end{proof}

\section{極小性に関する注意}
\label{formal-defos-section-minimality}

\noindent
本章の主定理は上の
Theorem \ref{formal-defos-theorem-presentation-deformation-groupoid} である。
これは，$\mathcal{C}_\Lambda$ 上で群亜群をファイバーとする
余ファイバー圏のうち，関手における滑らかなプロ表現可能群亜群による
表示をもつものを完全に記述する。本節では，
Sections \ref{formal-defos-section-minimal-versal} および
\ref{formal-defos-section-miniversal-objects-existence} で論じた極小性を用いて，
「極小」な滑らかなプロ表現可能表示を得る方法を簡単に論じる。

\begin{definition}
\label{formal-defos-definition-minimal-groupoid-in-functors}
$(U, R, s, t, c)$ を $\mathcal{C}_\Lambda$ 上の関手における
滑らかなプロ表現可能群亜群とする。
\begin{enumerate}
\item 群亜群 $(U(k[\epsilon]), R(k[\epsilon]), s, t, c)$ が
完全不連結，すなわち相異なる対象の間に射が存在しないとき，
$(U, R, s, t, c)$ は {\it 正規化されている}という。
% P10-E0500: the frozen spurious English article is article-neutral here; sidecar-only.
\item $U \to [U/R]$ が $[U/R]$ の極小 versal 形式対象によって
与えられるとき，$(U, R, s, t, c)$ は {\it 極小}であるという。
\end{enumerate}
\end{definition}

\noindent
二つの概念の違いは，条件 (\ref{formal-defos-equation-bijective}) と
(\ref{formal-defos-equation-bijective-orbits}) の違いに関係し，
$k' \subset k$ が分離的ならば消滅する。また，次の補題が示すように，
関手における正規化された滑らかなプロ表現可能群亜群は極小である。
正確な主張は次のとおりである。

\begin{lemma}
\label{formal-defos-lemma-characterize-minimal-groupoid-in-functors}
$(U, R, s, t, c)$ を $\mathcal{C}_\Lambda$ 上の関手における
滑らかなプロ表現可能群亜群とする。
\begin{enumerate}
\item $(U, R, s, t, c)$ が正規化されているための必要十分条件は，
射 $U \to [U/R]$ が接空間上の同型を誘導することである。
\item $(U, R, s, t, c)$ が極小であるための必要十分条件は，
$TU \to T[U/R]$ の核が
$\text{Der}_\Lambda(k, k) \to TU$ の像に含まれることである。
\end{enumerate}
\end{lemma}

\begin{proof}
(1) は定義から直ちに従う。
(2) を見るため，$\mathcal{F} = [U/R]$ と置く。$\mathcal{F}$ は
表示をもつので変形圏である；
Theorem \ref{formal-defos-theorem-presentation-deformation-groupoid} を参照せよ。
とくに (RS)，(S1)，(S2) を満たす；
Lemma \ref{formal-defos-lemma-RS-implies-S1-S2} を参照せよ。
極小 versal 形式対象は同型を除いて一意であることを思い出そう；
Lemma \ref{formal-defos-lemma-minimal-versal} を参照せよ。
Theorem \ref{formal-defos-theorem-miniversal-object-existence} により，
極小 versal 対象は (\ref{formal-defos-equation-bijective-orbits}) を満たす写像
$\underline{\xi} : \underline{R}|_{\mathcal{C}_\Lambda} \to \mathcal{F}$
を誘導する。
% P10-E0501: preserve the frozen over-F isomorphism assertion; proof gap remains sidecar-only.
$\mathcal{F}$ 上で $U \cong \underline{R}|_{\mathcal{C}_\Lambda}$ なので，
$TU \to T\mathcal{F} = T[U/R]$ は補題で述べた性質を満たす。
\end{proof}

\noindent
$\mathcal C_\Lambda$ 上の関手における極小プロ表現可能群亜群の商は，
自己同型でない自己同値をもたない。これを証明するため，まず次の補題に注意する。

\begin{lemma}
\label{formal-defos-lemma-surjective-morphism-prorepresentable-functor}
$U: \mathcal{C}_\Lambda \to \textit{Sets}$ をプロ表現可能関手とする。
$\varphi : U \to U$ を，$d\varphi : TU \to TU$ が同型であるような射とする。
このとき $\varphi$ は同型である。
\end{lemma}

\begin{proof}
ある $R \in \Ob(\widehat{\mathcal{C}}_\Lambda)$ に対して
$U \cong \underline{R}|_{\mathcal{C}_\Lambda}$ ならば，
$\varphi$ を完備化することで射 $\underline{R} \to \underline{R}$ を得る。
$f: R \to R$ を $\widehat{\mathcal{C}}_\Lambda$ における対応する射とすると，
$f$ は同型 $\text{Der}_\Lambda(R, k) \to \text{Der}_\Lambda(R, k)$ を誘導する；
Example \ref{formal-defos-example-tangent-space-map-prorepresentable-functor} を参照せよ。
とくに Lemma \ref{formal-defos-lemma-derivations-surjective} により $f$ は全射である。
Noether 環の全射自己準同型は同型なので
（Algebra, Lemma \ref{algebra-lemma-surjective-endo-noetherian-ring-is-iso}
を参照），$f$，したがって $\underline{R} \to \underline{R}$，したがって
$\varphi : U \to U$ は同型である。
\end{proof}

\begin{lemma}
\label{formal-defos-lemma-minimal-prorepresentable-groupoid-autoequivalence}
$(U, R, s, t, c)$ を $\mathcal{C}_\Lambda$ 上の関手における
極小な滑らかなプロ表現可能群亜群とする。
$\varphi : [U/R] \to [U/R]$ が，群亜群をファイバーとする
余ファイバー圏の同値ならば，$\varphi$ は同型である。
\end{lemma}

\begin{proof}
% P10-E0502: preserve the frozen strict quotient/groupoid morphism identification; gap sidecar-only.
射 $\varphi : [U/R] \to [U/R]$ は，
Definition \ref{formal-defos-definition-groupoid-in-functors} で定義した
$\mathcal{C}_\Lambda$ 上の関手における群亜群の射
$\varphi : (U, R, s, t, c) \to (U, R, s, t, c)$ と同じものである。
対応する射を $\phi : U \to U$ および $\psi : R \to R$ と書く。
図式
$$
\xymatrix{
& \text{Der}_\Lambda(k, k) \ar[dr]_\gamma \ar[dl]^\gamma \\
TU \ar[rr]_{d\phi} \ar[d] & & TU \ar[d]  \\
T[U/R] \ar[rr]^{d\varphi} & & T[U/R]
}
$$
は可換であり，$d\varphi$ は全単射であり，さらに
Lemma \ref{formal-defos-lemma-characterize-minimal-groupoid-in-functors} による
極小性の特徴付けがあるので，$d\phi$ は単射である
（したがって次元の理由により全単射である）。ゆえに
Lemma \ref{formal-defos-lemma-surjective-morphism-prorepresentable-functor} により
$\phi : U \to U$ は同型である。完全列
$$
0 \to \text{Inf}([U/R]) \to TR \to TU \oplus TU
$$
（Lemma \ref{formal-defos-lemma-deformation-functor-diagonal}）を用いる同様の議論により，
$\psi : R \to R$ が同型であることを証明できる。
% P10-E0503: supply the demonstrative subject missing in frozen English; sidecar-only.
しかしこれは，$R = U \times_{[U/R]} U$ であり，$\varphi$ と $\phi$ が
同型であることからも従う。
\end{proof}

\begin{lemma}
\label{formal-defos-lemma-minimal-prorepresentable-groupoid-equivalence}
$(U, R, s, t, c)$ と $(U', R', s', t', c')$ を
$\mathcal{C}_\Lambda$ 上の関手における極小な滑らかな
プロ表現可能群亜群とする。$\varphi : [U/R] \to [U'/R']$ が，
群亜群をファイバーとする余ファイバー圏の同値ならば，
$\varphi$ は同型である。
\end{lemma}

\begin{proof}
$\psi : [U'/R'] \to [U/R]$ を $\varphi$ の擬逆とする。
Lemma \ref{formal-defos-lemma-minimal-prorepresentable-groupoid-autoequivalence} により
$\psi \circ \varphi$ と $\varphi \circ \psi$ は同型なので，
$\varphi$ と $\psi$ は同型である。
\end{proof}

\noindent
次の補題は，本章で先に見た事柄のいくつかをまとめる。

\begin{lemma}
\label{formal-defos-lemma-minimal-groupoid-in-functors-construction}
$\mathcal{F}$ を，$\dim_k T\mathcal{F} <\infty$ および
$\dim_k \text{Inf}(\mathcal{F}) < \infty$ を満たす変形圏とする。
このとき $\mathcal{F}$ の極小 versal 形式対象 $\xi$ が存在する。
% P10-E0504: preserve the frozen ring/relation-functor R collision; sidecar-only.
$\xi$ が $R \in \Ob(\widehat{\mathcal{C}}_\Lambda)$ 上にあるとする。
$U = \underline{R}|_{\mathcal{C}_\Lambda}$ と置く。
$f = \underline{\xi} : U \to \mathcal{F}$ を付随する射とする。
$(U, R, s, t, c)$ を，Lemma \ref{formal-defos-lemma-presentation-construction} において
$f : U \to \mathcal{F}$ から構成された $\mathcal{C}_\Lambda$ 上の
関手における群亜群とする。このとき $(U, R, s, t, c)$ は
$\mathcal{C}_\Lambda$ 上の関手における極小な滑らかなプロ表現可能群亜群であり，
同値 $[U/R] \to \mathcal{F}$ が存在する。
\end{lemma}

\begin{proof}
$\mathcal{F}$ は変形圏なので (S1) と (S2) を満たす；
Lemma \ref{formal-defos-lemma-RS-implies-S1-S2} を参照せよ。
Lemma \ref{formal-defos-lemma-versal-object-existence} により versal 形式対象が存在する。
Lemma \ref{formal-defos-lemma-minimal-versal} により，補題の主張のとおりの
極小 versal 形式対象 $\xi/R$ が存在する。
% P10-E0505: Japanese supplies the natural clause boundary; frozen missing comma is sidecar-only.
$U = \underline{R}|_{\mathcal{C}_\Lambda}$ と置くと，付随する写像
$\underline{\xi} : U \to \mathcal{F}$ は滑らかである
（これは versal 形式対象の定義である）。
$(U, R, s, t, c)$ を，Lemma \ref{formal-defos-lemma-presentation-construction} において
写像 $\underline{\xi}$ から構成された関手における群亜群とする。
Lemma \ref{formal-defos-lemma-smooth-groupoid-in-functors-construction} により，
$(U, R, s, t, c)$ は関手における滑らかな群亜群であり，
$[U/R] \to \mathcal{F}$ は同値である。
Lemma \ref{formal-defos-lemma-prorepresentable-groupoid-in-functors-construction} により，
$(U, R, s, t, c)$ はプロ表現可能である。
% P10-E0506: preserve the frozen equivalence-as-equality formula; sidecar-only.
最後に，$U \to [U/R] = \mathcal{F}$ は極小 versal 形式対象 $\xi$ に
対応するので，$(U, R, s, t, c)$ は極小である。
\end{proof}

\noindent
関手における極小プロ表現可能群亜群による表示は，次の一意性を満たす。

\begin{lemma}
\label{formal-defos-lemma-minimal-presentations-equivalent}
% P10-E0507: the frozen missing English article is article-neutral here; sidecar-only.
$\mathcal{F}$ を $\mathcal{C}_\Lambda$ 上で群亜群をファイバーとする
余ファイバー圏とする。$\mathcal{F}$ が，関手における極小な滑らかな
プロ表現可能群亜群 $(U, R, s, t, c)$ および
$(U', R', s', t', c')$ による表示をもつと仮定する。
このとき $(U, R, s, t, c)$ と $(U', R', s', t', c')$ は同型である。
\end{lemma}

\begin{proof}
Lemma \ref{formal-defos-lemma-minimal-prorepresentable-groupoid-equivalence} と，
% P10-E0502: preserve the second frozen strict quotient/groupoid identification; sidecar-only.
射 $[U/R] \to [U'/R']$ は関手における群亜群の射と同じものであるという
観察から従う（Definition \ref{formal-defos-definition-quotient} における
$[U/R]$ の明示的構成による）。
\end{proof}

\noindent
以上をまとめると，次の定理を証明したことになる。

\begin{theorem}
\label{formal-defos-theorem-minimal-smooth-prorepresentable-presentations}
$\mathcal{F}$ を $\mathcal{C}_\Lambda$ 上で群亜群をファイバーとする
余ファイバー圏とする。次の条件を考える：
\begin{enumerate}
\item $\mathcal{F}$ は，$\mathcal{C}_\Lambda$ 上の関手における
正規化された滑らかなプロ表現可能群亜群による表示をもつ。
\item $\mathcal{F}$ は，$\mathcal{C}_\Lambda$ 上の関手における
滑らかなプロ表現可能群亜群による表示をもつ。
\item $\mathcal{F}$ は，$\mathcal{C}_\Lambda$ 上の関手における
極小な滑らかなプロ表現可能群亜群による表示をもつ。
\item $\mathcal{F}$ は次の条件を満たす：
\begin{enumerate}
\item $\mathcal{F}$ は変形圏である。
\item $\dim_k T\mathcal{F}$ は有限である。
\item $\dim_k \text{Inf}(\mathcal{F})$ は有限である。
\end{enumerate}
\end{enumerate}
このとき (2)，(3)，(4) は同値であり，(1) から従う。
$k' \subset k$ が分離的ならば，(1)，(2)，(3)，(4) はすべて同値である。
さらに，$\mathcal{F}$ の表示を与える関手における極小な滑らかな
プロ表現可能群亜群は，同型を除いて一意である。
\end{theorem}

\begin{proof}
Lemma \ref{formal-defos-lemma-characterize-minimal-groupoid-in-functors} により，
(1) は (3) を含意し，$k' \subset k$ が分離的ならば (3) と同値である。
(3) が (2) を含意することは明らかである。
Theorem \ref{formal-defos-theorem-presentation-deformation-groupoid} により，
(2) は (4) を含意する。
Lemma \ref{formal-defos-lemma-minimal-groupoid-in-functors-construction} により，
(4) は (3) を含意する。これですべての含意が証明された。
最後の一意性の主張は
Lemma \ref{formal-defos-lemma-minimal-presentations-equivalent} から従う。
\end{proof}

\section{versal 環の一意性}
\label{formal-defos-section-uniqueness}

\noindent
$\widehat{\mathcal{C}}_\Lambda$ の $R, S$ が与えられたとき，写像
$f, g : R \to S$ が {\it 形式的にホモトピック}であるとは，
ある $r \geq 0$ と，$\widehat{\mathcal{C}}_\Lambda$ における写像
$h : R \to R[[t_1, \ldots, t_r]]$ および
$k : R[[t_1, \ldots, t_r]] \to S$ が存在して，すべての $a \in R$ に対し
\begin{enumerate}
\item $h(a) \bmod (t_1, \ldots, t_r) = a$，
\item $f(a) = k(a)$，
\item $g(a) = k(h(a))$
\end{enumerate}
が成り立つこととする。$(r, h, k)$ を $f$ と $g$ の間の
{\it 形式ホモトピー}という。

\begin{lemma}
\label{formal-defos-lemma-homotopy}
形式的にホモトピックであることは，$\widehat{\mathcal{C}}_\Lambda$ における
射の集合上の同値関係である。
\end{lemma}

\begin{proof}
任意の $r \geq 1$ と二つの写像
$h_1, h_2 : R \to R[[t_1, \ldots, t_r]]$ が与えられ，すべての
$a \in R$ に対して
$h_1(a) \bmod (t_1, \ldots, t_r) =  h_2(a) \bmod (t_1, \ldots, t_r) = a$
を満たし，さらに写像 $k : R[[t_1, \ldots, t_r]] \to S$ が
与えられているとする。このとき $k \circ h_1$ と $k \circ h_2$ は
形式的にホモトピックであると主張する。
% P10-E0508: use the required noun 対称性; frozen wrong English word form remains sidecar-only.
この主張に内在する対称性から，形式的にホモトピックという概念が
対称であることが分かる。実際，写像
$$
\Psi :
R[[t_1, \ldots, t_r]]
\longrightarrow
R[[t_1, \ldots, t_r]],\quad
\sum a_I t^I \longmapsto \sum h_1(a_I)t^I
$$
は同型である。$a \in R$ に対して $h(a) = \Psi^{-1}(h_2(a))$ と置き，
$k' = k \circ \Psi$ と置けば，$(r, h, k')$ は $k \circ h_1$ と
$k \circ h_2$ の間の形式ホモトピーであり，主張が示される。
% P10-E0509: Japanese supplies the frozen missing terminal punctuation; sidecar-only.

\medskip\noindent
上のとおり三つの写像 $f_1, f_2, f_3 : R \to S$ と，
$f_1$ と $f_2$ の間の形式ホモトピー $(r_1, h_1, k_1)$，および
$f_3$ と $f_2$ の間の形式ホモトピー $(r_2, h_2, k_2)$（！）が
与えられているとする。座標を付け替えることにより，
$h_2 : R \to R[[t_{r_1 + 1}, \ldots, t_{r_1 + r_2}]]$
および
$k_2 : R[[t_{r_1 + 1}, \ldots, t_{r_1 + r_2}]] \to S$
と仮定してよい。適切な同型
$$
R[[t_1, \ldots, t_{r_1 + r_2}]]
\longrightarrow
R[[t_{r_1 + 1}, \ldots, t_{r_1 + r_2}]]
\widehat{\otimes}_{h_2, R, h_1}
R[[t_1, \ldots, t_{r_1}]]
$$
を選ぶことにより，これらの写像を可換図式
$$
\xymatrix{
R \ar[r]_{h_1} \ar[d]^{h_2} &
R[[t_1, \ldots, t_{r_1}]] \ar[d]^{h_2'} \\
R[[t_{r_1 + 1}, \ldots, t_{r_1 + r_2}]] \ar[r]^{h_1'} &
R[[t_1, \ldots, t_{r_1 + r_2}]]
}
$$
に組み込める。ここで $1 \leq i \leq r_1$ に対して
$h_2'(t_i) = t_i$ であり，
% P10-E0510: preserve the frozen upper bound r_2 rather than silently correcting it.
$r_1 + 1 \leq i \leq r_2$ に対して $h_1'(t_i) = t_i$ である。
詳細の一部は省略する。この図式は圏 $\widehat{\mathcal{C}}_\Lambda$ における
押し出しであり（Lemma \ref{formal-defos-lemma-CLambdahat-pushouts} の証明を参照），
$k_1 \circ h_1 = f_2 = k_2 \circ h_2$ なので，写像
$$
k : R[[t_1, \ldots, t_{r_1 + r_2}]] \to S
$$
であって $k_1 = k \circ h_2'$ および $k_2 = k \circ h_1'$ を満たすものが
存在する。$h = h_1' \circ h_2 = h_2' \circ h_1$ と置く。このとき
\begin{enumerate}
\item $k(h_1'(a)) = k_2(a) = f_3(a)$，
\item $k(h_2'(a)) = k_1(a) = f_1(a)$
\end{enumerate}
が成り立つ。証明の最初の段落の主張により，$f_1$ と $f_3$ は
形式的にホモトピックである。
\end{proof}

\begin{lemma}
\label{formal-defos-lemma-composition-homotopic}
圏 $\widehat{\mathcal{C}}_\Lambda$ において，
$f_1, f_2 : R \to S$ が形式的にホモトピックであり，
$g : S \to S'$ が射ならば，$g \circ f_1$ と $g \circ f_2$ は
形式的にホモトピックである。
\end{lemma}

\begin{proof}
実際，$(r, h, k)$ が $f_1$ と $f_2$ の間の形式ホモトピーならば，
$(r, h, g \circ k)$ は $g \circ f_1$ と $g \circ f_2$ の間の
形式ホモトピーである。
\end{proof}

\begin{lemma}
\label{formal-defos-lemma-versal-unique-up-to-homotopy}
$\mathcal{F}$ を $\mathcal{C}_\Lambda$ 上の変形圏で，
$\dim_k T\mathcal{F} < \infty$ および
$\dim_k \text{Inf}(\mathcal{F}) < \infty$ を満たすものとする。
$\xi$ を $R$ 上にある versal 形式対象とし，$\eta$ を $S$ 上にある
形式対象とする。このとき，
$$
f, g : R \to S
$$
であって $f_*\xi \cong \eta \cong g_*\xi$ を満たす任意の二つの写像は，
形式的にホモトピックである。
\end{lemma}

\begin{proof}
Theorem \ref{formal-defos-theorem-presentation-deformation-groupoid} とその証明により，
$\mathcal{F}$ は滑らかなプロ表現可能群亜群
$$
(\underline{R}, \underline{R_1}, s, t, c, e, i)|_{\mathcal{C}_\Lambda}
$$
% P10-E0511/P10-E0512: preserve lowercase lambda and the frozen incomplete F-clause.
による表示をもつ。これは $\mathcal{C}_\lambda$ 上の関手におけるもので，
しかも $\mathcal{F}$。
このとき写像 $s : R \to R_1$ と $t : R \to R_1$ は形式的に滑らかな
環写像であり，$e : R_1 \to R$ は切断である。とくに，ある
$r \geq 0$ に対して同型 $R_1 = R[[t_1, \ldots, t_r]]$ を選んで，
$s$ を埋め込み $R \subset R[[t_1, \ldots, t_r]]$ とし，
$t$ を，すべての $a \in R$ に対して
$h(a) \bmod (t_1, \ldots, t_r) = a$ を満たす写像
$h : R \to R[[t_1, \ldots, t_r]]$ に対応させることができる。
同型 $\alpha : f_*\xi \to g_*\xi$ の存在は，
$f = k \circ s$ および $g = k \circ t$ を満たす写像
$k : R_1 \to S$ が存在することにほかならない。これはまさに，
$(r, h, k)$ が $f$ と $g$ の間の形式ホモトピーであることを意味する。
\end{proof}

\begin{lemma}
\label{formal-defos-lemma-homotopic-minimal-prime}
圏 $\widehat{\mathcal{C}}_\Lambda$ において，
$f_1, f_2 : R \to S$ が形式的にホモトピックであり，
$\mathfrak p \subset R$ が極小素イデアルならば，イデアルとして
$f_1(\mathfrak p)S = f_2(\mathfrak p)S$ である。
\end{lemma}

\begin{proof}
$(r, h, k)$ が $f_1$ と $f_2$ の間の形式ホモトピーであると仮定する。
等式
$\mathfrak pR[[t_1, \ldots, t_r]] = h(\mathfrak p)R[[t_1, \ldots, t_r]]$
を主張する。この主張にさらに $k$ を合成すれば補題が従う。
主張を証明するため，写像
$\mathfrak p \mapsto \mathfrak pR[[t_1, \ldots, t_r]]$
が，$R$ の極小素イデアルと $R[[t_1, \ldots, t_r]]$ の極小素イデアルとの
間の全単射であることに注意する。最後に，$h$ は平坦なので
$h(\mathfrak p)R[[t_1, \ldots, t_r]]$ は極小素イデアルであり，
したがって今述べたことにより，ある極小素イデアル $\mathfrak q \subset R$
に対して $\mathfrak q R[[t_1, \ldots, t_r]]$ の形である。
しかし形式ホモトピーの定義により
$h \bmod (t_1, \ldots, t_r) = \text{id}_R$ なので，望むとおり
$\mathfrak q = \mathfrak p$ と結論する。
\end{proof}

\section{剰余体の変更}
\label{formal-defos-section-change-of-field}

\noindent
本節では，剰余体 $k$ を有限拡大で置き換えたときに何が起こるかを
手短に論じる。$\Lambda$ を Noether 環とし，$k$ を体として
$\Lambda \to k$ を有限環写像とする。本章全体を通して，剰余体が $k$ と
同一視された Artin 局所 $\Lambda$-代数の圏を
$\mathcal{C}_\Lambda$ と書いてきた；
Definition \ref{formal-defos-definition-CLambda} を参照せよ。
% P10-E0513: Japanese supplies the frozen subject-verb agreement naturally; sidecar-only.
しかし本節では $k$ を変更したときに何が起こるかを論じるので，
$k$ への依存を示すため，代わりに記法 $\mathcal{C}_{\Lambda, k}$ を用いる。

\begin{situation}
\label{formal-defos-situation-change-of-fields}
$\Lambda$ を Noether 環とし，$k$ と $l$ を体として
% P10-E0514/P10-E0515: render number agreement naturally; frozen English defects sidecar-only.
$\Lambda \to k \to l$ を有限環写像とする。したがって $l/k$ は有限体拡大である。
$\mathcal{C}_{\Lambda, l}$ の典型的対象を $B$，
$\mathcal{C}_{\Lambda, k}$ の典型的対象を $A$ と書く。次を定義する：
\begin{equation}
\label{formal-defos-equation-comparison}
\mathcal{C}_{\Lambda, l} \longrightarrow \mathcal{C}_{\Lambda, k},
\quad
B \longmapsto B \times_l k
\end{equation}
群亜群をファイバーとする余ファイバー圏
$p : \mathcal{F} \to \mathcal{C}_{\Lambda, k}$ が与えられたとき，
$\mathcal{F}_{l/k}(B) = \mathcal{F}(B \times_l k)$ と置くことにより，
付随する，群亜群をファイバーとする余ファイバー圏
$$
p_{l/k} : \mathcal{F}_{l/k} \longrightarrow \mathcal{C}_{\Lambda, l}
$$
を得る。
\end{situation}

\noindent
関手 (\ref{formal-defos-equation-comparison}) は意味をもつ。実際，
$B \times_l k \subset B$ なので
\begin{align*}
[k : k']\ \text{length}_{B \times_l k}(B \times_l k) & =
\text{length}_\Lambda(B \times_l k) \\
& \leq \text{length}_\Lambda(B) \\
& = [l : k']\ \text{length}_B(B) < \infty
\end{align*}
であり（Lemma \ref{formal-defos-lemma-length} 参照），したがって
$B \times_l k$ は Artin 環である
（Algebra, Lemma \ref{algebra-lemma-artinian-finite-length} 参照）。
ゆえに $B \times_l k$ は剰余体 $k$ をもつ Artin 局所環である。
(\ref{formal-defos-equation-comparison}) はファイバー積と可換することに注意せよ：
$$
(B_1 \times_B B_2) \times_l k =
(B_1 \times_l k) \times_{(B \times_l k)} (B_2 \times_l k)
$$
また，全射な環写像を全射な環写像に移す。
Remark \ref{formal-defos-remark-localize-cofibered-groupoid} の構成との混同を避けるため，
「高価な」記法 $\mathcal{F}_{l/k}$ を用いる。いくつかの初等的観察を挙げる。

\begin{lemma}
\label{formal-defos-lemma-elementary-properties-change-of-field}
記法と仮定は Situation \ref{formal-defos-situation-change-of-fields} のとおりとする。
\begin{enumerate}
\item $\overline{\mathcal{F}_{l/k}} = (\overline{\mathcal{F}})_{l/k}$ である。
\item $\mathcal{F}$ が前変形圏ならば，$\mathcal{F}_{l/k}$ は前変形圏である。
\item $\mathcal{F}$ が (S1) を満たすならば，$\mathcal{F}_{l/k}$ は (S1) を満たす。
\item $\mathcal{F}$ が (S2) を満たすならば，$\mathcal{F}_{l/k}$ は (S2) を満たす。
\item $\mathcal{F}$ が (RS) を満たすならば，$\mathcal{F}_{l/k}$ は (RS) を満たす。
\end{enumerate}
\end{lemma}

\begin{proof}
(1) は定義から直ちに従う。

\medskip\noindent
$\mathcal{F}_{l/k}(l) = \mathcal{F}(k)$ なので，(2) は定義から従う；
Definition \ref{formal-defos-definition-predeformation-category} を参照せよ。

\medskip\noindent
関手 (\ref{formal-defos-equation-comparison}) はファイバー積と可換し，全射写像を
全射写像に移すので，(3) が従う；
Definition \ref{formal-defos-definition-S1-S2} を参照せよ。

\medskip\noindent
(4) を示す。そのため，Definition \ref{formal-defos-definition-S1-S2} にあるような
$\mathcal{C}_{\Lambda, l}$ の図式
$$
\xymatrix{
          & l[\epsilon] \ar[d] \\
B  \ar[r] & l
}
$$
を考える。関手 (\ref{formal-defos-equation-comparison}) を適用すると
$$
\xymatrix{
          & k[l\epsilon] \ar[d] \\
B \times_l k  \ar[r] & k
}
$$
を得る。ここで $l\epsilon$ は有限次元 $k$-ベクトル空間
$l\epsilon \subset l[\epsilon]$ を表す。
Lemma \ref{formal-defos-lemma-S2-extensions} によれば，$\mathcal{F}$ に対する (S2) の
条件はこの図式についても成り立つ。したがって
$\mathcal{F}_{l/k}$ は (S2) を満たす。

\medskip\noindent
(5) は Lemma \ref{formal-defos-lemma-RS-2-categorical} の (2) における (RS) の
特徴付けと，(\ref{formal-defos-equation-comparison}) がファイバー積と可換することから従う。
\end{proof}

\noindent
次の補題は，とくに $\mathcal{F}$ が (S2) を満たす前変形圏である場合に
適用される；Lemma \ref{formal-defos-lemma-S1-S2-associated-functor} を参照せよ。

\begin{lemma}
\label{formal-defos-lemma-tangent-space-change-of-field}
記法と仮定は Situation \ref{formal-defos-situation-change-of-fields} のとおりとする。
$\mathcal{F}$ が前変形圏で，$\overline{\mathcal{F}}$ が (S2) を満たすと
仮定する。このとき，接空間の標準的な $l$-ベクトル空間同型
$$
T\mathcal{F} \otimes_k l \longrightarrow T\mathcal{F}_{l/k}
$$
が存在する。
\end{lemma}

\begin{proof}
Lemma \ref{formal-defos-lemma-elementary-properties-change-of-field} により，
$\mathcal{F}$ を $\overline{\mathcal{F}}$ で置き換えてよい。さらに，
$T\mathcal{F}$ と $T\mathcal{F}_{l/k}$ はそれぞれ標準的な
$k$-ベクトル空間構造と $l$-ベクトル空間構造をもつ；
Lemma \ref{formal-defos-lemma-tangent-space-vector-space} を参照せよ。このとき
$$
T\mathcal{F}_{l/k} = \mathcal{F}_{l/k}(l[\epsilon])
= \mathcal{F}(k[l\epsilon]) = T\mathcal{F} \otimes_k l
$$
であり，最後の等式は Lemma \ref{formal-defos-lemma-tangent-space-vector-space} による。
より一般に，有限次元 $l$-ベクトル空間 $V$ が与えられると
$$
\mathcal{F}_{l/k}(l[V]) = \mathcal{F}(k[V_k]) = T\mathcal{F} \otimes_k V_k
$$
である。ここで $V_k$ は $V$ を $k$-ベクトル空間とみたものを表す。
したがって，集合の圏への関手として，関手
$V \mapsto \mathcal{F}_{l/k}(l[V])$ と
$V \mapsto T\mathcal{F} \otimes_k V_k$ は標準的に同一視される。
Lemma \ref{formal-defos-lemma-linear-functor} により，各関手を $l$-線型関手にする方法は
高々一つである。ゆえにこれらの同型は $l$-ベクトル空間構造と両立し，
主張を得る。
\end{proof}

\begin{lemma}
\label{formal-defos-lemma-inf-aut-change-of-field}
記法と仮定は Situation \ref{formal-defos-situation-change-of-fields} のとおりとする。
$\mathcal{F}$ が変形圏であると仮定する。このとき，無限小自己同型空間の
標準的な $l$-ベクトル空間同型
$$
\text{Inf}(\mathcal{F}) \otimes_k l
\longrightarrow
\text{Inf}(\mathcal{F}_{l/k})
$$
が存在する。
\end{lemma}

\begin{proof}
$x_0 \in \Ob(\mathcal{F}(k))$ とし，
% P10-E0516: use the natural notation-introduction construction; frozen missing by is sidecar-only.
$l$ 上の $\mathcal{F}_{l/k}$ の対応する対象を $x_{l, 0}$ と書く。
$\text{Inf}(\mathcal{F}) = \text{Inf}_{x_0}(\mathcal{F})$ および
$\text{Inf}(\mathcal{F}_{l/k}) = \text{Inf}_{x_{l, 0}}(\mathcal{F}_{l/k})$
であることを思い出そう；Remark \ref{formal-defos-remark-trivial-aut-point} を参照せよ。
$\text{Inf}_{x_0}(\mathcal{F})$ 上のベクトル空間構造は，これを関手
$\mathit{Aut}(x_0)$ の接空間と同一視することから得られる。
この関手は，剰余体 $k$ をもつ Artin 局所 $k$-代数の圏
$\mathcal{C}_{k, k}$ 上で定義される。同様に，
$\text{Inf}_{x_{l, 0}}(\mathcal{F}_{l/k})$ は，剰余体 $l$ をもつ
Artin 局所 $l$-代数の圏 $\mathcal{C}_{l, l}$ 上で定義された
$\mathit{Aut}(x_{l, 0})$ の接空間である。定義を展開すると，
$\mathit{Aut}(x_{l, 0})$ は，$\mathcal{C}_{k, l}$ 上にある
$\mathit{Aut}(x_0)_{l/k}$ の $\mathcal{C}_{l, l}$ への制限である。
$\mathit{Aut}(x_0)_{l/k}$ を $\mathcal{C}_{k, l}$ 上の関手とみても
$\mathcal{C}_{l, l}$ 上の関手とみても接空間に違いはないので，補題は
Lemma \ref{formal-defos-lemma-tangent-space-change-of-field} と，
Lemma \ref{formal-defos-lemma-Aut-functor-RS} により $\mathit{Aut}(x_0)$ が (RS) を
満たすことから従う（したがって Lemma \ref{formal-defos-lemma-RS-implies-S1-S2}
により (S2) を満たす）。
\end{proof}

\begin{lemma}
\label{formal-defos-lemma-change-of-fields-smooth}
記法と仮定は Situation \ref{formal-defos-situation-change-of-fields} のとおりとする。
$\mathcal{F} \to \mathcal{G}$ が $\mathcal{C}_{\Lambda, k}$ 上で群亜群を
ファイバーとする余ファイバー圏の滑らかな射ならば，
$\mathcal{F}_{l/k} \to \mathcal{G}_{l/k}$ は
$\mathcal{C}_{\Lambda, l}$ 上で群亜群をファイバーとする余ファイバー圏の
滑らかな射である。
\end{lemma}

\begin{proof}
これは定義と，(\ref{formal-defos-equation-comparison}) が全射を保つことから直ちに従う。
\end{proof}

\noindent
$\mathcal{F}$ と $\mathcal{F}_{l/k}$ の関係については，さらに多くのこと
（とくに versal 変形の関係について）を述べられるので，必要に応じて
ここに追加する。

\begin{lemma}
\label{formal-defos-lemma-change-of-field-versal-ring}
記法と仮定は Situation \ref{formal-defos-situation-change-of-fields} のとおりとする。
$\xi$ を，$R \in \Ob(\widehat{\mathcal{C}}_{\Lambda, k})$ 上にある
$\mathcal{F}$ の versal 形式対象とする。このとき次が存在する：
\begin{enumerate}
\item $S \in \Ob(\widehat{\mathcal{C}}_{\Lambda, l})$ と局所
$\Lambda$-代数準同型 $R \to S$ であって，$\mathfrak m_S$-進位相について
形式的に滑らかであり，剰余体上で与えられた体拡大 $l/k$ を誘導するもの。
\item $S$ 上にある $\mathcal{F}_{l/k}$ の versal 形式対象。
\end{enumerate}
\end{lemma}

\begin{proof}
$S$ の構成。$R$-代数の全射 $R[x_1, \ldots, x_n] \to l$ を選ぶ。
その核は極大イデアル $\mathfrak m$ である。$S$ を Noether 環
$R[x_1, \ldots, x_n]$ の $\mathfrak m$-進完備化とする。このとき
Algebra, Lemma \ref{algebra-lemma-completion-Noetherian-Noetherian} により
$S$ は $\widehat{\mathcal{C}}_{\Lambda, l}$ に属する。写像 $R \to S$ は，
More on Algebra, Lemmas \ref{more-algebra-lemma-formally-smooth} および
\ref{more-algebra-lemma-formally-smooth-completion} と，
$R \to R[x_1, \ldots, x_n]$ が形式的に滑らかであることにより，
$\mathfrak m_S$-進位相について形式的に滑らかである。
% P10-E0517: use the grammatical proof-of citation; frozen missing preposition is sidecar-only.
（Lemma \ref{formal-defos-lemma-exists-smooth} の証明と比較せよ。）

\medskip\noindent
$\xi$ は versal なので，変換
$\underline{\xi} : \underline{R}|_{\mathcal{C}_{\Lambda, k}} \to \mathcal{F}$
は滑らかである。Lemma \ref{formal-defos-lemma-change-of-fields-smooth} により，誘導写像
$$
(\underline{R}|_{\mathcal{C}_{\Lambda, k}})_{l/k}
\longrightarrow
\mathcal{F}_{l/k}
$$
は滑らかである。したがって滑らかな射
$\underline{S}|_{\mathcal{C}_{\Lambda, l}} \to
(\underline{R}|_{\mathcal{C}_{\Lambda, k}})_{l/k}$
を構成すれば十分である。このような写像を与えることは，
$\mathcal{C}_{\Lambda, l}$ の各対象 $B$ に対して，$B$ に関手的な集合の写像
$$
\Mor_{\widehat{\mathcal{C}}_{\Lambda, l}}(S, B)
\longrightarrow
\Mor_{\widehat{\mathcal{C}}_{\Lambda, k}}(R, B \times_l k)
$$
を与えることを意味する。左辺の元 $\varphi : S \to B$ が与えられたとき，
これを，像が部分 $\Lambda$-代数 $B \times_l k$ に含まれる合成
$R \to S \to B$ に送る。この写像の滑らかさは，全射 $B' \to B$ と
可換図式
$$
\xymatrix{
S \ar[r] & B \ar@{=}[r] & B \\
R \ar[u] \ar[r] & B' \times_l k \ar[u] \ar[r] & B' \ar[u]
}
$$
が与えられたとき，外側の長方形に収まる環写像 $S \to B'$ を
見つけなければならないことを意味する。この写像の存在は，
$R \to S$ を $\mathfrak m_S$-進位相について形式的に滑らかとなるように
選んだことから保証される；
More on Algebra, Lemma \ref{more-algebra-lemma-lift-continuous} を参照せよ。
\end{proof}
